\documentclass[11pt]{article}
\usepackage[margin=1in]{geometry}
\usepackage[T1]{fontenc}
\usepackage[utf8]{inputenc}
\usepackage[french]{babel}
\usepackage{amsmath,amssymb,amsthm}
\usepackage[colorlinks=true,linkcolor=blue,urlcolor=blue]{hyperref}
\newtheorem{problem}{Problème}
\newenvironment{refs}{\par\small\noindent\emph{Références :}\begin{itemize}\setlength\itemsep{0pt}}{\end{itemize}\normalsize}
\title{Hors de la Caverne \\ \Large Casse-tête}
\author{}
\date{}
\begin{document}
\maketitle
\noindent\emph{Des problèmes, pas de réponses.} Chaque problème ci-dessous est posé
sans solution. Les références indiquent où trouver les connaissances nécessaires.
\medskip


\section*{Collège}

\begin{problem}[{Les trois boîtes mal étiquetées --- Collège}]
\textbf{PUZ-101.} \textbf{Problème.} Trois boîtes fermées portent les étiquettes \og{} pommes \fg{}, \og{} oranges \fg{} et
\og{} mélange \fg{}. L'une ne contient que des pommes, une autre que des oranges, la troisième un
mélange des deux --- mais les étiquettes ont été posées n'importe comment et, par malchance,
aucune des trois n'est correcte. Vous pouvez sortir un fruit de la boîte de votre choix, sans
regarder à l'intérieur.
\begin{enumerate}
\item[(a)] Déterminez le contenu des trois boîtes en ne sortant qu'un seul fruit, et démontrez que
votre procédé aboutit quel que soit le fruit que vous obtenez.
\item[(b)] Démontrez qu'après ce tirage la réponse est réellement forcée : il n'existe pas deux
répartitions différentes des contenus qui soient compatibles à la fois avec \og{} aucune
étiquette n'est correcte \fg{} et avec le fruit que vous avez vu.
\item[(c)] On affaiblit l'hypothèse : on vous garantit seulement qu'\textit{au moins} une étiquette est
fausse. Une boîte mixte peut vous donner l'un ou l'autre fruit, et c'est elle qui décide,
pas vous. Déterminez s'il existe encore une procédure qui garantisse la réponse en un nombre
fini de tirages, et démontrez ce que vous affirmez.
\end{enumerate}

\end{problem}

\noindent Ce casse-tête circule depuis des décennies dans les colonnes de jeux et
dans les entretiens d'embauche, et il doit sa longévité à un détail que la plupart des gens
lisent sans le voir : la phrase \og{} aucune étiquette n'est correcte \fg{} n'est pas une décoration
de l'histoire, c'est la donnée qui rend un tirage unique suffisant. Traduisez-la en langage de
permutations --- quelles répartitions des trois contenus sur les trois étiquettes restent
debout ? --- et le casse-tête devient un exercice de comptage plutôt qu'un exercice de chance.
La question (c) est là pour montrer ce que coûte l'affaiblissement d'une hypothèse : un énoncé
soluble en une observation peut devenir insoluble en un million, et la démonstration de cette
impossibilité se mène en imaginant un adversaire qui choisit les fruits au fur et à mesure,
uniquement pour vous garder dans le doute. C'est exactement la technique qu'emploient les
informaticiens pour minorer le nombre de questions nécessaires à un algorithme.

\begin{refs}
\item Contexte : Martin Gardner --- Wikipédia (\url{https://fr.wikipedia.org/wiki/Martin_Gardner})
\item Contexte : Raisonnement par l'absurde --- Wikipédia (\url{https://fr.wikipedia.org/wiki/Raisonnement_par_l%27absurde})
\item Voir aussi : PUZ-01 (\url{#PUZ-01}) ci-dessous, pour la même façon de compter ce qu'une seule observation peut trancher
\end{refs}

\bigskip

\begin{problem}[{L'escargot au fond du puits --- Collège, short}]
\textbf{PUZ-102.} \textbf{Problème.} Un escargot part du fond d'un puits de $ 10 $ mètres :
chaque jour il grimpe de $ 3 $ mètres, chaque nuit il en redescend $ 2 $. Déterminez le
jour où il atteint le bord, et démontrez que votre compte est exact --- c'est-à-dire que ni la
veille ni le lendemain ne conviennent.
\end{problem}

\noindent Les puits, les fosses et les murs rongés par les rats peuplent les recueils
d'arithmétique du Moyen Âge : le \textit{Liber abaci} de Leonardo Fibonacci, en 1202, contient
la version du lion tombé au fond d'une fosse, qui remonte d'une certaine fraction le jour et
reglisse d'une autre la nuit. Ce qui a fait traverser huit siècles à ce problème, ce n'est pas
la soustraction $ 3 - 2 $ mais le fait que le dernier jour ne ressemble pas aux autres : le
régime régulier cesse de s'appliquer au moment précis où l'on en a besoin. C'est la même
erreur d'une unité que l'on commet en comptant les piquets d'une clôture, en numérotant les
pages d'un livre ou en écrivant une boucle informatique, et savoir la traquer vaut mieux que
savoir diviser.

\begin{refs}
\item Contexte : Liber abaci --- Wikipédia (\url{https://fr.wikipedia.org/wiki/Liber_abaci})
\item Contexte : Leonardo Fibonacci --- Wikipédia (\url{https://fr.wikipedia.org/wiki/Leonardo_Fibonacci})
\end{refs}

\bigskip

\begin{problem}[{Le carré magique de Luo Shu --- Collège}]
\textbf{PUZ-103.} \textbf{Problème.} Un carré magique d'ordre $ 3 $ est un tableau de $ 3 \times 3 $ cases
rempli avec les entiers de $ 1 $ à $ 9 $, chacun employé une fois et une seule, de sorte
que les trois lignes, les trois colonnes et les deux diagonales aient toutes la même
somme.
\begin{enumerate}
\item[(a)] Déterminez, avec démonstration, la valeur de cette somme commune, puis celle que porte
nécessairement la case centrale.
\item[(b)] Complétez le carré ci-dessous et démontrez qu'il n'admet qu'une seule complétion :
\[ \begin{array}{ccc} 8 & 1 & ? \\ ? & ? & ? \\ ? & ? & 2 \end{array} \]
\item[(c)] Démontrez qu'il existe exactement huit carrés magiques d'ordre $ 3 $, et qu'ils se
déduisent tous d'un même carré par rotations et symétries.
\item[(d)] On remplace $ 1, \dots, 9 $ par neuf entiers distincts quelconques, la définition étant
inchangée. Démontrez que la case centrale vaut toujours le tiers de la somme magique, et
construisez une infinité de tels carrés.
\end{enumerate}

\end{problem}

\noindent Le carré de Luo Shu est le plus ancien carré magique connu : la tradition
chinoise le fait apparaître sur la carapace d'une tortue sortie de la rivière Luo, et il
figure dans les textes cosmologiques bien avant d'être vu comme un objet mathématique. On le
retrouve ensuite partout --- dans le monde arabe, dans l'Inde médiévale, puis en Europe, où
Albrecht Dürer place un carré magique d'ordre $ 4 $ dans sa gravure \textit{Melencolia I} de
$ 1514 $, daté par les deux cases centrales de sa dernière ligne. La beauté de l'ordre
$ 3 $ est qu'il ne laisse aucune place au hasard : les huit alignements se recoupent
tellement que presque tout est forcé dès les premières lignes de raisonnement, et la question
(c) demande de transformer cette impression en démonstration. C'est un excellent premier
contact avec l'idée de classification --- non pas exhiber un objet, mais prouver qu'on les a
tous.

\begin{refs}
\item Contexte : Carré de Luo Shu --- Wikipédia (\url{https://fr.wikipedia.org/wiki/Carr%C3%A9_de_Luo_Shu})
\item Contexte : Carré magique (mathématiques) --- Wikipédia (\url{https://fr.wikipedia.org/wiki/Carr%C3%A9_magique_(math%C3%A9matiques)})
\item Voir aussi : PUZ-16 (\url{#PUZ-16}) ci-dessous, la variante avec des carrés parfaits, toujours ouverte
\end{refs}

\bigskip

\begin{problem}[{Les quatre poids de Bachet --- Collège}]
\textbf{PUZ-104.} \textbf{Problème.} Une balance à deux plateaux ne fait que comparer : elle dit quel côté
descend, ou que les deux s'équilibrent. Vous cherchez un jeu de quatre poids, de masses
entières en kilogrammes, qui permette de peser n'importe quelle masse entière de $ 1 $ à
$ 40 $ kg : la marchandise est posée sur un plateau, et vous avez le droit de répartir vos
poids sur les \textit{deux} plateaux à votre guise.
\begin{enumerate}
\item[(a)] Trouvez un tel jeu de quatre poids, et démontrez qu'il pèse effectivement toutes les
masses entières de $ 1 $ à $ 40 $.
\item[(b)] Démontrez qu'aucun jeu de quatre poids ne peut faire mieux : quelles que soient les quatre
masses choisies, il est impossible de peser toutes les masses entières de $ 1 $ à
$ 41 $.
\item[(c)] Reprenez tout avec la règle du marchand : les poids ne peuvent aller que sur le plateau
opposé à la marchandise. Déterminez le meilleur jeu de quatre poids et démontrez la borne
correspondante.
\item[(d)] Généralisez les deux régimes à $ n $ poids : donnez dans chaque cas la plus grande masse
$ M(n) $ telle que toutes les masses entières de $ 1 $ à $ M(n) $ soient pesables, et
démontrez vos deux formules.
\end{enumerate}

\end{problem}

\noindent Claude-Gaspard Bachet de Méziriac pose ce problème dans ses \textit{Problèmes
plaisants et délectables qui se font par les nombres}, en $ 1612 $, le premier recueil
imprimé de mathématiques récréatives digne de ce nom --- c'est aussi l'homme qui traduisit
Diophante en latin, dans la marge duquel Fermat écrira sa fameuse note. Des versions du
problème des poids circulaient déjà chez les marchands médiévaux, mais Bachet en fait une
question de principe, et la réponse est l'une des plus jolies rencontres entre le commerce et
la numération : autoriser un poids à changer de plateau, c'est lui donner trois emplois
possibles au lieu de deux, et compter le nombre de résultats accessibles revient à écrire les
nombres dans une base. La question (b) est la vraie question du site : trouver un jeu qui
marche est affaire de tâtonnement, démontrer qu'aucun jeu ne fait mieux demande de raisonner
sur tous les jeux à la fois.

\begin{refs}
\item Contexte : Claude-Gaspard Bachet de Méziriac --- Wikipédia (\url{https://fr.wikipedia.org/wiki/Claude-Gaspard_Bachet_de_M%C3%A9ziriac})
\item Contexte : Système ternaire --- Wikipédia (\url{https://fr.wikipedia.org/wiki/Syst%C3%A8me_ternaire})
\item Contexte : Problème de pesée --- Wikipédia (\url{https://fr.wikipedia.org/wiki/Probl%C3%A8me_de_pes%C3%A9e})
\end{refs}

\bigskip

\begin{problem}[{Casser la tablette de chocolat --- Collège, short}]
\textbf{PUZ-105.} \textbf{Problème.} Une tablette de chocolat est un rectangle de
$ m \times n $ carrés séparés par des rainures, et un coup consiste à prendre un morceau et
à le casser en deux en suivant une rainure, d'un bord à l'autre. Déterminez le nombre de
cassures nécessaires pour réduire la tablette en $ mn $ carrés isolés, et démontrez que ce
nombre ne dépend jamais de l'ordre dans lequel on casse.
\end{problem}

\noindent Voilà le premier exemple d'invariant qu'on donne aux élèves dans les cercles
mathématiques russes, et il est difficile d'en trouver un meilleur : la stratégie ne sert à
rien, l'ingéniosité ne sert à rien, parce qu'une certaine quantité varie de la même façon à
chaque coup quel que soit le coup. Le jour où l'on comprend qu'il faut suivre cette
quantité-là plutôt que la forme des morceaux, on a appris un réflexe qui resservira devant
des problèmes autrement plus sérieux --- un invariant transforme une infinité de scénarios en
une seule ligne de calcul. Le raisonnement se rédige aussi bien par récurrence que par
comptage direct, et comparer les deux rédactions est instructif.

\begin{refs}
\item Contexte : Invariant --- Wikipédia (\url{https://fr.wikipedia.org/wiki/Invariant})
\item Contexte : Raisonnement par récurrence --- Wikipédia (\url{https://fr.wikipedia.org/wiki/Raisonnement_par_r%C3%A9currence})
\item Lecture : Dmitri Fomin, Sergey Genkin \& Ilia Itenberg, \textit{Mathematical Circles (Russian Experience)}, chapitre sur les invariants
\end{refs}

\bigskip

\begin{problem}[{L'héritage des dix-sept chameaux --- Collège}]
\textbf{PUZ-106.} \textbf{Problème.} Un homme laisse dix-sept chameaux à ses trois fils : la moitié à l'aîné,
le tiers au deuxième, le neuvième au cadet. Comme dix-sept n'est divisible ni par deux, ni par
trois, ni par neuf, les fils se querellent. Un voisin prête alors un chameau : le partage se
fait sur dix-huit bêtes, chacun emporte sa part entière, et le voisin repart avec son
chameau.
\begin{enumerate}
\item[(a)] Effectuez le partage sur dix-huit bêtes, et vérifiez que chaque fils reçoit un nombre
entier de chameaux et qu'il en reste exactement un à rendre.
\item[(b)] Chaque fils a-t-il reçu exactement la part que le testament lui promettait ? Comparez et
expliquez précisément d'où vient l'écart.
\item[(c)] Formalisez ce qui fait marcher le tour : on cherche trois fractions $ \frac{1}{a} $,
$ \frac{1}{b} $, $ \frac{1}{c} $ avec $ 2 \leq a \leq b \leq c $ et un troupeau de
$ N $ bêtes tels que l'emprunt d'une seule bête donne à chacun un nombre entier de bêtes et
laisse exactement une bête à rendre. Écrivez cette condition sous forme d'équation.
\item[(d)] Trouvez tous les testaments possibles, et démontrez que votre liste est complète. Traitez
ensuite le cas de deux héritiers, puis de quatre.
\end{enumerate}

\end{problem}

\noindent L'histoire se raconte de l'Iran au Maghreb comme un conte de sagesse, avec
un cadi ou un vieillard à la place du voisin ; la plus ancienne trace imprimée qu'on lui
connaisse en Europe se trouve chez Niccolò Fontana Tartaglia, au milieu du seizième siècle,
dans son \textit{General trattato di numeri et misure}. Le tour n'est pas un tour de magie mais
une identité entre fractions unitaires, du genre de celles que les scribes de l'Égypte
ancienne manipulaient déjà : dire laquelle, et pourquoi elle permet au chameau emprunté de
revenir, c'est exactement le travail des questions (b) et (c). La question (d) est un vrai
petit problème de théorie des nombres --- chercher toutes les décompositions d'un nombre en
somme de fractions unitaires est une activité qui remonte au papyrus Rhind et qui n'a jamais
cessé d'occuper les mathématiciens.

\begin{refs}
\item Contexte : Fraction égyptienne --- Wikipédia (\url{https://fr.wikipedia.org/wiki/Fraction_%C3%A9gyptienne})
\item Contexte : Niccolò Fontana Tartaglia --- Wikipédia (\url{https://fr.wikipedia.org/wiki/Niccol%C3%B2_Fontana_Tartaglia})
\end{refs}

\bigskip

\begin{problem}[{Le verre d'eau et le verre de vin --- Collège, short}]
\textbf{PUZ-107.} \textbf{Problème.} Un verre contient un décilitre d'eau, un autre un
décilitre de vin ; on prend une cuillerée d'eau qu'on verse dans le vin, on remue autant qu'on
veut, puis on reprend une cuillerée du mélange qu'on reverse dans l'eau. Y a-t-il alors plus
de vin dans l'eau que d'eau dans le vin, et pouvez-vous le démontrer sans rien supposer sur
l'homogénéité du mélange ?
\end{problem}

\noindent C'est un classique des recueils de récréations anglais de la fin du dix-neuvième
siècle, et sa réputation vient de ce que les calculs de proportions, qui semblent la voie
naturelle, sont à la fois pénibles et inutiles : la bonne démonstration est une comptabilité,
pas un calcul de concentrations, et elle tient en deux lignes qui ne supposent rien du tout
sur la façon dont on a remué. On peut même laisser la cuillerée entièrement non mélangée, ou
la mélanger parfaitement, ou n'importe quoi entre les deux, sans changer la réponse --- c'est
ce genre de robustesse qui signale qu'une quantité se conserve quelque part. Le même
raisonnement, transposé, sert en chimie et en comptabilité en partie double.

\begin{refs}
\item Contexte : Problème du mélange eau-vin --- Wikipédia (en anglais) (\url{https://en.wikipedia.org/wiki/Wine/water_mixing_problem})
\item Contexte : Conservation de la masse --- Wikipédia (\url{https://fr.wikipedia.org/wiki/Conservation_de_la_masse})
\end{refs}

\bigskip

\begin{problem}[{L'éléphant de Cao Chong --- Collège}]
\textbf{PUZ-108.} \textbf{Problème.} On veut connaître la masse d'un éléphant. On dispose d'une rivière, d'une
barque assez grande pour le porter, d'un tas de pierres, d'un morceau de craie et d'une
balance ordinaire qui ne pèse que de petites charges --- une pierre à la fois.
\begin{enumerate}
\item[(a)] Décrivez une méthode qui donne la masse de l'éléphant, et justifiez chaque étape : énoncez
en particulier, avec précision, le principe physique qui autorise la comparaison que vous
faites.
\item[(b)] Démontrez ce principe dans le cas dont vous vous servez : deux chargements qui font
flotter la même barque exactement à la même ligne d'eau ont la même masse. Vous ne supposerez
que ceci --- la barque flotte à l'équilibre, et l'eau a partout la même masse volumique.
\item[(c)] Version chiffrée. À la ligne de flottaison, la barque a une section horizontale constante
de $ 6 $ m$^2$, et chargée de l'éléphant elle s'enfonce de $ 40 $ cm de plus qu'à vide.
Calculez la masse de l'éléphant, dites où exactement vous avez utilisé l'hypothèse \og{} section
constante \fg{}, et expliquez ce que devient le calcul si la coque s'évase vers le haut.
\item[(d)] On opère maintenant dans un bassin étroit, dont le niveau monte sensiblement quand la
barque s'enfonce. La méthode de la question (a) reste-t-elle exacte ? Démontrez votre
réponse.
\end{enumerate}

\end{problem}

\noindent Cao Chong, fils de Cao Cao, est mort à douze ans en l'an $ 208 $, au début
de la période des Trois Royaumes ; les \textit{Chroniques des Trois Royaumes} de Chen Shou
racontent que, tout enfant, il résolut le problème de la pesée d'un éléphant offert à son père
en le faisant monter dans une barque, en marquant la ligne d'eau, puis en remplaçant l'animal
par des pierres jusqu'à retrouver la même marque. L'histoire est aujourd'hui encore l'une des
plus connues des écoliers chinois, et rien n'indique que la Chine de cette époque connaissait
Archimède, mort quatre siècles plus tôt à Syracuse : la même idée physique a été trouvée deux
fois. Ce qu'il y a de mathématique ici, c'est le procédé de substitution --- on ne mesure pas
l'objet, on lui substitue quelque chose de mesurable qui lui est égal pour la seule propriété
qui compte ---, et c'est le geste de base de toute la métrologie, des étalons du kilogramme aux
balances de laboratoire.

\begin{refs}
\item Contexte : Cao Chong --- Wikipédia (\url{https://fr.wikipedia.org/wiki/Cao_Chong})
\item Contexte : Poussée d'Archimède --- Wikipédia (\url{https://fr.wikipedia.org/wiki/Pouss%C3%A9e_d%27Archim%C3%A8de})
\item Voir aussi : PUZ-59 (\url{#PUZ-59}) ci-dessous, où la même poussée décide du niveau d'un bassin
\end{refs}

\bigskip

\begin{problem}[{La corde autour de la Terre --- Collège, short}]
\textbf{PUZ-109.} \textbf{Problème.} Une corde tendue fait exactement le tour de l'équateur
d'une planète parfaitement sphérique, en épousant le sol. On l'allonge d'un mètre et on la
dispose en cercle à hauteur constante au-dessus du sol : calculez cette hauteur, puis
démontrez qu'elle est la même pour la Terre, pour un ballon de football et pour le Soleil.
\end{problem}

\noindent On fait remonter ce casse-tête à William Whiston, successeur de Newton à
Cambridge, qui le glisse dans un ouvrage du tout début du dix-huitième siècle ; depuis, il
n'a jamais cessé d'être posé, parce qu'il ne se démode pas --- l'écart entre ce que la formule
donne et ce que l'intuition annonce reste violent même quand on connaît la réponse. Le calcul
est à la portée d'un élève de cinquième, et c'est bien là ce qui rend le problème
philosophique : vous devez arbitrer entre une formule que vous savez par cœur et une
conviction qui hurle le contraire, et l'un des deux se trompe. Une variante bien plus
difficile consiste à ne pas soulever la corde partout, mais à la tirer en un seul point pour
la décoller du sol : la hauteur obtenue est alors surprenante en sens inverse, et sa
détermination sort de l'arithmétique élémentaire.

\begin{refs}
\item Contexte : Circonférence --- Wikipédia (\url{https://fr.wikipedia.org/wiki/Circonf%C3%A9rence})
\item Contexte : William Whiston --- Wikipédia (\url{https://fr.wikipedia.org/wiki/William_Whiston})
\item Contexte : Pi --- Wikipédia (\url{https://fr.wikipedia.org/wiki/Pi})
\end{refs}

\bigskip

\begin{problem}[{Le champ de Didon --- Collège}]
\textbf{PUZ-110.} \textbf{Problème.} Vous disposez d'une corde de longueur $ L $ et vous voulez enclore le
champ d'aire la plus grande possible.
\begin{enumerate}
\item[(a)] Parmi tous les rectangles de périmètre $ L $, déterminez celui dont l'aire est maximale,
et démontrez qu'aucun autre rectangle ne fait mieux.
\item[(b)] Le champ borde maintenant une rivière rectiligne, et l'on n'a pas besoin de clôturer ce
côté-là. Déterminez le meilleur rectangle, démontrez-le, et comparez son aire à celle
obtenue en (a).
\item[(c)] Démontrez que le disque de périmètre $ L $ est plus grand que tous les rectangles de
périmètre $ L $ --- la valeur $ \pi \approx 3{,}14 $ suffit. Traitez de même le cas de la
rivière, en comparant le demi-disque au meilleur rectangle de (b).
\item[(d)] Le théorème isopérimétrique affirme que le disque bat toute autre figure de même
périmètre. Énoncez-le proprement, puis expliquez ce qui manquait à la démonstration de Jakob
Steiner en $ 1838 $ et ce qu'il a fallu lui ajouter.
\end{enumerate}

\end{problem}

\noindent Virgile raconte dans l'\textit{Énéide} que Didon, fuyant Tyr, obtint sur la
côte d'Afrique autant de terre qu'une peau de bœuf pouvait en couvrir ; elle fit découper la
peau en lanières très fines, les mit bout à bout, et enferma dans cette corde la colline qui
devint Carthage. Les mathématiciens ont gardé le nom : le \og{} problème de Didon \fg{} est le premier
problème d'optimisation de forme de l'histoire, et Zénodore, vers le deuxième siècle avant
notre ère, démontrait déjà que le cercle l'emporte sur les polygones. La question (d) est une
leçon d'honnêteté mathématique : Steiner produit en $ 1838 $ de superbes arguments de
symétrisation montrant que toute figure autre que le disque peut être améliorée, mais cela ne
démontre pas que l'optimum existe --- Weierstrass devra fournir plus tard l'argument
d'existence. Un raisonnement qui prouve \og{} si un maximum existe, c'est celui-là \fg{} n'est pas une
démonstration tant que la parenthèse n'est pas refermée.

\begin{refs}
\item Contexte : Isopérimétrie --- Wikipédia (\url{https://fr.wikipedia.org/wiki/Isop%C3%A9rim%C3%A9trie})
\item Contexte : Didon --- Wikipédia (\url{https://fr.wikipedia.org/wiki/Didon})
\item Contexte : Jakob Steiner --- Wikipédia (\url{https://fr.wikipedia.org/wiki/Jakob_Steiner})
\end{refs}

\bigskip

\begin{problem}[{Le carré qui gagne un carreau --- Collège}]
\textbf{PUZ-111.} \textbf{Problème.} Sur du papier quadrillé, prenez le carré $ 8 \times 8 $ dont les coins
opposés sont $ (0,0) $ et $ (8,8) $, et découpez-le ainsi : une coupe verticale en
$ x = 5 $ sur toute la hauteur ; dans la bande de gauche, la coupe droite qui joint
$ (0,3) $ à $ (5,5) $ ; dans la bande de droite, la coupe droite qui joint $ (5,0) $ à
$ (8,8) $. Vous obtenez deux trapèzes et deux triangles. Rangez maintenant ces quatre pièces
dans un rectangle $ 13 \times 5 $ de coins $ (0,0) $ et $ (13,5) $ : un triangle en
$ (0,0), (8,0), (8,3) $, un trapèze en $ (0,0), (0,5), (5,5), (5,2) $, le second triangle
en $ (5,2), (5,5), (13,5) $ et le second trapèze en $ (8,0), (13,0), (13,5), (8,3) $.
\begin{enumerate}
\item[(a)] Vérifiez que les quatre pièces découpent exactement le carré --- ni trou, ni recouvrement ---
et calculez leur aire totale. Vérifiez ensuite qu'aucune des quatre ne dépasse du rectangle et
qu'elles ne se recouvrent pas.
\item[(b)] Le rectangle a pour aire $ 65 $. D'où vient le carreau supplémentaire ? Démontrez que
les points $ (0,0) $, $ (8,3) $ et $ (13,5) $ ne sont pas alignés, puis calculez
exactement l'aire de la partie du rectangle que les quatre pièces laissent découverte.
\item[(c)] Recommencez avec le carré $ 13 \times 13 $ et le rectangle $ 8 \times 21 $ : coupe
verticale en $ x = 8 $, coupe de $ (0,5) $ à $ (8,8) $ à gauche, coupe de $ (8,0) $ à
$ (13,13) $ à droite. Cette fois l'écart change de signe : dites lequel, expliquez ce que
cela signifie concrètement pour les pièces, et démontrez que dans les deux cas l'écart vaut
exactement une unité d'aire.
\item[(d)] Démontrez l'identité de Cassini $ F_{n-1}F_{n+1} - F_n^2 = (-1)^n $ pour la suite de
Fibonacci $ 1, 1, 2, 3, 5, 8, 13, 21, \dots $, et déduisez-en que le tour se refait
indéfiniment. Expliquez enfin, en comparant les pentes $ F_{n-2}/F_{n-1} $ et
$ F_{n-1}/F_n $, pourquoi l'illusion devient de plus en plus convaincante quand $ n $
grandit.
\end{enumerate}

\end{problem}

\noindent Le paradoxe est publié par Oskar Schlömilch en $ 1868 $ ; Sam Loyd s'en
attribuera plus tard la paternité, comme de beaucoup d'autres, et Lewis Carroll, qui était
professeur de mathématiques à Oxford, s'y est intéressé de près. Sa version en triangle, due à
Paul Curry vers $ 1953 $, a fait le tour du monde sous le nom de triangle de Curry. Ce qu'il
y a d'instructif ici, c'est que la démonstration ne consiste pas à trouver l'erreur du
découpage --- le découpage du carré est parfaitement exact --- mais à découvrir que l'assemblage
en rectangle est \textit{presque} exact, et que le \og{} presque \fg{} a une aire mesurable : une
lamelle si allongée que l'œil la prend pour un trait de crayon. Derrière l'illusion se cache
une identité algébrique sur la suite de Fibonacci, et c'est pourquoi le même truc marche avec
$ 3, 5, 8 $, avec $ 5, 8, 13 $, avec $ 8, 13, 21 $, toujours à une unité près.

\begin{refs}
\item Contexte : Paradoxe du carré manquant --- Wikipédia (\url{https://fr.wikipedia.org/wiki/Paradoxe_du_carr%C3%A9_manquant})
\item Contexte : Identités de Cassini, de Catalan et de Vajda --- Wikipédia (\url{https://fr.wikipedia.org/wiki/Identit%C3%A9s_de_Cassini,_de_Catalan_et_de_Vajda})
\item Contexte : Suite de Fibonacci --- Wikipédia (\url{https://fr.wikipedia.org/wiki/Suite_de_Fibonacci})
\end{refs}

\bigskip

\begin{problem}[{La mouche entre les deux trains --- Collège}]
\textbf{PUZ-112.} \textbf{Problème.} Deux trains roulent l'un vers l'autre sur la même voie, chacun à
$ 50 $ km/h, séparés au départ par $ 100 $ km. Une mouche quitte au même instant le
pare-brise du premier à $ 75 $ km/h, atteint le second, fait aussitôt demi-tour, revient au
premier, repart, et ainsi de suite jusqu'à être écrasée entre les deux.
\begin{enumerate}
\item[(a)] Calculez la distance totale parcourue par la mouche, en justifiant chaque étape --- à
commencer par l'instant de la collision.
\item[(b)] Calculez la longueur du premier trajet de la mouche, puis celle du deuxième. Démontrez que
ces longueurs forment une suite géométrique dont vous déterminerez la raison, et expliquez
comment une somme d'une infinité de longueurs toutes strictement positives peut être
finie.
\item[(c)] Démontrez que la mouche fait une infinité de demi-tours avant la collision, alors que la
durée totale du vol est finie.
\item[(d)] Renversez le temps. Montrez que, connaissant le lieu et l'instant de la collision ainsi que
la vitesse de la mouche, on ne peut pas reconstituer son point de départ : plusieurs histoires
différentes aboutissent exactement au même écrasement.
\end{enumerate}

\end{problem}

\noindent L'anecdote est la plus célèbre des mathématiques du vingtième siècle : on
pose le problème à John von Neumann, qui donne la réponse presque immédiatement ; son
interlocuteur, un peu déçu, remarque que la plupart des gens essaient de sommer la série au
lieu de voir l'astuce, et von Neumann répond qu'il a justement sommé la série. Les deux voies
sont légitimes et le problème vaut par leur confrontation : l'une regarde le temps, l'autre
regarde la distance, et la question (b) vous demande de vérifier qu'elles donnent bien le même
nombre. La question (c) est un authentique paradoxe de Zénon apprivoisé --- une infinité
d'événements logés dans un intervalle fini ---, et la question (d) montre qu'une équation du
mouvement peut se laisser remonter dans un sens et pas dans l'autre, ce qui est une remarque
plus profonde qu'elle n'en a l'air.

\begin{refs}
\item Contexte : John von Neumann --- Wikipédia (\url{https://fr.wikipedia.org/wiki/John_von_Neumann})
\item Contexte : Série géométrique --- Wikipédia (\url{https://fr.wikipedia.org/wiki/S%C3%A9rie_g%C3%A9om%C3%A9trique})
\item Contexte : Paradoxes de Zénon --- Wikipédia (\url{https://fr.wikipedia.org/wiki/Paradoxes_de_Z%C3%A9non})
\end{refs}

\bigskip

\begin{problem}[{Le tournoi et le deuxième meilleur --- Collège}]
\textbf{PUZ-113.} \textbf{Problème.} Huit joueurs disputent un tournoi à élimination directe : quatre matchs au
premier tour, deux au deuxième, une finale. On suppose qu'il existe un classement strict des
huit joueurs, que le meilleur des deux gagne toujours, et qu'il n'y a pas de match nul.
\begin{enumerate}
\item[(a)] Démontrez que le vainqueur du tournoi est bien le meilleur des huit. Montrez en revanche
que le finaliste malheureux n'est pas nécessairement le deuxième, en exhibant un classement et
un tableau où le vrai deuxième est éliminé dès le premier tour.
\item[(b)] Le tournoi a coûté sept matchs. Organisez des matchs supplémentaires qui désignent à coup
sûr le deuxième meilleur joueur, pour un total de neuf matchs, et démontrez que votre méthode
ne se trompe jamais.
\item[(c)] Démontrez que huit matchs ne suffisent pas : aucune organisation, même adaptative --- chaque
match étant choisi en fonction des résultats précédents ---, ne peut garantir de désigner à la
fois le premier et le deuxième en huit matchs. C'est la moitié difficile du problème.
\item[(d)] Généralisez à $ n $ joueurs : donnez une méthode en $ n - 2 + \lceil \log_2 n \rceil $
matchs, et démontrez qu'on ne peut pas faire mieux.
\end{enumerate}

\end{problem}

\noindent En $ 1883 $, Charles Dodgson --- que le monde connaît sous le nom de Lewis
Carroll --- publie une brochure sur les tournois de tennis pour dénoncer une injustice bien
réelle : la méthode d'attribution du deuxième prix est fausse, et le joueur qui la reçoit n'est
en général pas le deuxième meilleur. Le problème a été repris sous forme mathématique par Hugo
Steinhaus, et le nombre exact de matchs nécessaires a été établi par S. S. Kislitsyn en
$ 1964 $ ; Donald Knuth en donne le traitement de référence dans le troisième volume de
\textit{The Art of Computer Programming}, où le tournoi devient un algorithme de sélection et le
match une comparaison. La question (c) est le cœur du sujet : pour minorer le nombre de matchs,
on ne raisonne pas sur les tournois possibles mais on imagine un adversaire qui, à chaque
match, choisit le vainqueur de façon à retarder le plus longtemps possible le moment où vous
saurez --- technique dite de l'adversaire, aujourd'hui centrale en complexité algorithmique.

\begin{refs}
\item Contexte : Tournoi à élimination directe --- Wikipédia (\url{https://fr.wikipedia.org/wiki/Tournoi_%C3%A0_%C3%A9limination_directe})
\item Contexte : Lewis Carroll --- Wikipédia (\url{https://fr.wikipedia.org/wiki/Lewis_Carroll})
\item Lecture : Donald E. Knuth, \textit{The Art of Computer Programming}, vol. 3, \S{} 5.3.3 (sélection par comparaisons)
\end{refs}

\bigskip

\begin{problem}[{Le treize tombe-t-il souvent un vendredi ? --- Collège}]
\textbf{PUZ-114.} \textbf{Problème.} On travaille dans le calendrier grégorien, celui qui est en usage
aujourd'hui : une année est bissextile si elle est multiple de $ 4 $, sauf si elle est
multiple de $ 100 $, sauf si elle est multiple de $ 400 $.
\begin{enumerate}
\item[(a)] Démontrez qu'un mois contient un vendredi $ 13 $ si et seulement si son premier jour est
un dimanche.
\item[(b)] Démontrez que toute année, commune ou bissextile, contient au moins un vendredi $ 13 $.
Déterminez ensuite le nombre maximal de vendredis $ 13 $ que peut compter une même année, et
démontrez votre réponse.
\item[(c)] Démontrez que $ 400 $ années grégoriennes comptent un nombre entier de semaines.
Déduisez-en que le calendrier se répète à l'identique tous les $ 400 $ ans, et donc que la
question \og{} quel jour de la semaine le $ 13 $ tombe-t-il le plus souvent ? \fg{} a bien un
sens.
\item[(d)] Répondez-y : déterminez, avec démonstration, le jour de la semaine le plus fréquent pour un
$ 13 $, et le moins fréquent. Le comptage est long mais fini --- organisez-le pour qu'il tienne
en une page.
\end{enumerate}

\end{problem}

\noindent La peur du vendredi $ 13 $ est récente : elle ne se répand en Europe et en
Amérique qu'au dix-neuvième siècle, en soudant deux superstitions plus anciennes et
indépendantes, celle du vendredi et celle du nombre treize. Le fait mathématique, lui, est
attesté depuis les années $ 1930 $, quand quelqu'un a pris la peine de compter : dans le
cycle grégorien de $ 400 $ ans, le $ 13 $ tombe un vendredi un peu plus souvent que
n'importe quel autre jour de la semaine. Rien dans la réforme de $ 1582 $ ne visait cet
effet ; il résulte du hasard heureux --- la question (c) --- qui fait que la longueur du cycle
grégorien est un multiple de sept, ce qui n'était nullement obligé et qui rend toute la
question calculable à la main. C'est un bel exemple de problème où l'arithmétique modulaire
sert à remplacer un calcul infini par un tableau fini.

\begin{refs}
\item Contexte : Vendredi treize --- Wikipédia (\url{https://fr.wikipedia.org/wiki/Vendredi_treize})
\item Contexte : Calendrier grégorien --- Wikipédia (\url{https://fr.wikipedia.org/wiki/Calendrier_gr%C3%A9gorien})
\item Contexte : Congruence sur les entiers --- Wikipédia (\url{https://fr.wikipedia.org/wiki/Congruence_sur_les_entiers})
\end{refs}

\bigskip

\begin{problem}[{Les allumettes à déplacer --- Collège}]
\textbf{PUZ-115.} \textbf{Problème.} On écrit des nombres en chiffres romains avec des allumettes : le
$ \mathrm{I} $ coûte une allumette, le $ \mathrm{V} $ et le $ \mathrm{X} $ en coûtent
deux chacun, le signe $ + $ en coûte deux, le signe $ - $ une seule, et le signe $ = $
deux. Sur la table on lit l'égalité fausse
\[ \mathrm{XI} + \mathrm{I} = \mathrm{X}. \]
Un coup consiste à retirer une allumette de la figure et à la reposer ailleurs : le nombre
total d'allumettes ne change donc pas, il est interdit d'en ajouter, d'en retirer ou d'en
laisser une qui ne serve à rien, et ce qui reste sur la table doit être une égalité bien
formée --- un nombre, un signe d'opération, un nombre, le signe $ = $, un nombre.
\begin{enumerate}
\item[(a)] Rendez l'égalité vraie en déplaçant une seule allumette.
\item[(b)] Dressez la liste de toutes les égalités vraies que l'on peut obtenir à partir de cette
figure en un seul déplacement, et démontrez que votre liste est complète. C'est ici que
l'invariant \og{} le nombre d'allumettes ne change pas \fg{} fait le gros du travail, en réduisant une
recherche vague à une vérification finie.
\item[(c)] On n'utilise plus que les chiffres $ \mathrm{I} $, $ \mathrm{V} $, $ \mathrm{X} $,
écrits selon les règles habituelles de la numération romaine, ce qui donne les nombres de
$ 1 $ à $ 39 $. Déterminez le plus grand nombre que l'on puisse écrire avec six
allumettes, et démontrez qu'aucun autre ne fait mieux.
\item[(d)] Même question avec sept, puis huit, puis neuf allumettes. Déterminez à partir de combien
d'allumettes on peut écrire le plus grand nombre disponible, et démontrez-le.
\end{enumerate}

\end{problem}

\noindent Les casse-tête d'allumettes envahissent la presse populaire européenne à la
fin du dix-neuvième siècle, à l'époque où l'allumette devient l'objet le plus banal d'une
poche, et Henry Dudeney comme Sam Loyd en publient par centaines. On les prend pour de simples
jeux d'observation, à tort : la méthode qui les résout proprement est celle-là même qu'on
emploie en combinatoire sérieuse --- d'abord un invariant qui élimine d'un coup l'immense
majorité des figures imaginables, ensuite un examen exhaustif du petit nombre de cas qui
survivent. La question (c) ajoute une optimisation authentique, où l'on doit non seulement
exhiber une écriture économique mais démontrer qu'aucune autre ne fait mieux, en majorant ce
qu'une allumette peut rapporter.

\begin{refs}
\item Contexte : Numération romaine --- Wikipédia (\url{https://fr.wikipedia.org/wiki/Num%C3%A9ration_romaine})
\item Contexte : Casse-tête d'allumettes --- Wikipédia (en anglais) (\url{https://en.wikipedia.org/wiki/Matchstick_puzzle})
\item Contexte : Sam Loyd --- Wikipédia (\url{https://fr.wikipedia.org/wiki/Sam_Loyd})
\end{refs}

\bigskip

\begin{problem}[{La côte et la vitesse moyenne --- Collège, short}]
\textbf{PUZ-116.} \textbf{Problème.} Vous montez une côte de deux kilomètres à la vitesse
moyenne de $ 15 $ km/h. À quelle vitesse devez-vous la redescendre pour que votre vitesse
moyenne sur l'aller-retour atteigne $ 30 $ km/h --- et si aucune vitesse ne convient,
démontrez-le.
\end{problem}

\noindent Le psychologue Max Wertheimer, fondateur de la théorie de la forme, aimait
soumettre des énigmes à ses amis pour observer comment ils pensaient ; il posa celle-ci à
Albert Einstein, et l'anecdote, qu'il rapporte dans \textit{Productive Thinking}, veut
qu'Einstein soit d'abord tombé dans le piège. Le piège est le mot \og{} moyenne \fg{} : une vitesse
moyenne n'est pas la moyenne des vitesses, c'est la distance totale divisée par le temps
total, et ces deux quantités ne se comportent pas du tout de la même façon quand on les
combine. Une fois l'énoncé traduit en temps plutôt qu'en vitesses, tout devient limpide --- y
compris le fait qu'une réponse puisse ne pas exister, ce qui est une conclusion parfaitement
respectable et demande, elle aussi, une démonstration.

\begin{refs}
\item Contexte : Moyenne harmonique --- Wikipédia (\url{https://fr.wikipedia.org/wiki/Moyenne_harmonique})
\item Contexte : Max Wertheimer --- Wikipédia (\url{https://fr.wikipedia.org/wiki/Max_Wertheimer})
\end{refs}

\bigskip


\section*{Lycée}

\begin{problem}[{Neuf boules et une balance --- Lycée}]
\textbf{PUZ-01.} \textbf{Problème.} Vous disposez de neuf boules d'apparence identique. Huit pèsent
exactement le même poids ; une est légèrement plus lourde. Votre seul instrument est une
balance à plateaux, qui compare le contenu de ses deux plateaux et vous indique quel côté
est le plus lourd, ou qu'ils s'équilibrent. Elle ne donne aucune valeur numérique.
\begin{enumerate}
\item[(a)] Trouvez une procédure qui identifie la boule lourde en utilisant la balance exactement
deux fois, et démontrez qu'elle fonctionne toujours --- quels que soient les résultats.
\item[(b)] Démontrez qu'une seule pesée ne suffit pas. Démontrez ensuite la borne générale : avec
$ k $ pesées vous ne pouvez pas distinguer plus de $ 3^{k} $ possibilités, parce que
chaque pesée n'a que trois issues. Déduisez-en le plus grand nombre de boules qu'une
procédure à $ k $ pesées peut traiter dans cette version \og{} une boule lourde \fg{}.
\item[(c)] Changez l'énoncé : on sait que la boule intruse a un poids différent, mais on ne vous
dit \textit{pas} si elle est plus lourde ou plus légère. Montrez que deux pesées ne
suffisent alors que pour trois boules, et expliquez exactement où l'argument de comptage
de la question (b) doit être modifié.
\end{enumerate}

\end{problem}

\noindent C'est l'ancêtre de tous les casse-tête de pesée, et sa maison naturelle
est la théorie de l'information. Chaque usage de la balance produit l'un de trois
résultats, donc $ k $ pesées peuvent porter au plus $ \log_2 3^k $ bits d'information ;
une énigme à $ N $ réponses équiprobables demande $ 3^k \geq N $. Ce qui rend la
question (a) satisfaisante, c'est que la borne n'est pas seulement respectée mais
exactement atteinte --- neuf vaut précisément $ 3^2 $, de sorte que la procédure doit tirer
toute l'information de chaque pesée et ne peut se permettre aucune comparaison gaspillée.
La question (c) est là où l'intuition de la plupart des gens casse : ignorer le sens coûte
plus cher qu'il n'y paraît, parce que les issues doivent désormais identifier à la fois
\textit{quelle} boule et \textit{dans quel sens}, et la comptabilité change. La version en
taille industrielle --- douze pièces, sens inconnu, trois pesées --- a tellement circulé
pendant la Seconde Guerre mondiale qu'on plaisantait en y voyant une arme secrète destinée
à faire perdre son temps aux mathématiciens alliés.

\begin{refs}
\item Contexte : Problème de pesée --- Wikipédia (\url{https://fr.wikipedia.org/wiki/Probl%C3%A8me_de_pes%C3%A9e})
\item Contexte : Théorie de l'information --- Wikipédia (\url{https://fr.wikipedia.org/wiki/Th%C3%A9orie_de_l%27information})
\item Lecture : Martin Gardner, \textit{The Scientific American Book of Mathematical Puzzles and Diversions} (1959)
\item Voir aussi : Cryptographie et théorie de l'information (\url{applied-crypto.html}), pour l'entropie comme borne de comptage
\end{refs}

\bigskip

\begin{problem}[{Douze pièces, sens inconnu --- Lycée}]
\textbf{PUZ-02.} \textbf{Problème.} Douze pièces sont identiques, à ceci près qu'exactement l'une d'elles
est fausse et diffère des autres par le poids. Vous ignorez si elle est plus lourde ou
plus légère. Vous disposez d'une balance à plateaux et de trois pesées.
\begin{enumerate}
\item[(a)] Trouvez la pièce fausse \textit{et} déterminez si elle est lourde ou légère, en trois
pesées. Démontrez que votre procédure traite tous les cas.
\item[(b)] Il y a $ 24 $ réponses possibles (douze pièces, deux sens) et trois pesées donnent
$ 27 $ issues. Expliquez pourquoi la marge est si mince, et démontrez que treize pièces
est le vrai maximum pour trois pesées si vous devez aussi annoncer le sens --- et qu'avec
treize vous pouvez trouver la pièce sans toujours pouvoir nommer le sens.
\item[(c)] Concevez une procédure où les trois pesées sont décidées \textit{à l'avance}, avant de
voir la moindre issue. Une telle solution non adaptative existe ; en trouver une est plus
difficile que la version adaptative, et c'est la forme qui se généralise aux codes
correcteurs d'erreurs.
\end{enumerate}

\end{problem}

\noindent C'est le casse-tête le plus difficile que la plupart des gens résolvent
sans aide, et le moment d'illumination --- comprendre qu'une pièce peut passer d'un plateau
à l'autre entre deux pesées, ou être mise de côté entièrement --- reste véritablement en
mémoire. La question (c) est la profonde : un schéma de pesée non adaptatif attribue à
chaque pièce une suite fixe de positions, c'est-à-dire précisément un mot de code, et
l'exigence que toutes les issues soient distinguables est une condition de distance
minimale. Casse-tête de pesée et codes correcteurs d'erreurs sont le même sujet vu de deux
côtés, lien rendu explicite par les travaux de Richard Hamming à la fin des années 1940.
Freeman Dyson en a publié une analyse élégante en 1946, ayant reçu l'énigme alors qu'il
travaillait à la recherche opérationnelle du temps de guerre.

\begin{refs}
\item Contexte : Problème de pesée --- Wikipédia (\url{https://fr.wikipedia.org/wiki/Probl%C3%A8me_de_pes%C3%A9e})
\item Contexte : Code de Hamming --- Wikipédia (\url{https://fr.wikipedia.org/wiki/Code_de_Hamming})
\item Article : Freeman J. Dyson, \og{} The problem of the pennies \fg{}, \textit{The Mathematical Gazette} 30 (1946)
\item Voir aussi : Cryptographie et théorie de l'information (\url{applied-crypto.html})
\end{refs}

\bigskip

\begin{problem}[{Deux mèches, quarante-cinq minutes --- Lycée, short}]
\textbf{PUZ-03.} \textbf{Problème.} Vous avez deux mèches et un briquet. Chaque mèche
brûle exactement une heure d'un bout à l'autre, mais aucune ne brûle à vitesse uniforme,
de sorte qu'une demi-mèche ne fait pas une demi-heure. Mesurez quarante-cinq minutes, et
démontrez que votre méthode est exacte quelle que soit l'irrégularité de la
combustion.
\end{problem}

\noindent L'astuce, une fois vue, est inoubliable, et la raison pour laquelle elle
marche mérite d'être énoncée précisément : brûler une mèche par les deux bouts la consume
en la moitié de sa \textit{durée totale de combustion}, quelle que soit la répartition de
cette durée le long de sa longueur. C'est un énoncé sur une intégrale que l'on divise par
deux, non sur le milieu de la mèche --- et garder ces deux choses distinctes, c'est tout le
casse-tête.

\begin{refs}
\item Contexte : Mèche (pyrotechnie) --- Wikipédia (\url{https://fr.wikipedia.org/wiki/M%C3%A8che_%28pyrotechnie%29})
\item Lecture : Peter Winkler, \textit{Mathematical Puzzles: A Connoisseur's Collection}
\end{refs}

\bigskip

\begin{problem}[{Les cent casiers --- Lycée, short}]
\textbf{PUZ-04.} \textbf{Problème.} Un couloir compte $ 100 $ casiers fermés. L'élève
$ 1 $ bascule l'état de tous les casiers, l'élève $ 2 $ celui d'un casier sur deux,
l'élève $ 3 $ celui d'un casier sur trois, et ainsi de suite jusqu'à l'élève
$ 100 $. Déterminez, avec démonstration, exactement quels casiers sont ouverts à la
fin.
\end{problem}

\noindent Un casier finit ouvert précisément lorsque son numéro a un nombre impair
de diviseurs, et la réponse transforme un casse-tête en petit théorème : les diviseurs
s'apparient selon $ d \leftrightarrow n/d $, donc le compte est impair exactement quand
un diviseur est apparié avec lui-même. Le casse-tête vous demande en réalité de découvrir
pourquoi les carrés parfaits sont particuliers, et c'est l'un des exemples les plus purs
d'argument de comptage qu'un enfant de douze ans peut mener à bout et qu'un théoricien des
nombres respecte encore.

\begin{refs}
\item Contexte : Fonction somme des diviseurs --- Wikipédia (\url{https://fr.wikipedia.org/wiki/Fonction_somme_des_puissances_k-i%C3%A8mes_des_diviseurs})
\item Voir aussi : Théorie des nombres (\url{pure-number-theory.html})
\end{refs}

\bigskip

\begin{problem}[{Le pont et la lampe --- Lycée}]
\textbf{PUZ-05.} \textbf{Problème.} Quatre personnes doivent traverser de nuit un pont branlant. Le pont
supporte au plus deux personnes à la fois, et quiconque traverse doit porter l'unique
lampe. Seules, les quatre mettent $ 1 $, $ 2 $, $ 5 $ et $ 10 $ minutes ; une
paire traverse à l'allure du plus lent. Quelqu'un doit toujours rapporter la lampe.
\begin{enumerate}
\item[(a)] Faites passer les quatre en $ 17 $ minutes.
\item[(b)] Démontrez que $ 17 $ est optimal --- aucun ordonnancement ne fait mieux.
\item[(c)] La stratégie gloutonne évidente, où la personne la plus rapide fait à chaque fois la
navette avec la lampe, prend $ 19 $ minutes. Expliquez précisément pourquoi la gloutonnerie
échoue ici, et dégagez le principe général : les deux plus lents doivent traverser ensemble,
de sorte que le plus grand de leurs temps ne soit payé qu'une fois.
\end{enumerate}

\end{problem}

\noindent La question (b) est la partie mathématique. Produire un ordonnancement en
$ 17 $ minutes est affaire d'essais ; démontrer qu'aucun ordonnancement ne fait mieux
demande de raisonner sur tous les ordonnancements à la fois, et la voie habituelle consiste
à observer que la structure de toute solution est forcée --- les traversées alternent avec
les retours, de sorte qu'avec quatre personnes il y a exactement trois traversées aller et
deux retours. Ce casse-tête est un favori en entretien précisément parce que l'écart entre
(a) et (b) sépare ceux qui ont trouvé une réponse de ceux qui la comprennent. C'est aussi
une véritable instance de problème d'ordonnancement où l'heuristique gloutonne naturelle
est démontrablement sous-optimale, raison pour laquelle il figure dans les cours
d'algorithmique.

\begin{refs}
\item Contexte : Problème du pont et de la lampe --- Wikipédia (en anglais) (\url{https://en.wikipedia.org/wiki/Bridge_and_torch_problem})
\item Contexte : Algorithme glouton --- Wikipédia (\url{https://fr.wikipedia.org/wiki/Algorithme_glouton})
\item Voir aussi : Optimisation et théorie des jeux (\url{applied-optimization.html})
\end{refs}

\bigskip

\begin{problem}[{Le loup, la chèvre et le chou --- Lycée}]
\textbf{PUZ-06.} \textbf{Problème.} Un fermier doit faire traverser une rivière à un loup, une chèvre et
un chou dans une barque qui ne porte que le fermier et au plus un des trois. Laissés seuls
ensemble, le loup mange la chèvre et la chèvre mange le chou.
\begin{enumerate}
\item[(a)] Trouvez une suite de traversées qui fonctionne.
\item[(b)] Modélisez le casse-tête par un graphe : chaque sommet est une configuration licite ---
qui se trouve sur quelle rive, et où est la barque --- et l'on relie deux sommets lorsqu'une
seule traversée transforme l'un en l'autre. Dessinez ce graphe et identifiez chaque
solution comme un chemin dedans.
\item[(c)] Démontrez que votre solution est la plus courte, et déterminez combien de solutions
distinctes de longueur minimale existent.
\end{enumerate}

\end{problem}

\noindent Ce casse-tête a plus de mille ans --- il apparaît dans les
\textit{Propositiones ad Acuendos Juvenes}, un recueil attribué à Alcuin d'York vers l'an
800, ce qui en fait l'un des plus anciens problèmes récréatifs consignés en Europe. La
question (b) est la raison de sa présence sur ce site : dès que vous dessinez le graphe des
configurations, une histoire de bétail devient un objet fini sur lequel vous pouvez
raisonner exhaustivement, et des questions comme \og{} existe-t-il une solution plus
courte ? \fg{} deviennent des questions de longueur de chemin. Cette traduction --- du récit à
l'espace d'états --- est exactement ce que fait un informaticien quand il planifie une
recherche, et elle vaut d'être effectuée à la main au moins une fois.

\begin{refs}
\item Contexte : Problème du loup, de la chèvre et du chou --- Wikipédia (\url{https://fr.wikipedia.org/wiki/Probl%C3%A8me_du_loup%2C_de_la_ch%C3%A8vre_et_du_chou})
\item Contexte : Espace d'états --- Wikipédia (\url{https://fr.wikipedia.org/wiki/Espace_d%27%C3%A9tats})
\item Voir aussi : Combinatoire et théorie des graphes (\url{pure-combinatorics.html}), Mathématiques discrètes et informatique (\url{applied-discrete-cs.html})
\end{refs}

\bigskip

\begin{problem}[{Chevaliers, valets et une seule question --- Lycée}]
\textbf{PUZ-07.} \textbf{Problème.} Sur une île, chaque habitant est soit un chevalier, qui dit toujours
la vérité, soit un valet, qui ment toujours.
\begin{enumerate}
\item[(a)] Vous rencontrez deux insulaires, A et B. A déclare : \og{} Nous sommes tous les deux des
valets. \fg{} Déterminez ce qu'est chacun d'eux, et démontrez que votre déduction est la seule
possibilité.
\item[(b)] Vous arrivez à une bifurcation dont une branche mène au salut. Un seul insulaire s'y
tient, et vous pouvez poser une seule question à réponse oui-ou-non. Trouvez une question
qui révèle la route sûre, et démontrez qu'elle fonctionne que l'insulaire soit chevalier
ou valet.
\item[(c)] Généralisez : montrez qu'une question de la forme \og{} si je vous demandais $ Q $,
répondriez-vous oui ? \fg{} fait toujours émerger la vérité au sujet de $ Q $, et expliquez
pourquoi composer un menteur avec lui-même produit de l'honnêteté.
\end{enumerate}

\end{problem}

\noindent Raymond Smullyan a bâti une carrière sur ces insulaires, et la question
(c) met à nu le mécanisme qui les sous-tend tous : la réponse d'un valet est la négation de
la vérité, de sorte qu'enfouir la question d'un niveau supplémentaire applique la négation
deux fois, et le mensonge s'annule lui-même. C'est une observation de théorie des groupes
déguisée --- les comportements de véracité et de mensonge se composent comme les éléments
d'un groupe d'ordre deux. La même astuce autoréférentielle anime l'énigme logique la plus
difficile jamais posée, le problème des trois dieux de George Boolos, où les dieux
répondent dans une langue inconnue.

\begin{refs}
\item Contexte : Chevaliers et valets --- Wikipédia (en anglais) (\url{https://en.wikipedia.org/wiki/Knights_and_knaves})
\item Contexte : L'Énigme la plus difficile du monde --- Wikipédia (\url{https://fr.wikipedia.org/wiki/L%27%C3%89nigme_la_plus_difficile_du_monde})
\item Lecture : Raymond Smullyan, \textit{What Is the Name of This Book?} (1978) --- publié en français sous le titre \textit{Quel est le titre de ce livre ?}
\item Voir aussi : Logique, théorie des ensembles et fondements (\url{pure-foundations.html})
\end{refs}

\bigskip

\begin{problem}[{Paver un damier avec des pièces 1$\times$4 --- Lycée, short}]
\textbf{PUZ-08.} \textbf{Problème.} On veut paver un plateau $ 10 \times 10 $ par des
pièces $ 1 \times 4 $ placées horizontalement ou verticalement. Démontrez que c'est
impossible.
\end{problem}

\noindent Contrairement à l'échiquier mutilé, un argument à deux couleurs ne suffit
pas ici ; il vous faut un coloriage plus astucieux, et le trouver est tout le casse-tête.
Les arguments de coloriage sont la forme la moins coûteuse de démonstration d'impossibilité
en mathématiques --- vous inventez une quantité que tout coup licite préserve, puis vous
montrez que la cible la viole --- et apprendre à chercher le bon invariant plutôt que la
bonne construction est un véritable changement dans la manière d'attaquer les
problèmes.

\begin{refs}
\item Contexte : Polyomino --- Wikipédia (\url{https://fr.wikipedia.org/wiki/Polyomino})
\item Voir aussi : L'art de la démonstration (\url{proofs.html}), pour l'échiquier mutilé
\end{refs}

\bigskip

\begin{problem}[{Deux brocs et aucune graduation --- Lycée}]
\textbf{PUZ-17.} \textbf{Problème.} Vous disposez de deux brocs sans graduation, de capacités entières $ a $ et $ b $ litres, d'un robinet fournissant de l'eau sans limite, et d'un évier. Les seules opérations autorisées sont : remplir un broc à ras bord, vider complètement un broc, et verser d'un broc dans l'autre jusqu'à ce que la source soit vide ou que la destination soit pleine.
\begin{enumerate}
\item[(a)] Avec $ a = 3 $ et $ b = 5 $, laissez exactement $ 4 $ litres dans l'un des brocs. Écrivez la suite complète des opérations et justifiez chaque étape.
\item[(b)] Déterminez, avec démonstration, exactement quelles quantités $ d $ sont mesurables --- c'est-à-dire pour quels $ d $ une suite finie d'opérations autorisées laisse exactement $ d $ litres dans un broc. Démontrez les deux sens : que tout $ d $ de ce type est atteignable, et qu'aucun autre $ d $ ne l'est.
\item[(c)] La version de Tartaglia n'a ni robinet ni évier, de sorte que le total est conservé : un broc de $ 8 $ litres est plein de vin, et des brocs de $ 5 $ et $ 3 $ litres sont vides. Déterminez, avec démonstration, toutes les répartitions de vin atteignables et, en particulier, décidez si le vin peut être partagé en deux parts égales.
\end{enumerate}

\end{problem}

\noindent Les casse-tête de transvasement sont médiévaux, et le problème du vin $ 8 $--$ 5 $--$ 3 $ de la question (c) est assez ancien pour figurer dans les \textit{Problèmes plaisants et délectables} de Claude Gaspard Bachet de Méziriac, en 1612 ; la version $ 3 $--$ 5 $ de la question (a) a touché un très large public grâce à \textit{Une journée en enfer} en 1995. Les mathématiques derrière la question (b) sont la théorie des équations diophantiennes linéaires : tout état atteignable est une combinaison entière des deux capacités, si bien que la question est de savoir quels entiers de telles combinaisons produisent --- c'est le contenu du théorème de Bachet-Bézout, que Bachet lui-même connaissait dans le cas entier bien avant que Bézout n'attache son nom à la version polynomiale. La question (c) est un problème réellement différent, car la conservation du total vous confine à un ensemble d'états de dimension deux ; M. C. K. Tweedie a montré en 1939 que toute la famille des casse-tête à total conservé cède à une seule image géométrique, bel exemple de changement de représentation rendant inutile une recherche exhaustive.

\begin{refs}
\item Contexte : Casse-tête de transvasement --- Wikipédia (en anglais) (\url{https://en.wikipedia.org/wiki/Water_pouring_puzzle})
\item Contexte : Théorème de Bachet-Bézout --- Wikipédia (\url{https://fr.wikipedia.org/wiki/Th%C3%A9or%C3%A8me_de_Bachet-B%C3%A9zout})
\item Article : M. C. K. Tweedie, \og{} A graphical method of solving Tartaglian measuring puzzles \fg{}, \textit{The Mathematical Gazette} 23 (1939), 278--282
\item Voir aussi : Théorie des nombres (\url{pure-number-theory.html})
\end{refs}

\bigskip

\begin{problem}[{Les missionnaires et les cannibales --- Lycée}]
\textbf{PUZ-18.} \textbf{Problème.} Trois missionnaires et trois cannibales se tiennent sur une rive de rivière. Une barque porte au plus deux personnes et ne peut pas traverser à vide : quelqu'un doit ramer dans chaque sens. Si, à un instant quelconque, les cannibales sont plus nombreux que les missionnaires sur l'une des rives où se trouve au moins un missionnaire, les missionnaires sont mangés.
\begin{enumerate}
\item[(a)] Trouvez une suite de traversées qui amène les six sains et saufs sur l'autre rive.
\item[(b)] Formez le graphe dont les sommets sont les états licites --- le nombre de missionnaires et le nombre de cannibales sur la rive de départ, ainsi que la position de la barque --- et dont les arêtes relient les états qui diffèrent d'une traversée. Démontrez que votre solution est un plus court chemin dans ce graphe, et comptez combien de solutions de longueur minimale il existe.
\item[(c)] Déterminez, avec démonstration, si quatre missionnaires et quatre cannibales peuvent traverser dans une barque de deux places. Déterminez ensuite le plus grand $ n $ pour lequel $ n $ missionnaires et $ n $ cannibales peuvent traverser quand la barque en porte trois.
\end{enumerate}

\end{problem}

\noindent Ce casse-tête descend d'un problème de passage de rivière des \textit{Propositiones ad Acuendos Juvenes}, le recueil attribué à Alcuin d'York vers l'an 800, où il apparaît sous la forme de trois frères voyageant chacun avec une sœur, avec une barque de deux places. Il a circulé dans le nord de l'Europe à partir du treizième siècle comme problème des \og{} maris jaloux \fg{}, et n'a acquis ses missionnaires et ses cannibales qu'à la fin du dix-neuvième. Son importance moderne est autre : Saul Amarel l'a utilisé en 1968 comme exemple central d'un article influent sur la façon dont le choix de la représentation détermine si un problème de recherche est traitable, et il illustre depuis, dans tous les manuels, la recherche dans un espace d'états en intelligence artificielle. La question (c) est là où le casse-tête cesse d'être une histoire pour devenir un théorème --- la capacité de la barque et le nombre de paires interagissent d'une manière qu'Ian Pressman et David Singmaster ont analysée avec soin en 1989.

\begin{refs}
\item Contexte : Problèmes de passage de rivière --- Wikipédia (\url{https://fr.wikipedia.org/wiki/Probl%C3%A8mes_de_passage_de_rivi%C3%A8re})
\item Contexte : Propositiones ad acuendos juvenes --- Wikipédia (\url{https://fr.wikipedia.org/wiki/Propositiones_ad_acuendos_juvenes})
\item Article : Ian Pressman \& David Singmaster, \og{} `The Jealous Husbands' and `The Missionaries and Cannibals' \fg{}, \textit{The Mathematical Gazette} 73 (1989)
\item Voir aussi : Combinatoire et théorie des graphes (\url{pure-combinatorics.html}), Mathématiques discrètes et informatique (\url{applied-discrete-cs.html})
\end{refs}

\bigskip

\begin{problem}[{Une seule lecture sur une balance numérique --- Lycée}]
\textbf{PUZ-19.} \textbf{Problème.} Dix sacs contiennent chacun un grand nombre de pièces. Les pièces authentiques pèsent exactement $ 10 $ grammes ; toutes les pièces d'exactement un sac sont fausses et pèsent $ 9 $ grammes. Votre instrument est une balance numérique qui affiche le poids total exact de ce que vous y posez, et vous ne pouvez l'utiliser qu'une fois.
\begin{enumerate}
\item[(a)] Identifiez le sac de fausses pièces en une seule pesée, et démontrez que votre procédure distingue les dix possibilités.
\item[(b)] Supposons maintenant qu'un nombre inconnu des dix sacs --- peut-être aucun, peut-être tous --- contienne des fausses pièces. Déterminez, avec démonstration, si une pesée suffit encore, et donnez en conséquence soit une procédure, soit un argument d'impossibilité.
\item[(c)] Modifiez l'énoncé pour que ce soient des pièces individuelles, et non des sacs entiers, qui puissent être fausses : parmi $ n $ pièces, un sous-ensemble inconnu est léger, et chaque pesée affiche le poids total exact de tout sous-ensemble que vous choisissez. Supposez que tous les sous-ensembles à peser doivent être fixés à l'avance. Démontrez les meilleures bornes supérieure et inférieure que vous pouvez sur le nombre de pesées nécessaires, en fonction de $ n $.
\end{enumerate}

\end{problem}

\noindent Une balance à deux plateaux renvoie l'un de trois symboles : elle porte donc au plus $ \log_2 3 $ bits par usage ; une balance qui affiche un nombre peut en principe en porter un nombre non borné, et le casse-tête consiste à convertir cette capacité brute en un code utilisable. Les questions (a) et (b) relèvent d'un folklore d'origine incertaine, circulant depuis des décennies comme question d'entretien, mais la question (c) est un vrai problème de recherche : c'est le problème de pesée de pièces posé par Paul Erdős et Alfréd Rényi en 1963, dont la réponse asymptotique a demandé plusieurs années et plusieurs auteurs. La même idée --- regrouper les échantillons pour qu'une mesure résolve d'un coup de nombreux objets --- a été inventée indépendamment par Robert Dorfman en 1943 pour dépister la syphilis chez les soldats américains, et fait aujourd'hui l'objet de la théorie des tests groupés, revenue sur le devant de la scène lors des campagnes de dépistage à grande échelle.

\begin{refs}
\item Contexte : Tests groupés --- Wikipédia (\url{https://fr.wikipedia.org/wiki/Tests_group%C3%A9s})
\item Contexte : Problème de pesée --- Wikipédia (\url{https://fr.wikipedia.org/wiki/Probl%C3%A8me_de_pes%C3%A9e})
\item Article : P. Erdős \& A. Rényi, \og{} On two problems of information theory \fg{}, \textit{Publications of the Mathematical Institute of the Hungarian Academy of Sciences} 8 (1963), 229--243
\item Voir aussi : Cryptographie et théorie de l'information (\url{applied-crypto.html})
\end{refs}

\bigskip

\begin{problem}[{Cent pièces dans le noir --- Lycée, short}]
\textbf{PUZ-20.} \textbf{Problème.} Cent pièces sont posées sur une table dans une pièce entièrement obscure ; exactement dix d'entre elles sont côté pile visible, et vous ne pouvez déterminer par le toucher ni par aucun autre moyen dans quel sens une pièce est tournée, bien que vous puissiez déplacer les pièces et les retourner. Séparez-les en deux groupes, pas nécessairement de même taille, de sorte que les deux groupes contiennent le même nombre de piles, et démontrez que votre procédure fonctionne pour toute disposition initiale possible.
\end{problem}

\noindent Le difficile n'est pas de trouver une manœuvre astucieuse, mais de remarquer quel genre d'objet on vous demande : puisque vous n'apprenez jamais rien sur la disposition, votre procédure est une recette fixe qui doit réussir simultanément pour les $ \binom{100}{10} $ dispositions. Cela vous force à chercher une quantité dont le comportement sous vos coups est le même quoi que vous ne puissiez voir --- un argument d'invariance mené à l'envers, et une bonne première rencontre avec l'idée qu'une démonstration peut maîtriser tous les cas à la fois sans en examiner aucun.

\begin{refs}
\item Contexte : Invariant (mathématiques) --- Wikipédia (\url{https://fr.wikipedia.org/wiki/Invariant})
\item Voir aussi : L'art de la démonstration (\url{proofs.html})
\end{refs}

\bigskip

\begin{problem}[{Le traiteur paresseux et le pâtissier négligent --- Lycée}]
\textbf{PUZ-21.} \textbf{Problème.} Une crêpe est un disque. Une coupe est une droite entière qui la traverse, et les morceaux sont les régions connexes qui restent ; les morceaux ne sont jamais réarrangés entre deux coupes.
\begin{enumerate}
\item[(a)] Déterminez, avec démonstration, le plus grand nombre de morceaux que l'on peut obtenir avec $ n $ coupes, et démontrez votre formule par récurrence sur $ n $.
\item[(b)] Une autre question : marquez $ n $ points sur la circonférence d'un disque, en position générale de sorte que trois des cordes qu'ils déterminent ne passent jamais par un même point intérieur, et tracez les $ \binom{n}{2} $ cordes. Déterminez, avec démonstration, le nombre de régions en lesquelles le disque est découpé, et comparez votre réponse à $ 2^{n-1} $.
\item[(c)] En dimension trois : un gâteau cubique est coupé par $ n $ plans, toujours sans réarrangement. Déterminez, avec démonstration, le plus grand nombre de morceaux, et expliquez le lien entre cette réponse et celle trouvée en (a).
\end{enumerate}

\end{problem}

\noindent La question (a) est la plus ancienne des trois et parle en réalité des arrangements de droites du plan, sujet ouvert par Jakob Steiner en 1826 et aujourd'hui standard en géométrie algorithmique, où la complexité d'un arrangement détermine le temps d'exécution de bien des algorithmes. La question (b) est la configuration dont Leo Moser s'est servi en 1949 pour mettre en garde contre la généralisation à partir d'une poignée de cas calculés ; c'est l'exemple le plus cité en mathématiques d'une récurrence qui n'existe pas, et l'intérêt de l'exercice est d'obtenir le compte par une méthode qui ne repose pas sur le repérage d'un motif. La question (c) appartient à la même famille, et les trois réponses forment ensemble un petit motif de coefficients binomiaux qui dit quelque chose de structurel sur la dimension.

\begin{refs}
\item Contexte : Suite du traiteur paresseux --- Wikipédia (\url{https://fr.wikipedia.org/wiki/Suite_du_traiteur_paresseux})
\item Contexte : Découpage d'un disque en régions --- Wikipédia (en anglais) (\url{https://en.wikipedia.org/wiki/Dividing_a_circle_into_areas})
\item Contexte : Suite des nombres gâteaux --- Wikipédia (\url{https://fr.wikipedia.org/wiki/Suite_des_nombres_g%C3%A2teaux})
\item Voir aussi : Combinatoire et théorie des graphes (\url{pure-combinatorics.html})
\end{refs}

\bigskip

\begin{problem}[{Des tétrominos en T sur un plateau --- Lycée}]
\textbf{PUZ-22.} \textbf{Problème.} Un tétromino en T est le polyomino à quatre cases formé d'une rangée de trois cases avec une quatrième case accolée au milieu de l'un des longs côtés. Les tuiles peuvent être tournées, ne peuvent pas se chevaucher, et doivent recouvrir le plateau exactement, sans laisser de case découverte.
\begin{enumerate}
\item[(a)] Pavez un plateau $ 4 \times 4 $ avec quatre tétrominos en T.
\item[(b)] Démontrez qu'un plateau $ 10 \times 10 $ n'admet aucun pavage par des tétrominos en T, bien que ses $ 100 $ cases soient divisibles par $ 4 $.
\item[(c)] Déterminez, avec démonstration, exactement quels rectangles $ m \times n $ admettent un pavage par des tétrominos en T.
\end{enumerate}

\end{problem}

\noindent Solomon Golomb a forgé le mot \og{} polyomino \fg{} lors d'un exposé de 1953 au Harvard Mathematics Club et a passé les quarante années suivantes à faire de ces formes une discipline ; la question (c) a été réglée par D. W. Walkup en 1965, dans une note de trois pages de l'\textit{American Mathematical Monthly}. Ce qui rend le tétromino en T instructif, c'est que l'obstruction naïve échoue : le compte des cases est correct, et l'argument ordinaire du damier noir et blanc qui tue le plateau mutilé ne tue pas immédiatement celui-ci, de sorte qu'il faut réfléchir plus sérieusement à ce à quoi sert un coloriage. Une démonstration d'impossibilité de ce genre est une affirmation sur une infinité de pavages candidats, établie sans en examiner aucun, et apprendre à fabriquer la bonne pondération est un ajout définitif à la boîte à outils de qui résout des problèmes.

\begin{refs}
\item Contexte : Tétromino --- Wikipédia (\url{https://fr.wikipedia.org/wiki/T%C3%A9tromino})
\item Contexte : Polyomino --- Wikipédia (\url{https://fr.wikipedia.org/wiki/Polyomino})
\item Article : D. W. Walkup, \og{} Covering a rectangle with T-tetrominoes \fg{}, \textit{American Mathematical Monthly} 72 (1965), 986--988
\item Lecture : Solomon W. Golomb, \textit{Polyominoes: Puzzles, Patterns, Problems, and Packings} (2e éd., 1994)
\end{refs}

\bigskip

\begin{problem}[{L'anniversaire de Cheryl --- Lycée}]
\textbf{PUZ-31.} \textbf{Problème.} Albert et Bernard veulent connaître la date d'anniversaire de Cheryl.
Elle leur dit que c'est l'une de dix dates : 15 mai, 16 mai, 19 mai ; 17 juin, 18 juin ;
14 juillet, 16 juillet ; 14 août, 15 août, 17 août. Elle dit ensuite à Albert le mois, en
privé, et à Bernard le quantième, en privé. Chacun sait que l'autre a été informé, et
chacun sait que l'autre raisonne parfaitement. Ce qui suit est dit à voix haute, dans
l'ordre.

\noindent Albert : \og{} Je ne sais pas quand Cheryl est née, et je sais que Bernard ne le sait pas
non plus. \fg{} Bernard : \og{} Au début je ne savais pas, mais maintenant je sais. \fg{} Albert :
\og{} Alors je le sais aussi. \fg{}
\begin{enumerate}
\item[(a)] Déterminez la date d'anniversaire de Cheryl, et justifiez chaque élimination : pour
chacune des dix dates, dites laquelle des trois déclarations l'écarte, et pourquoi.
\item[(b)] Écrivez chaque déclaration comme une condition exacte. La première phrase d'Albert,
par exemple, est un énoncé qui ne porte que sur le mois ; rendez cela précis, et énoncez
ce que la phrase de Bernard affirme au sujet du quantième, compte tenu de ce qu'Albert a
déjà dit.
\item[(c)] Le casse-tête est fragile. Construisez une liste de dates de votre cru sur laquelle
les mêmes trois déclarations ne pourraient pas toutes être faites de façon véridique, puis
une seconde liste sur laquelle elles peuvent toutes être faites tout en laissant deux dates
possibles. Concluez que l'unicité de la réponse est une propriété de la liste particulière
de Cheryl, non du format.
\end{enumerate}

\end{problem}

\noindent Ce problème est paru comme question 24 de l'olympiade mathématique des
écoles de Singapour et d'Asie, le 8 avril 2015, rédigé par le Dr Joseph Yeo Boon Wooi, du
National Institute of Education de Singapour ; un présentateur de télévision, Kenneth Kong,
l'a mis en ligne deux jours plus tard et il a circulé dans le monde entier en une semaine,
en grande partie parce que quantité d'adultes étaient convaincus qu'il ne pouvait pas être
résolu. Il peut l'être, et la raison en est le moteur de toute cette famille de problèmes :
une déclaration d'\textit{ignorance} est informative. Chaque locuteur, en avouant ne pas
savoir, apprend à l'autre que certaines possibilités sont absentes, et les possibilités
sont élaguées tour après tour jusqu'à ce qu'une seule date survive. C'est exactement le
mécanisme des insulaires aux yeux bleus de PUZ-11, mais assez petit pour que l'argument
tout entier tienne sur une page --- ce qui en fait le meilleur premier exercice de
raisonnement sur ce que les autres savent.

\begin{refs}
\item Contexte : Anniversaire de Cheryl --- Wikipédia (\url{https://fr.wikipedia.org/wiki/Anniversaire_de_Cheryl})
\item Contexte : Logique épistémique dynamique --- Wikipédia (en anglais) (\url{https://en.wikipedia.org/wiki/Dynamic_epistemic_logic})
\item Voir aussi : PUZ-11 ci-dessus, et Logique, théorie des ensembles et fondements (\url{pure-foundations.html})
\end{refs}

\bigskip

\begin{problem}[{L'énigme du zèbre --- Lycée}]
\textbf{PUZ-32.} \textbf{Problème.} Cinq maisons s'alignent. Chacune est d'une couleur différente,
habitée par un homme d'une nationalité différente, et chaque homme boit une boisson
différente, fume une marque de cigarettes différente et garde un animal différent. On vous
donne quinze faits.

\noindent (1) L'Anglais habite la maison rouge. (2) L'Espagnol possède le chien. (3) On boit du
café dans la maison verte. (4) L'Ukrainien boit du thé. (5) La maison verte est
immédiatement à droite de la maison ivoire. (6) Le fumeur d'Old Gold possède les
escargots. (7) On fume des Kools dans la maison jaune. (8) On boit du lait dans la maison
du milieu. (9) Le Norvégien habite la première maison. (10) L'homme qui fume des
Chesterfield habite la maison voisine de celle de l'homme au renard. (11) On fume des
Kools dans la maison voisine de celle où l'on garde le cheval. (12) Le fumeur de Lucky
Strike boit du jus d'orange. (13) Le Japonais fume des Parliament. (14) Le Norvégien
habite à côté de la maison bleue. (15) Il y a cinq maisons.
\begin{enumerate}
\item[(a)] Déterminez qui boit de l'eau et qui possède le zèbre, et démontrez que les quinze
faits n'admettent qu'une seule affectation complète.
\item[(b)] Reformulez le casse-tête comme un problème de satisfaction de contraintes : prenez
cinq variables par catégorie (couleur, nationalité, boisson, marque, animal), chacune à
valeurs dans les positions de maison $ \{1,2,3,4,5\} $, imposez une contrainte de
différence deux à deux à l'intérieur de chaque catégorie, et traduisez chaque fait en
contrainte. Montrez que le nombre d'affectations à parcourir aveuglément est
$ (5!)^{5} = 24\,883\,200\,000 $, et décrivez ce que fait gagner l'application locale de
chaque contrainte.
\item[(c)] Ni l'eau ni le zèbre ne sont mentionnés dans aucun des quinze faits, et pourtant le
casse-tête interroge sur les deux. Expliquez pourquoi c'est légitime. Construisez ensuite
un petit casse-tête de même forme --- trois catégories sur trois positions, contraintes de
différence deux à deux, et une poignée de faits \og{} à côté de \fg{} et \og{} même maison \fg{} --- qui
possède exactement deux solutions, puis un autre qui n'en possède aucune.
\end{enumerate}

\end{problem}

\noindent L'énigme a été imprimée dans \textit{Life International} le 17 décembre
1962, la réponse paraissant en mars suivant ; on l'attribue souvent à Einstein ou à Lewis
Carroll, sans le moindre élément de preuve dans l'un ou l'autre cas. Son importance réelle
est d'être devenue l'exemple traité standard du problème de satisfaction de contraintes, le
modèle dans lequel un ensemble fini de variables à domaines finis doit satisfaire des
relations données. Ce cadre couvre le coloriage de graphes, la construction d'emplois du
temps, le sudoku et la vérification de circuits, et il possède une théorie profonde : Feder
et Vardi ont conjecturé dans les années 1990 que tout langage de contraintes fixé est soit
résoluble en temps polynomial, soit NP-complet, sans rien entre les deux, et la conjecture
a été démontrée indépendamment par Andrei Bulatov et Dmitriy Zhuk en 2017. La question (b)
est le moment où une histoire de cigarettes devient cette discipline : une fois les
contraintes écrites, la question \og{} comment explorer cela efficacement ? \fg{} ne parle plus de
zèbres.

\begin{refs}
\item Contexte : Énigme d'Einstein (énigme du zèbre) --- Wikipédia (\url{https://fr.wikipedia.org/wiki/%C3%89nigme_d%27Einstein})
\item Contexte : Problème de satisfaction de contraintes --- Wikipédia (\url{https://fr.wikipedia.org/wiki/Probl%C3%A8me_de_satisfaction_de_contraintes})
\item Lecture : Stuart Russell \& Peter Norvig, \textit{Artificial Intelligence: A Modern Approach}, le chapitre sur les problèmes de satisfaction de contraintes
\item Voir aussi : Mathématiques discrètes et informatique (\url{applied-discrete-cs.html})
\end{refs}

\bigskip

\begin{problem}[{Les vingt questions --- Lycée, short}]
\textbf{PUZ-33.} \textbf{Problème.} Un joueur choisit secrètement un élément d'un
ensemble fini $ S $ connu des deux ; l'autre doit l'identifier en posant des questions de
la forme \og{} votre élément appartient-il à $ A $ ? \fg{} pour toute partie
$ A \subseteq S $ de son choix, chaque réponse étant véridique. Déterminez, avec
démonstration, le plus petit nombre de questions qui suffise toujours, en fonction de
$ |S| $, les questions ultérieures pouvant dépendre des réponses antérieures.
\end{problem}

\noindent Le jeu de société est ancien --- Hannah More raconte y avoir joué lors
d'un dîner londonien dans les années 1780, et son cousin hongrois, le Barkochba, se joue à
Budapest depuis 1911 --- mais les mathématiques sont celles de Shannon. Une réponse par oui
ou par non porte au plus un bit, de sorte que $ q $ questions peuvent séparer au plus
$ 2^{q} $ possibilités ; c'est la même comptabilité que celle qui régit la balance de
PUZ-01, où chaque pesée porte $ \log_2 3 $ bits au lieu d'un. La moitié facile consiste à
exhiber une stratégie. La moitié qui mérite d'être traitée avec soin est la borne
inférieure, car elle doit mettre en échec toutes les stratégies adaptatives à la fois, et
pas seulement les évidentes.

\begin{refs}
\item Contexte : Le jeu des vingt questions --- Wikipédia (en anglais) (\url{https://en.wikipedia.org/wiki/Twenty_questions})
\item Contexte : Entropie de Shannon --- Wikipédia (\url{https://fr.wikipedia.org/wiki/Entropie_de_Shannon})
\item Voir aussi : Cryptographie et théorie de l'information (\url{applied-crypto.html})
\end{refs}

\bigskip

\begin{problem}[{Le problème de Josèphe --- Lycée}]
\textbf{PUZ-34.} \textbf{Problème.} Numérotez $ n $ personnes $ 1, 2, \ldots, n $ dans le sens des
aiguilles d'une montre autour d'un cercle. En commençant à compter à la personne $ 1 $,
une personne sur deux encore dans le cercle est retirée --- la personne $ 2 $ d'abord,
puis $ 4 $, et ainsi de suite tour après tour --- jusqu'à ce qu'il n'en reste qu'une.
Notons $ J(n) $ le numéro de ce survivant.
\begin{enumerate}
\item[(a)] Calculez $ J(n) $ pour $ n = 1, 2, \ldots, 16 $ et décrivez le motif que vous
observez.
\item[(b)] Démontrez une récurrence exprimant $ J(2m) $ et $ J(2m+1) $ en fonction de
$ J(m) $, et servez-vous-en pour obtenir une forme close de $ J(n) $ valable pour tout
$ n $. Démontrez cette forme close par récurrence.
\item[(c)] Retirez maintenant une personne sur $ k $ au lieu d'une sur deux. Démontrez que
\[ J_k(n) = \big( J_k(n-1) + k - 1 \big) \bmod n \; + \; 1, \qquad J_k(1) = 1, \]
et expliquez pourquoi cela donne un algorithme plutôt qu'une formule. Pour quels $ k $
pouvez-vous trouver quoi que ce soit d'aussi net que la réponse à la question (b) ?
\end{enumerate}

\end{problem}

\noindent Flavius Josèphe raconte dans \textit{La Guerre des Juifs} (livre III,
chapitre 8) comment, après la chute de Yodfat en 67, lui et quarante compagnons se sont
cachés dans une grotte et ont convenu de mourir par tirage au sort plutôt que de se rendre
aux Romains --- et comment lui et un autre survécurent, \og{} par chance ou peut-être par la main
de Dieu \fg{}. La formalisation par comptage est postérieure : Bachet de Méziriac en a donné
une version en 1612, et une refonte médiévale place quinze chrétiens et quinze Turcs sur un
navire en train de sombrer, un homme sur neuf étant jeté par-dessus bord. Ce qui fait du cas
$ k=2 $ un petit classique, c'est que la réponse, une fois trouvée, est étonnamment simple
--- assez simple pour que Graham, Knuth et Patashnik ouvrent \textit{Concrete Mathematics} avec
elle, s'en servant pour enseigner comment deviner un motif à partir de données, le démontrer
par récurrence, puis voir pourquoi la démonstration fonctionne. La question (c) est le
contrepoids honnête : le cas général $ k $ est calculable mais obstinément informe.

\begin{refs}
\item Contexte : Problème de Josèphe --- Wikipédia (\url{https://fr.wikipedia.org/wiki/Probl%C3%A8me_de_Jos%C3%A8phe})
\item Lecture : Ronald Graham, Donald Knuth \& Oren Patashnik, \textit{Concrete Mathematics}, \S{}1.3
\item Voir aussi : Combinatoire et théorie des graphes (\url{pure-combinatorics.html}), L'art de la démonstration (\url{proofs.html})
\end{refs}

\bigskip

\begin{problem}[{Les deux enveloppes, scellées --- Lycée}]
\textbf{PUZ-45.} \textbf{Problème.} Deux enveloppes scellées contiennent de l'argent ; l'une en contient
exactement deux fois plus que l'autre. Vous en prenez une au hasard et, sans l'ouvrir, on
vous propose d'échanger.
\begin{enumerate}
\item[(a)] Voici un argument en faveur de l'échange. Soit $ A $ la somme contenue dans votre
enveloppe. L'autre enveloppe contient $ 2A $ ou $ A/2 $, chacun avec probabilité
$ 1/2 $, de sorte que son contenu espéré vaut
\[ \tfrac{1}{2}(2A) + \tfrac{1}{2}\!\left(\tfrac{A}{2}\right) = \tfrac{5}{4}A > A . \]
Le même argument s'applique à l'autre enveloppe, si bien que chacune vaut plus que l'autre.
Localisez l'étape exacte à partir de laquelle ceci cesse d'être une probabilité valide, et
justifiez votre verdict en nommant un espace probabilisé, les variables aléatoires qui y
sont définies, et la quantité sur laquelle on conditionne réellement.
\item[(b)] Construisez un modèle honnête : soit $ X $ la variable aléatoire donnant la plus
petite somme, et supposons que les deux enveloppes contiennent $ X $ et $ 2X $.
Démontrez que si $ E[X] $ est finie, le gain espéré de l'échange est nul.
\item[(c)] Montrez que la question (b) ne liquide pas entièrement le casse-tête. Construisez une
loi pour $ X $ telle que, après avoir ouvert votre enveloppe et vu la somme $ a $,
l'espérance conditionnelle de l'autre enveloppe dépasse strictement $ a $ --- et cela pour
toute valeur de $ a $ que vous pourriez voir. Démontrez ensuite que toute loi ayant cette
propriété vérifie $ E[X] = \infty $, et comparez la situation au paradoxe de
Saint-Pétersbourg.
\end{enumerate}

\end{problem}

\noindent La forme la plus ancienne en est le paradoxe des cravates de Maurice
Kraitchik, en 1943, où deux hommes soutiennent chacun que la cravate de l'autre a plus de
chances d'être la plus chère ; la version aux cartes d'Erwin Schrödinger est parvenue à
l'imprimé via le \textit{Miscellany} de Littlewood en 1953, Martin Gardner a popularisé un jeu
de portefeuilles en 1958 en admettant qu'il ne savait pas dire proprement ce qui clochait,
et Barry Nalebuff a donné la formulation moderne aux enveloppes en 1989. Ce qui vaut à ce
casse-tête sa longue vie, c'est la question (c) : l'argument naïf n'est pas seulement
bâclé, il est \textit{presque} réalisable, et les lois qui le réalisent sont exactement celles
d'espérance infinie, là où la valeur espérée cesse d'être un guide pour l'action. La
littérature n'est toujours pas unanime sur la phrase du sophisme qui est la coupable, bien
que les mathématiques soient entièrement réglées dès qu'un modèle est fixé --- bonne
illustration de la différence entre une question de probabilité et une question de
mots.

\begin{refs}
\item Contexte : Paradoxe des deux enveloppes --- Wikipédia (\url{https://fr.wikipedia.org/wiki/Paradoxe_des_deux_enveloppes})
\item Contexte : Paradoxe de Saint-Pétersbourg --- Wikipédia (\url{https://fr.wikipedia.org/wiki/Paradoxe_de_Saint-P%C3%A9tersbourg})
\item Article : Barry Nalebuff, \og{} Puzzles: The other person's envelope is always greener \fg{}, \textit{Journal of Economic Perspectives} 3 (1989), 171--181
\item Voir aussi : Probabilités (\url{applied-probability.html}), pour le paradoxe de Saint-Pétersbourg traité en entier
\item Voir aussi : PUZ-41 (\url{#PUZ-41}), la version dans laquelle vous ouvrez d'abord l'enveloppe --- là où le paradoxe devient réellement subtil.
\end{refs}

\bigskip

\begin{problem}[{Le jeu de Penney --- Lycée}]
\textbf{PUZ-46.} \textbf{Problème.} Le joueur A annonce une suite de trois résultats de pile ou face --- disons
PFP --- et la montre au joueur B, qui annonce alors une autre suite de trois. Une pièce
équilibrée est lancée de façon répétée jusqu'à ce que l'une des deux suites apparaisse en
trois lancers consécutifs ; celui qui l'a annoncée gagne.
\begin{enumerate}
\item[(a)] Démontrez que B, qui joue en second, peut toujours répondre par une suite qui apparaît
la première avec une probabilité strictement supérieure à $ 1/2 $, quel qu'ait été le
choix de A. Déterminez la meilleure réponse à chacune des huit suites de longueur trois, et
démontrez chacune de vos huit affirmations.
\item[(b)] Déduisez-en que \og{} bat \fg{} n'est pas une relation transitive sur les suites : exhibez
$ S_1, S_2, S_3, S_4 $ de longueur trois, chacune battant la suivante et $ S_4 $ battant
$ S_1 $. Démontrez ensuite qu'il n'existe aucun cycle de longueur trois parmi les huit
suites, de sorte que quatre est le plus court cycle disponible.
\item[(c)] Calculez deux fois l'une des probabilités de victoire, par des voies véritablement
différentes, et vérifiez que les réponses concordent : une fois par analyse au premier pas
sur l'automate fini dont les états enregistrent quelle portion de chaque suite est
actuellement appariée, et une fois par un argument de martingale du type employé pour les
temps d'attente de motifs.
\item[(d)] Montrez que le phénomène ne porte pas seulement sur les temps d'attente. Démontrez que,
parmi les suites de longueur trois, aucun renversement ne se produit : si $ S $ a un temps
d'attente espéré strictement plus court que $ T $ considéré isolément, alors $ T $ ne
gagne pas la course en face-à-face. Trouvez ensuite des suites $ S $ et $ T $ de
longueur \textit{quatre} pour lesquelles $ S $ a le temps d'attente espéré le plus court et
pourtant $ T $ apparaît la première avec une probabilité strictement supérieure à
$ 1/2 $, et expliquez quelle quantité, autre que les deux temps d'attente, décide de la
course.
\end{enumerate}

\end{problem}

\noindent Walter Penney a publié ce jeu dans le \textit{Journal of Recreational
Mathematics} en octobre 1969, et Martin Gardner l'a répandu dans \textit{Scientific
American} comme fleuron de ses paradoxes non transitifs. Son charme est d'être un jeu à
pièce équilibrée dans lequel le second joueur possède un avantage décisif, ce qui heurte
l'instinct même que heurtent les dés intransitifs : nous attendons de \og{} a plus de chances de
gagner contre \fg{} que ce soit un ordre, et ce n'en est pas un. Les mathématiques
sous-jacentes tiennent à la corrélation entre motifs qui se chevauchent --- PPP porte PP à la
fois comme préfixe et comme suffixe, et se fait donc concurrence à elle-même, alors que PPF
non --- et cette structure de chevauchement est exactement ce que l'algorithme du nombre
directeur de Conway condense en une formule. La question (d) est la forme la plus nette de
la surprise : temps d'attente espéré et victoire en face-à-face sont des comparaisons
différentes, et une suite peut gagner l'une et perdre l'autre.

\begin{refs}
\item Contexte : Jeu de Penney --- Wikipédia (en anglais) (\url{https://en.wikipedia.org/wiki/Penney%27s_game})
\item Contexte : Jeu intransitif --- Wikipédia (en anglais) (\url{https://en.wikipedia.org/wiki/Intransitive_game})
\item Lecture : Martin Gardner, \textit{The Colossal Book of Mathematics} (chapitre sur les paradoxes non transitifs)
\item Voir aussi : Probabilités (\url{applied-probability.html}), pour les temps d'attente de motifs par arrêt optionnel, et pour les dés intransitifs
\end{refs}

\bigskip

\begin{problem}[{Quatorze, quinze, et le coup qui ne vient jamais --- Lycée, short}]
\textbf{PUZ-47.} \textbf{Problème.} Le taquin contient des tuiles numérotées de $ 1 $ à
$ 15 $ dans un cadre $ 4 \times 4 $ comportant une case vide, et un coup fait glisser
orthogonalement une tuile dans la case vide. Démontrez qu'en partant des tuiles rangées dans
l'ordre, la case vide en dernier, aucune suite de coups n'atteint jamais la position où
$ 14 $ et $ 15 $ sont échangés et où tout le reste est à sa place initiale.
\end{problem}

\noindent Noyes Palmer Chapman, receveur des postes à Canastota, dans l'État de New
York, semble en avoir fabriqué la première version vers 1874 ; Sam Loyd s'en est attribué
l'invention à partir de 1891 et a offert un prix de mille dollars pour une solution à la
position où précisément $ 14 $ et $ 15 $ sont permutés, laquelle est impossible ---
l'argent n'a jamais couru le moindre risque. Woolsey Johnson et William Story l'avaient déjà
démontré en 1879, dans l'une des toutes premières applications de la signature d'une
permutation à une question extérieure à l'algèbre : exactement la moitié des $ 16! $
dispositions sont atteignables, et les deux moitiés ne se rencontrent jamais. L'outil à
saisir est l'homomorphisme $ \operatorname{sgn} : S_{16} \to \{\pm 1\} $, assorti de
l'observation qu'un simple glissement fait quelque chose de prévisible à la position de la
case vide autant qu'à la permutation.

\begin{refs}
\item Contexte : Taquin --- Wikipédia (\url{https://fr.wikipedia.org/wiki/Taquin})
\item Contexte : Signature d'une permutation --- Wikipédia (\url{https://fr.wikipedia.org/wiki/Signature_d%27une_permutation})
\item Article : Wm. W. Johnson \& W. E. Story, \og{} Notes on the ``15'' puzzle \fg{}, \textit{American Journal of Mathematics} 2 (1879), 397--404
\item Article : Aaron F. Archer, \og{} A modern treatment of the 15 puzzle \fg{}, \textit{American Mathematical Monthly} 106 (1999), 793--799
\end{refs}

\bigskip

\begin{problem}[{Le jeu de Nim, premier jeu jamais résolu --- Lycée}]
\textbf{PUZ-48.} \textbf{Problème.} Le Nim se joue avec plusieurs tas de jetons. Un coup retire un nombre
strictement positif de jetons d'un seul tas, et le joueur qui prend le tout dernier jeton
gagne. Pour des tas de tailles $ a_1, \dots, a_k $, définissez la \textit{somme de Nim}
$ a_1 \oplus \cdots \oplus a_k $ comme le résultat de l'écriture des tailles en binaire et
de l'addition des colonnes modulo $ 2 $, sans retenue.
\begin{enumerate}
\item[(a)] Démontrez le théorème de Bouton : le joueur qui doit jouer perd sous jeu parfait si et
seulement si $ a_1 \oplus \cdots \oplus a_k = 0 $. Votre démonstration doit également
produire le coup gagnant.
\item[(b)] Déduisez-en le cas à deux tas, et énoncez la stratégie obtenue en une phrase qu'un
enfant pourrait suivre et utiliser.
\item[(c)] Résolvez le jeu de soustraction à un tas où un coup retire $ 1 $, $ 2 $ ou $ 3 $
jetons : déterminez quelles tailles de tas sont perdantes pour le joueur qui doit jouer,
démontrez-le, et généralisez au retrait d'entre $ 1 $ et $ m $ jetons.
\item[(d)] Nim misère : c'est maintenant le joueur qui prend le dernier jeton qui \textit{perd}.
Déterminez les positions perdantes, démontrez votre réponse, et identifiez précisément
quelle étape de votre argument de la question (a) s'effondre et où les deux jeux commencent
à diverger.
\end{enumerate}

\end{problem}

\noindent Charles Bouton, à Harvard, a publié la théorie complète dans les
\textit{Annals of Mathematics} en 1901-1902, et a donné son nom au jeu ; des variantes se
jouaient depuis très longtemps, et le jeu ressemble de près au \textit{jiǎn shízǐ} chinois,
\og{} ramasser des pierres \fg{}, bien que l'origine ne soit pas établie. L'article de Bouton
compte bien au-delà du jeu : c'est la première fois que quelqu'un prend un jeu à deux
joueurs de pure adresse et en produit une stratégie complète, démontrable et calculable, et
il l'a fait en trouvant un invariant arithmétique --- la somme de Nim --- qu'aucun coup ne peut
préserver à moins que celui qui joue ne le veuille. L'idée s'est révélée si féconde que tout
jeu impartial se révèle être un tas de Nim déguisé, ce qui fait le contenu de PUZ-52.
Westinghouse a construit le Nimatron, une machine à relais qui jouait au Nim contre les
visiteurs de l'Exposition universelle de New York, l'une des toutes premières machines à
jouer jamais construites.

\begin{refs}
\item Contexte : Jeux de Nim --- Wikipédia (\url{https://fr.wikipedia.org/wiki/Jeux_de_Nim})
\item Contexte : Nimber --- Wikipédia (\url{https://fr.wikipedia.org/wiki/Nimber})
\item Article : C. L. Bouton, \og{} Nim, a game with a complete mathematical theory \fg{}, \textit{Annals of Mathematics} (2) 3 (1901--1902), 35--39
\item Lecture : Elwyn Berlekamp, John Conway \& Richard Guy, \textit{Winning Ways for Your Mathematical Plays}, vol. 1
\end{refs}

\bigskip

\begin{problem}[{La barque, l'ancre et le niveau du bassin --- Lycée}]
\textbf{PUZ-59.} \textbf{Problème.} Une petite barque flotte dans un bassin fermé de section horizontale
$ A $ ; aucune eau n'entre ni ne sort. Dans la barque repose une ancre en fer de masse
$ m $ et de masse volumique $ \rho_a $, supérieure à la masse volumique $ \rho_w $ de
l'eau. Le batelier soulève l'ancre par-dessus bord ; elle coule et vient reposer au fond du
bassin.
\begin{enumerate}
\item[(a)] Déterminez, avec démonstration, si le niveau de l'eau dans le bassin monte, descend ou
reste le même, et donnez exactement la variation de niveau en fonction de $ m $,
$ \rho_a $, $ \rho_w $ et $ A $.
\item[(b)] Refaites le calcul lorsque l'objet jeté par-dessus bord est un bloc de liège de masse
$ m $ et de masse volumique inférieure à $ \rho_w $, qui s'éloigne en flottant à la
surface. Démontrez votre réponse.
\item[(c)] Refaites-le encore lorsque l'ancre n'est pas lâchée mais descendue au bout d'une corde
et laissée suspendue à la barque, sans toucher le fond. Déterminez le niveau dans ce cas et
démontrez-le.
\item[(d)] Enfin, la barque prend l'eau et coule au fond avec l'ancre encore à bord. Déterminez la
variation de niveau, puis énoncez un principe unique dont découlent les quatre réponses,
avec une démonstration que c'est bien le cas.
\end{enumerate}

\end{problem}

\noindent Archimède a écrit \textit{Des corps flottants} vers 250 av. J.-C., et c'est
le plus ancien ouvrage conservé qui mérite le nom de physique mathématique : il part d'un
postulat sur la pression dans un fluide et en déduit, comme théorèmes, les conditions dans
lesquelles les corps flottent, coulent et reposent stablement. La proposition 5 du livre I ---
qu'un corps flottant déplace son propre poids de fluide --- est tout ce casse-tête, et la
difficulté n'est pas la physique mais la comptabilité, car \og{} déplace son propre poids \fg{} et
\og{} déplace son propre volume \fg{} sont des énoncés différents et le casse-tête passe de l'un à
l'autre sans prévenir. Les architectes navals appellent encore \textit{déplacement} la masse
d'un navire, précisément pour cette raison. Un lecteur qui a terminé les quatre questions
devrait pouvoir dire aussitôt, et démontrer, ce qu'il advient du niveau de la mer quand la
glace de mer flottante fond, et pourquoi une calotte reposant sur la terre ferme est une
autre question --- ce qui n'est pas une curiosité mais l'un des calculs porteurs de la science
du climat.

\begin{refs}
\item Contexte : Poussée d'Archimède --- Wikipédia (\url{https://fr.wikipedia.org/wiki/Pouss%C3%A9e_d%27Archim%C3%A8de})
\item Contexte : Flottabilité --- Wikipédia (\url{https://fr.wikipedia.org/wiki/Flottabilit%C3%A9})
\item Lecture : Archimède, \textit{Des corps flottants}, livre I --- dans T. L. Heath (dir.), \textit{The Works of Archimedes} (1897)
\item Lecture : Jearl Walker, \textit{The Flying Circus of Physics} --- un livre entier de questions de cette forme
\end{refs}

\bigskip

\begin{problem}[{Le singe et le chasseur --- Lycée}]
\textbf{PUZ-60.} \textbf{Problème.} Une sarbacane est pointée \textit{directement} sur un singe suspendu à une
branche, à une distance horizontale $ d $ et à une hauteur $ h $ au-dessus de la bouche
du canon. À l'instant où la fléchette quitte le canon à la vitesse $ v_0 $, le singe lâche
prise et tombe. La pesanteur est uniforme, d'intensité $ g $, et le sol se trouve à une
hauteur $ H $ sous la branche.
\begin{enumerate}
\item[(a)] En travaillant dans le référentiel du sol, résolvez les deux trajectoires et déterminez,
avec démonstration, l'ensemble exact des vitesses initiales $ v_0 $ pour lesquelles la
fléchette atteint le singe avant que le singe n'atteigne le sol. Exprimez la condition sous
forme d'inégalité en $ d $, $ h $, $ H $ et $ g $.
\item[(b)] Passez maintenant à un référentiel lui-même en chute libre, et obtenez le même résultat
une seconde fois. Identifiez précisément l'étape où votre argument utilise l'égalité des
masses inerte et gravitationnelle, et dites ce qui changerait si la fléchette et le singe
avaient des rapports différents entre les deux.
\item[(c)] Rétablissez la résistance de l'air sous la forme d'un frottement linéaire, de sorte que
chaque corps subisse une accélération $ -b\vec{v} $ avec la même constante $ b $.
Déterminez si la conclusion de la question (a) survit, et démontrez votre réponse.
\item[(d)] Remplacez le frottement linéaire par un frottement quadratique, une accélération
$ -c\,|\vec{v}|\,\vec{v} $ avec la même constante $ c $ pour les deux corps. Montrez par
un argument explicite que le raisonnement de la question (b) échoue désormais, et calculez
la distance de manqué au premier ordre en $ c $.
\end{enumerate}

\end{problem}

\noindent C'est la plus ancienne démonstration de cours de physique encore
pratiquée : un électroaimant retient un singe en peluche, la sarbacane coupe le circuit en
tirant, et toute la salle regarde. Son contenu est la décomposition galiléenne du mouvement
des projectiles dans les \textit{Discorsi} de 1638 en un mouvement horizontal uniforme composé
avec un mouvement vertical uniformément accéléré --- la première fois que quelqu'un traitait
une trajectoire courbe comme la somme de deux trajectoires simples. La question (b) est une
miniature de ce qu'Einstein appelait en 1907 la pensée la plus heureuse de sa vie : un
observateur en chute libre ne sent pas son propre poids, et la gravitation peut donc être
localement éliminée par changement de référentiel. Les questions (c) et (d) sont là parce
que l'astuce de la chute libre est souvent enseignée comme si elle était inconditionnelle,
et elle ne l'est pas : elle survit exactement aux forces qui dépendent des deux corps de la
même manière, et il revient au lecteur de découvrir par lui-même quelle loi de frottement
qualifie. Ce site ne vous dira pas ce que montre la démonstration ; demandez à un professeur
de physique de la faire, et prédisez le résultat par écrit d'abord.

\begin{refs}
\item Contexte : Le singe et le chasseur --- Wikipédia (en anglais) (\url{https://en.wikipedia.org/wiki/Monkey_and_hunter})
\item Contexte : Principe d'équivalence --- Wikipédia (\url{https://fr.wikipedia.org/wiki/Principe_d%27%C3%A9quivalence})
\item Cours : MIT OCW 8.01SC --- Classical Mechanics (\url{https://ocw.mit.edu/courses/8-01sc-classical-mechanics-fall-2016/})
\item Voir aussi : Physique mathématique (\url{applied-physics.html}), pour les nombres impairs de Galilée et la naissance de la cinématique
\end{refs}

\bigskip

\begin{problem}[{La bicyclette tirée en arrière par sa pédale basse --- Lycée, short}]
\textbf{PUZ-61.} \textbf{Problème.} Une bicyclette est debout, maintenue verticale mais
libre de rouler, une pédale au point le plus bas de son cercle ; une corde attachée à cette
pédale est tirée horizontalement vers l'arrière, assez doucement pour que les deux roues
roulent sans glisser et que la machine ne bascule jamais. Déterminez, avec démonstration,
dans quel sens la bicyclette se déplace, et trouvez la condition exacte, portant sur la
longueur de manivelle $ r $, le rayon $ R $ de la roue motrice et le rapport de
transmission $ g $ --- tours de pédale par tour de roue ---, qui décide de la réponse.
\end{problem}

\noindent D. E. Daykin, à Reading, a mis la question par écrit sous le titre \og{} The
bicycle problem \fg{} dans \textit{Mathematics Magazine} en janvier 1972 ; Martin Gardner l'a
reprise dans \textit{Mathematical Carnival} (1975) et a rapporté que ses lecteurs se
divisaient nettement, plusieurs d'entre eux ne changeant d'avis qu'après être allés à
l'atelier pour essayer. L'outil qui tranche est l'axe instantané de rotation d'une roue qui
roule --- le point du solide qui est momentanément au repos --- d'où la vitesse de tout autre
point découle par un simple produit vectoriel ; de manière équivalente, on peut se demander
quelle courbe la pédale décrit par rapport au sol pendant que la bicyclette roule, à savoir
une trochoïde, et si cette courbe recule jamais. La condition demandée dans la seconde
moitié de la question est le contenu véritable : le casse-tête a un seuil, et une bicyclette
dont le braquet le dépasse ne se comporte pas comme une bicyclette dont le braquet lui est
inférieur, raison pour laquelle discuter de \og{} la \fg{} réponse sans nommer un braquet, c'est
discuter de rien.

\begin{refs}
\item Contexte : Axe instantané de rotation --- Wikipédia (\url{https://fr.wikipedia.org/wiki/Axe_instantan%C3%A9_de_rotation})
\item Contexte : Trochoïde --- Wikipédia (\url{https://fr.wikipedia.org/wiki/Trocho%C3%AFde})
\item Article : D. E. Daykin, \og{} The bicycle problem \fg{}, \textit{Mathematics Magazine} 45 (1972), no 1, p. 1
\item Lecture : Martin Gardner, \textit{Mathematical Carnival} (1975), p. 182
\end{refs}

\bigskip

\begin{problem}[{Le slinky qui tombe --- Lycée}]
\textbf{PUZ-62.} \textbf{Problème.} Un slinky de masse totale $ M $ et de constante de raideur $ k $ pend
au repos par sa spire supérieure, étiré par son propre poids, et est lâché sans vitesse
initiale à $ t = 0 $. Modélisez-le par un ressort sans masse de constante $ k $ portant
une masse répartie uniformément selon une coordonnée $ s \in [0,1] $ qui étiquette les
spires de haut en bas.
\begin{enumerate}
\item[(a)] Trouvez la configuration d'équilibre avant le lâcher : la position de la spire $ s $
en fonction de $ s $, et l'allongement statique total en fonction de $ M $, $ g $ et
$ k $.
\item[(b)] Démontrez qu'à partir de l'instant du lâcher le centre de masse obéit exactement à
\[ x_{\mathrm{cm}}(t) \;=\; x_{\mathrm{cm}}(0) + \tfrac{1}{2} g t^{2} \]
et déduisez-en une contrainte rigoureuse reliant le mouvement de la spire supérieure à celui
de la spire inférieure.
\item[(c)] Écrivez l'équation du mouvement des perturbations longitudinales de ce modèle,
identifiez la vitesse à laquelle elles se propagent, et calculez le temps qu'une
perturbation met pour aller de la spire supérieure à la spire inférieure. Comparez ce temps
à celui que le slinky met pour tomber de sa propre longueur étirée initiale.
\item[(d)] Déterminez, avec démonstration, ce que fait la spire inférieure entre l'instant du
lâcher et l'instant où la spire supérieure la rejoint. Énoncez exactement sur quelle
idéalisation du modèle repose votre conclusion, et décrivez en quoi un slinky réel --- dont
les spires se heurtent et ne peuvent pas se traverser --- doit s'écarter du modèle.
\end{enumerate}

\end{problem}

\noindent Richard James, ingénieur naval travaillant sur des suspensions à ressort
pour instruments de bord, a fait tomber un ressort de torsion d'une étagère en 1943 et l'a
regardé marcher ; lui et sa femme Betty en ont vendu quatre cents en quatre-vingt-dix
minutes chez Gimbels, à Philadelphie, en 1945. La version en chute a été analysée par M. G.
Calkin dans l'\textit{American Journal of Physics} en 1993, puis de nouveau, avec vidéo à
haute cadence et un bien meilleur traitement des spires, par R. C. Cross et M. S. Wheatland
dans la même revue en 2012 --- leur apport est que les spires situées derrière la perturbation
mettent un temps fini à se refermer, là où le modèle antérieur rendait cela instantané. Ce
qui vaut à cette question une page ici, c'est que deux arguments complètement différents s'y
appliquent et doivent être conciliés : le théorème du centre de masse, qui est exact et dit
quelque chose de l'objet entier, et l'équation des ondes, qui est approchée et dit quelque
chose des parties. La question (d) est là où ils se rencontrent. Prédisez la réponse, puis
filmez un slinky avec un téléphone en cadence rapide ; l'équipement coûte quelques euros et
l'expérience prend un après-midi.

\begin{refs}
\item Contexte : Slinky --- Wikipédia (\url{https://fr.wikipedia.org/wiki/Slinky})
\item Contexte : Équation des ondes --- Wikipédia (\url{https://fr.wikipedia.org/wiki/%C3%89quation_des_ondes})
\item Article : M. G. Calkin, \og{} Motion of a falling spring \fg{}, \textit{American Journal of Physics} 61 (1993), 261--264
\item Article : R. C. Cross \& M. S. Wheatland, \og{} Modeling a falling slinky \fg{}, \textit{American Journal of Physics} 80 (2012), 1051--1060 --- arXiv:1208.4629 (\url{https://arxiv.org/abs/1208.4629}) (attention : le résumé donne la réponse à la question (d))
\end{refs}

\bigskip

\begin{problem}[{L'oiseau buveur, ou comment encadrer une estimation --- Lycée}]
\textbf{PUZ-63.} \textbf{Problème.} L'oiseau buveur est un jouet de verre scellé : une tête recouverte de
feutre, un bulbe formant le corps, et un tube qui les relie, contenant du dichlorométhane et
sa propre vapeur, et rien d'autre. Mouillez la tête de feutre, posez l'oiseau à côté d'un
verre d'eau, et il bascule en avant, trempe son bec, se redresse, et recommence --- pendant
des jours.
\begin{enumerate}
\item[(a)] Identifiez la source chaude et la source froide du cycle et exprimez chaque température
en fonction de grandeurs que vous pourriez mesurer dans la pièce. Établissez la borne de
Carnot sur le rendement de toute machine fonctionnant entre elles, et évaluez-la pour une
humidité intérieure typique.
\item[(b)] Estimez la puissance mécanique que délivre le jouet. Choisissez vos données d'entrée ---
masse, géométrie du pivot, amplitude, période --- et, pour chacune, donnez un intervalle que
vous êtes prêt à défendre plutôt qu'un nombre unique ; propagez les intervalles dans le
calcul et donnez la puissance en watts avec des bornes explicites.
\item[(c)] Estimez le débit d'évaporation de l'eau à la tête, convertissez-le en puissance
thermique, et bornez ainsi le rendement réel du jouet. Comparez avec la question (a),
donnez le rapport avec son incertitude, et identifiez la donnée d'entrée unique qui domine
la largeur de votre intervalle.
\item[(d)] Prédisez, avec justification, ce que fait l'oiseau dans chacune des situations
suivantes, et concevez une expérience qui distinguerait votre prédiction de sa négation :
enfermé dans un bocal d'air déjà saturé de vapeur d'eau ; dans un fort courant d'air ; avec
le verre d'eau retiré ; dans une enceinte à vide.
\end{enumerate}

\end{problem}

\noindent Miles V. Sullivan, chimiste aux Bell Telephone Laboratories, a breveté la
forme moderne en 1946 (brevet américain 2 402 463), et Arthur M. Hillery en avait breveté
une version l'année précédente ; mais le jouet est bien plus ancien que l'un et l'autre, et
l'on a montré à Albert et Elsa Einstein un \og{} petit oiseau insatiable \fg{} chinois à leur
arrivée à Shanghai en 1922. L'oiseau est le meilleur argument de bureau qu'une machine
thermique a besoin d'une \textit{différence} de température, et non simplement de chaleur, et
trouver d'où vient cette différence --- elle n'est fournie par rien que l'on branche --- est
tout le casse-tête. Les questions (b) et (c) forment un problème de Fermi au sens strict où
Enrico Fermi l'entendait. On se souvient de lui pour avoir demandé combien d'accordeurs de
piano travaillent à Chicago, mais ce qu'il enseignait réellement était la discipline du
report des bornes : donnez un intervalle pour chaque entrée, propagez-les, et sachez ensuite
dire quelle entrée vous devriez aller mesurer sérieusement. À Trinity, en juillet 1945, il a
laissé tomber des bouts de papier au passage de l'onde de choc et a estimé sur-le-champ la
puissance de la première bombe atomique, à un facteur près que les mesures instrumentées ont
ensuite respecté.

\begin{refs}
\item Contexte : Oiseau buveur --- Wikipédia (\url{https://fr.wikipedia.org/wiki/Oiseau_buveur})
\item Contexte : Estimation de Fermi --- Wikipédia (\url{https://fr.wikipedia.org/wiki/Estimation_de_Fermi})
\item Contexte : Cycle de Carnot --- Wikipédia (\url{https://fr.wikipedia.org/wiki/Cycle_de_Carnot})
\item Voir aussi : Physique mathématique (\url{applied-physics.html}), pour la thermodynamique derrière la borne de Carnot
\end{refs}

\bigskip

\begin{problem}[{Le cube peint --- Lycée}]
\textbf{PUZ-73.} \textbf{Problème.} Un cube plein de côté $ n \ge 2 $ est construit à partir de
$ n^{3} $ cubes unités. Toute sa surface extérieure est peinte, puis il est de nouveau
démonté en cubes unités.
\begin{enumerate}
\item[(a)] Déterminez, avec démonstration, combien de petits cubes portent de la peinture sur
exactement trois faces, sur exactement deux, sur exactement une, et sur aucune, chacun en
fonction de $ n $. Vérifiez que vos quatre décomptes totalisent $ n^{3} $.
\item[(b)] Démontrez algébriquement l'identité
\[ n^{3} \;=\; (n-2)^{3} + 6(n-2)^{2} + 12(n-2) + 8 \]
puis dites exactement quel trait du cube compte chacun des quatre termes. Généralisez à
$ d $ dimensions : démontrez que le nombre de cellules unités du $ d $-cube de côté
$ n $ qui touchent le bord exactement selon une face de dimension $ k $ vaut
$ 2^{\,d-k}\binom{d}{k}(n-2)^{k} $, et déduisez-en l'identité
\[ n^{d} \;=\; \sum_{k=0}^{d} \binom{d}{k} 2^{\,d-k} (n-2)^{k} . \]
\item[(c)] Une autre question sur le même objet. Une coupe est une unique tranche plane à travers
tout ce qui se trouve alors devant la lame ; entre deux coupes, les morceaux peuvent être
réempilés à votre guise. Déterminez, avec démonstration, le plus petit nombre de coupes qui
divise un cube $ 3 \times 3 \times 3 $ en ses $ 27 $ cubes unités, puis déterminez ce
plus petit nombre pour un cube $ n \times n \times n $ général.
\end{enumerate}

\end{problem}

\noindent Les questions (a) et (b) sont le même énoncé raconté deux fois, et c'est
tout l'intérêt de les mettre côte à côte : un exercice de comptage qu'un enfant peut faire
en regardant un cube se révèle être la décomposition en couches d'un hypercube, où le bord
du $ d $-cube se scinde en $ 2^{\,d-k}\binom{d}{k} $ faces de chaque dimension $ k $.
Ludwig Schläfli a dégagé cette structure combinatoire pour les polytopes en toutes
dimensions entre 1850 et 1852, dans un manuscrit si en avance sur son public qu'il n'a été
publié intégralement qu'en 1901, neuf ans après sa mort ; l'identité binomiale de la
question (b) est l'ombre arithmétique de sa classification. La question (c) est d'une tout
autre nature et forme la moitié la plus difficile. Produire une suite courte de coupes est
affaire d'essais, mais montrer qu'aucune suite n'est plus courte oblige à maîtriser d'un
coup tout réempilement possible, et ce genre de minoration se démontre presque toujours en
trouvant une caractéristique du travail achevé dont une seule coupe ne peut fournir qu'une
quantité bornée. Notez ce qui change d'une question à l'autre : (a) et (b) portent sur un
objet fixé, tandis que (c) porte sur une recherche non bornée parmi des procédures, et c'est
la seconde qui a fait de l'optimisation combinatoire une discipline.

\begin{refs}
\item Contexte : Hypercube --- Wikipédia (\url{https://fr.wikipedia.org/wiki/Hypercube})
\item Contexte : Formule du binôme de Newton --- Wikipédia (\url{https://fr.wikipedia.org/wiki/Formule_du_bin%C3%B4me_de_Newton})
\item Lecture : H. S. M. Coxeter, \textit{Regular Polytopes} (3e éd., Dover, 1973), chap. VII sur Schläfli
\item Voir aussi : Combinatoire et théorie des graphes (\url{pure-combinatorics.html}), Géométrie (\url{pure-geometry.html})
\end{refs}

\bigskip

\begin{problem}[{Les poignées de main, et ce qu'une soirée ne peut pas faire --- Lycée}]
\textbf{PUZ-74.} \textbf{Problème.} Lors d'une réunion, certaines paires de personnes se serrent la main.
Personne ne se serre la main à soi-même, et aucune paire de personnes ne se serre la main
plus d'une fois.
\begin{enumerate}
\item[(a)] Démontrez que le nombre de personnes qui serrent un nombre impair de mains est pair.
Énoncez le résultat comme un théorème sur les graphes et démontrez-le dans ce langage.
\item[(b)] Démontrez que si au moins deux personnes sont présentes, alors deux d'entre elles
serrent le même nombre de mains.
\item[(c)] Voici maintenant $ n $ couples mariés qui se rencontrent, l'hôte et l'hôtesse compris,
de sorte que $ 2n $ personnes sont présentes. Personne ne serre la main de son propre
conjoint. L'hôte demande à chacune des $ 2n-1 $ autres personnes combien de mains elle a
serrées, et reçoit $ 2n-1 $ réponses différentes. Déterminez, avec démonstration, combien
de mains l'hôtesse a serrées, en fonction de $ n $, puis déterminez combien l'hôte
lui-même en a serrées.
\item[(d)] Décidez, avec démonstration, si la réponse à la question (c) reste forcée lorsque la
règle sur les conjoints est supprimée, et si elle reste forcée lorsque l'hôte est autorisé à
figurer parmi les personnes qu'il interroge.
\end{enumerate}

\end{problem}

\noindent La question (a) est le lemme des poignées de main, et c'est le plus
ancien théorème de la théorie des graphes : Euler a démontré la formule de la somme des
degrés dans l'article de 1736 sur les ponts de Königsberg qui a fondé la discipline, et tout
argument de parité ultérieur sur les graphes en descend. La question (b) est le principe des
tiroirs appliqué à un ensemble de degrés possibles trop petit d'un élément, et si elle vaut
d'être traitée avec soin, c'est que le principe des tiroirs évident échoue d'exactement un ---
il vous faut remarquer pourquoi deux degrés particuliers ne peuvent pas coexister. La
question (c) circule partout dans les recueils de problèmes sous le nom de \og{} problème de la
soirée \fg{}, sans attribution établie, et elle mérite sa réputation : les données semblent bien
trop maigres pour déterminer quoi que ce soit, la réponse se révèle indépendante de savoir
quel couple est lequel, et la démonstration naturelle est une récurrence sur $ n $ où l'on
épluche les réponses extrêmes deux par deux. La question (d) est là parce qu'un casse-tête
dont vous n'avez pas testé les hypothèses est un casse-tête que vous n'avez pas compris ;
découvrez laquelle des conditions est porteuse.

\begin{refs}
\item Contexte : Lemme des poignées de main --- Wikipédia (\url{https://fr.wikipedia.org/wiki/Lemme_des_poign%C3%A9es_de_main})
\item Contexte : Principe des tiroirs --- Wikipédia (\url{https://fr.wikipedia.org/wiki/Principe_des_tiroirs})
\item Lecture : Miklós Bóna, \textit{A Walk Through Combinatorics}, chapitres sur la récurrence et sur les graphes
\item Voir aussi : Combinatoire et théorie des graphes (\url{pure-combinatorics.html}), L'art de la démonstration (\url{proofs.html})
\end{refs}

\bigskip

\begin{problem}[{Les zéros à la fin d'une factorielle --- Lycée, short}]
\textbf{PUZ-75.} \textbf{Problème.} Déterminez, avec démonstration, combien de zéros
figurent à la fin du développement décimal de $ 100! $, et donnez une formule pour le
nombre de zéros terminaux de $ n! $ valable pour tout entier strictement positif $ n $.
Démontrez ensuite que tout entier positif ou nul n'apparaît pas comme un tel décompte, et
caractérisez exactement quels entiers apparaissent.
\end{problem}

\noindent La formule qu'on vous demande est celle de Legendre : l'exposant d'un
nombre premier $ p $ dans $ n! $ vaut
\[ \nu_p(n!) \;=\; \sum_{i \ge 1} \left\lfloor \frac{n}{p^{i}} \right\rfloor
\;=\; \frac{n - s_p(n)}{p - 1}, \]
où $ s_p(n) $ est la somme des chiffres de $ n $ en base $ p $ ; Adrien-Marie Legendre
l'a consignée dans sa \textit{Théorie des nombres}, et on l'appelle aussi formule de Polignac.
La seconde forme est la plus intéressante, car elle convertit une question de divisibilité
en une question de chiffres, et la dernière partie du problème --- quels décomptes sont sautés
--- se résout entièrement en regardant les développements en base cinq. La même idée de
comptage de chiffres est le théorème de Kummer de 1852, selon lequel l'exposant de $ p $
dans $ \binom{m+n}{n} $ est le nombre de retenues lorsqu'on additionne $ m $ et $ n $
en base $ p $, l'un des faits les plus élégants de la théorie élémentaire des nombres, et
invisible tant qu'on n'écrit pas la factorielle de la seconde façon.

\begin{refs}
\item Contexte : Formule de Legendre --- Wikipédia (\url{https://fr.wikipedia.org/wiki/Formule_de_Legendre})
\item Contexte : Zéro terminal --- Wikipédia (en anglais) (\url{https://en.wikipedia.org/wiki/Trailing_zero})
\item Voir aussi : Théorie des nombres (\url{pure-number-theory.html})
\end{refs}

\bigskip

\begin{problem}[{L'araignée et la mouche --- Lycée}]
\textbf{PUZ-76.} \textbf{Problème.} Une pièce rectangulaire mesure $ 30 $ pieds de long, $ 12 $ pieds de
large et $ 12 $ pieds de haut. Une araignée se tient sur l'un des murs d'extrémité
$ 12 \times 12 $, au milieu de ce mur et à un pied sous le plafond. Une mouche se tient
sur le mur d'extrémité opposé, au milieu et à un pied au-dessus du sol, et ne bouge pas.
L'araignée peut ramper sur les murs, le sol et le plafond, mais ne peut pas se laisser
tomber dans les airs.
\begin{enumerate}
\item[(a)] Déterminez, avec démonstration, la longueur du plus court chemin de l'araignée à la
mouche.
\item[(b)] Démontrez le principe sur lequel repose tout le problème : un plus court chemin sur la
surface d'une boîte est un segment de droite dans tout dépliage plan de la suite de faces
qu'il traverse. Énumérez ensuite toutes les suites de faces qu'un candidat au plus court
chemin pourrait emprunter, dépliez chacune, et justifiez que votre liste est complète ---
c'est là, et non dans l'arithmétique, que sont les mathématiques.
\item[(c)] Une fourmi marche sur la surface d'un cube unité, d'un sommet au sommet diagonalement
opposé. Déterminez la longueur d'un plus court chemin, et déterminez combien de plus courts
chemins distincts il existe.
\item[(d)] Montrez que le dépliage évident n'est pas toujours le bon : trouvez une boîte de côtés
$ a, b, c $ et deux points de sa surface pour lesquels le plus court chemin de surface ne
se trouve \textit{pas} dans le dépliage qu'une première tentative choisirait. Expliquez
ensuite précisément pourquoi une géodésique sur une boîte ne peut jamais passer par un
sommet, et pourquoi ce fait rend le problème combinatoire plutôt qu'affaire de calcul
différentiel.
\end{enumerate}

\end{problem}

\noindent Henry Ernest Dudeney a posé ce problème dans un journal anglais en 1903,
et c'est le plus connu de ses casse-tête géométriques pour la bonne raison que la première
réponse de presque tout le monde est fausse, et fausse par un chemin qui paraît manifestement
optimal. La technique du dépliage est bien plus ancienne que le casse-tête :
l'\textit{Underweysung der Messung} d'Albrecht Dürer, en 1525, dessine les solides platoniciens
et plusieurs solides archimédiens comme des patrons plans à découper et à plier,
apparemment la première fois que quelqu'un mettait des polyèdres sur le papier de cette
manière. Ce que Dudeney ajoute, c'est l'observation qu'un patron n'est pas seulement une aide
à la fabrication mais un changement de coordonnées dans lequel \og{} le plus court \fg{} devient
\og{} tout droit \fg{}. La question (d) ouvre sur un sujet vivant --- savoir si tout polyèdre convexe
possède seulement un patron qui ne se recouvre pas lui-même est un problème non résolu,
repris sur cette page en PUZ-86 (\url{#PUZ-86}) --- et il vaut la peine de savoir que
l'innocente affaire d'aplatir une boîte résiste aux mathématiciens depuis cinq siècles.

\begin{refs}
\item Contexte : Patron (géométrie) --- Wikipédia (\url{https://fr.wikipedia.org/wiki/Patron_%28g%C3%A9om%C3%A9trie%29}), y compris sa section sur les plus courts chemins
\item Contexte : Spider and Fly Problem --- Wolfram MathWorld (en anglais) (\url{https://mathworld.wolfram.com/SpiderandFlyProblem.html}) (attention : donne la réponse)
\item Lecture : Henry E. Dudeney, \textit{Amusements in Mathematics} (1917) --- son propre recueil des casse-tête de journaux
\item Lecture : Erik D. Demaine \& Joseph O'Rourke, \textit{Geometric Folding Algorithms: Linkages, Origami, Polyhedra} (Cambridge, 2007), partie III
\item Voir aussi : PUZ-86 (\url{#PUZ-86}) ci-dessous, et Géométrie (\url{pure-geometry.html})
\end{refs}

\bigskip

\begin{problem}[{Quand les aiguilles d'une horloge se superposent --- Lycée, short}]
\textbf{PUZ-77.} \textbf{Problème.} Sur une horloge analogique ordinaire, les deux
aiguilles tournent régulièrement, la grande aiguille exactement douze fois plus vite que la
petite. Déterminez, avec démonstration, tous les instants d'une période de douze heures où
les deux aiguilles pointent exactement dans la même direction, démontrez que votre liste est
complète, puis décidez --- avec démonstration --- si une trotteuse tournant régulièrement les
rejoint jamais à un instant autre que midi.
\end{problem}

\noindent Les casse-tête sur les aiguilles d'une horloge comptent parmi les plus
anciens de la littérature populaire, et Dudeney comme Sam Loyd en ont imprimé plusieurs.
Dépouillée du cadran, la première moitié est un énoncé sur des angles modulo un tour complet :
les deux aiguilles se superposent exactement lorsque la différence de leurs angles est un
multiple entier de $ 2\pi $, de sorte qu'une question géométrique devient une congruence
linéaire et que les coïncidences tombent à des instants régulièrement espacés dont
l'espacement n'est pas un nombre entier de minutes. Ce dernier fait est ce qui surprend, et
ce qui fait d'un cadran d'horloge une correcte première rencontre avec l'arithmétique modulo
un. La question des trois aiguilles est d'une tout autre nature : elle demande si trois
fonctions affines du temps peuvent partager une solution commune, et une fois les deux
conditions écrites vous constaterez que la réponse tient à une question de rationalité
plutôt qu'à quoi que ce soit d'horloger.

\begin{refs}
\item Contexte : Problème de l'angle des aiguilles --- Wikipédia (en anglais) (\url{https://en.wikipedia.org/wiki/Clock_angle_problem})
\item Contexte : Arithmétique modulaire --- Wikipédia (\url{https://fr.wikipedia.org/wiki/Arithm%C3%A9tique_modulaire})
\item Voir aussi : Théorie des nombres (\url{pure-number-theory.html})
\end{refs}

\bigskip

\begin{problem}[{Le plateau amputé et le triomino en L --- Lycée}]
\textbf{PUZ-78.} \textbf{Problème.} Un \textit{triomino en L} est la forme de trois carrés unités joints en L ;
un \textit{triomino en I} en compte trois alignés. Un plateau est \textit{amputé} lorsqu'on lui a
retiré exactement une case. Les pièces peuvent être tournées, ne peuvent pas se chevaucher,
et doivent recouvrir chaque case restante exactement une fois.
\begin{enumerate}
\item[(a)] Démontrez que, pour tout $ n \ge 1 $ et pour tout choix de la case retirée, le plateau
amputé $ 2^{n} \times 2^{n} $ peut être pavé par des triominos en L. Votre démonstration
doit être une récurrence dont vous puissiez décrire le pas inductif en une phrase.
\item[(b)] Déterminez, avec démonstration, exactement quels rectangles $ m \times n $ --- sans case
retirée --- peuvent être pavés par des triominos en L. Les deux sens sont exigés : une
construction pour les rectangles qui marchent et un argument d'impossibilité pour les
autres.
\item[(c)] Déterminez, avec démonstration, exactement quels plateaux amputés $ m \times n $
peuvent être pavés par des triominos en L. Là où la réponse dépend de la case
\textit{précisément} retirée, dites exactement comment, et démontrez-le.
\item[(d)] Remplacez maintenant le L par le triomino droit en I. Déterminez quels rectangles
$ m \times n $ il pave, et démontrez que le théorème de la question (a) échoue pour lui :
trouvez, pour chaque $ n \ge 1 $, les positions de la case retirée pour lesquelles le
plateau amputé $ 2^{n} \times 2^{n} $ n'admet aucun pavage par des triominos en I. Le
coloriage qui tranche n'est pas le noir et blanc.
\end{enumerate}

\end{problem}

\noindent Solomon Golomb a démontré l'énoncé de la question (a) en 1954, à l'époque
même où il forgeait le mot \og{} polyomino \fg{} ; Martin Gardner l'a porté à un très large lectorat,
et il est depuis la réclame de référence, en classe, pour la récurrence, parce que le pas
inductif contient exactement une idée et que cette idée est visible dès qu'on l'a. Ce qui
rend le reste de ce problème digne d'être posé, c'est que (a) n'est pas représentative. Le
cas $ 2^{n} \times 2^{n} $ est le cas facile ; les rectangles généraux demandent une
construction plus un véritable argument d'impossibilité, et la version amputée de la
question (c) demande les deux ensemble et comporte des exceptions. La question (d) est la
leçon la plus tranchante qui soit sur les arguments de coloriage : deux formes ayant le même
nombre de cases peuvent se comporter tout autrement, et l'invariant qui tue l'une est
invisible au coloriage ordinaire du damier. Comparez PUZ-08 (\url{#PUZ-08}) et
PUZ-22 (\url{#PUZ-22}) sur cette page, où la même difficulté apparaît pour la pièce
$ 1 \times 4 $ et pour le tétromino en T.

\begin{refs}
\item Contexte : Triomino --- Wikipédia (\url{https://fr.wikipedia.org/wiki/Triomino})
\item Contexte : Polyomino --- Wikipédia (\url{https://fr.wikipedia.org/wiki/Polyomino})
\item Lecture : Solomon W. Golomb, \textit{Polyominoes: Puzzles, Patterns, Problems, and Packings} (2e éd., Princeton, 1994)
\item Voir aussi : Combinatoire et théorie des graphes (\url{pure-combinatorics.html}) pour le cas $ 2^{n} \times 2^{n} $ à lui seul, et PUZ-22 (\url{#PUZ-22}) ci-dessus
\end{refs}

\bigskip

\begin{problem}[{Trois interrupteurs, une seule visite --- Lycée, short}]
\textbf{PUZ-87.} \textbf{Problème.} Trois interrupteurs situés au rez-de-chaussée sont
reliés un à un à trois lampes à incandescence identiques dans un grenier sans fenêtre ; vous
pouvez manœuvrer les interrupteurs aussi longtemps que vous le voulez, selon le motif que
vous voulez, mais vous ne pouvez monter au grenier qu'une seule fois et ne jamais y
retourner. Trouvez une procédure qui identifie toujours quel interrupteur commande quelle
lampe, démontrez qu'elle est correcte, puis déterminez le plus grand nombre $ n $ de
lampes qu'une visite unique pourrait résoudre si chaque lampe présente $ m $ états
distinguables à un observateur qui entre.
\end{problem}

\noindent Ce casse-tête relève du folklore --- il circule comme question d'entretien
technique depuis au moins les années 1990 et figure dans la plupart des recueils modernes
sans auteur attribué --- mais la comptabilité qui le sous-tend est le registre même qui régit
PUZ-01 et PUZ-33. Vous n'avez droit qu'à un regard, de sorte que tout ce que vous faites
auparavant doit organiser les choses pour que l'unique observation que vous finirez par faire
prenne une valeur différente pour chacun des $ n! $ câblages possibles ; la seconde moitié
de la question vous demande d'écrire cette condition sous forme d'inégalité et de la
résoudre. Si ce casse-tête est plus qu'un exercice de comptage, c'est qu'il vous force à dire
ce qu'\textit{est} une observation. Un canal porte l'information dont ses sorties peuvent
réellement être distinguées, non l'information que son concepteur voulait lui faire porter,
et cet écart n'est pas une plaisanterie : c'est tout le sujet des attaques par canal
auxiliaire en cryptographie, où l'on a lu des secrets dans la durée, la consommation
électrique, le son et la température de machines qui ne transmettaient rien du tout.

\begin{refs}
\item Contexte : Théorie de l'information --- Wikipédia (\url{https://fr.wikipedia.org/wiki/Th%C3%A9orie_de_l%27information})
\item Contexte : Attaque par canal auxiliaire --- Wikipédia (\url{https://fr.wikipedia.org/wiki/Attaque_par_canal_auxiliaire})
\item Voir aussi : PUZ-01 (\url{#PUZ-01}) et PUZ-33 (\url{#PUZ-33}) pour la même borne de comptage, et Cryptographie et théorie de l'information (\url{applied-crypto.html})
\end{refs}

\bigskip

\begin{problem}[{Les pirates se partagent le butin --- Lycée}]
\textbf{PUZ-88.} \textbf{Problème.} Cinq pirates, hiérarchisés par ancienneté et se connaissant les uns les
autres, doivent partager $ 100 $ pièces d'or indivisibles. Le pirate survivant le plus
ancien propose un partage ; tous les pirates survivants, le proposant compris, votent
ensuite. Si au moins la moitié vote pour, le partage est adopté et la partie s'arrête ; sinon
le proposant est jeté par-dessus bord et le suivant dans l'ordre d'ancienneté fait une
proposition, sous la même règle. Chaque pirate est parfaitement rationnel, sait que tous les
autres le sont, et préfère --- dans cet ordre de priorité --- survivre, puis recevoir plus d'or,
puis, toutes choses égales par ailleurs, voir un compagnon jeté par-dessus bord.
\begin{enumerate}
\item[(a)] Déterminez, avec démonstration, ce que propose le pirate le plus ancien et montrez que
c'est accepté. Construisez l'argument de bas en haut à partir du cas d'un seul pirate.
\item[(b)] Démontrez que votre réponse est l'unique équilibre parfait en sous-jeux du jeu sous forme
extensive à information parfaite correspondant. Montrez ensuite que la préférence de
troisième rang fait un vrai travail : remplacez \og{} préfère voir un pirate jeté par-dessus
bord \fg{} par \og{} préfère ne pas le voir \fg{}, et déterminez comment la réponse change.
\item[(c)] Généralisez à $ n $ pirates et $ g $ pièces. Déterminez, avec démonstration, pour
quels $ n $ exactement le proposant survit, et décrivez le motif de survie une fois que
$ n $ dépasse $ 2g $.
\item[(d)] Changez la règle de vote de sorte qu'une proposition ne passe qu'à la majorité stricte
des survivants, la voix du proposant ne départageant plus les égalités. Refaites (a) et (c)
sous cette règle et dites précisément quelle étape de votre argument précédent échoue.
\end{enumerate}

\end{problem}

\noindent Ian Stewart a présenté ce casse-tête à un large public dans sa chronique
de \textit{Scientific American} de mai 1999, \og{} A puzzle for pirates \fg{}, où il rapportait aussi
l'extension de Steve Omohundro à un nombre quelconque de pirates et de pièces ; une version
était parue l'année précédente dans la contribution de Bruce Talbot Coram à \textit{The Theory of
Institutional Design}. Les mathématiques en jeu sont le raisonnement rétrograde sur un jeu
fini à information parfaite --- la méthode qu'Ernst Zermelo a introduite en 1913 dans le
premier théorème jamais démontré sur les échecs --- et le casse-tête est la démonstration la
plus nette de ce que le raisonnement rétrograde donne vu de l'intérieur : le pouvoir du
pirate le plus ancien sur les autres est entièrement emprunté à un sous-jeu qui ne sera
jamais joué, et le vote de chaque pirate est émis contre une hypothèse. Le résultat est
célèbre pour être à la fois inattaquable et incroyable, ce qui lui vaut de figurer aux côtés
du jeu du mille-pattes dans la littérature sur la question de savoir si les gens raisonnent
réellement à rebours. La question (c) est là où cela cesse d'être une anecdote : la réponse
pour $ n $ général a une véritable structure, et la trouver exige d'organiser la récurrence
plutôt que d'allonger le tableau.

\begin{refs}
\item Contexte : Jeu du pirate --- Wikipédia (\url{https://fr.wikipedia.org/wiki/Jeu_du_pirate})
\item Contexte : Raisonnement rétrograde --- Wikipédia (\url{https://fr.wikipedia.org/wiki/Raisonnement_r%C3%A9trograde})
\item Lecture : Ian Stewart, \og{} A puzzle for pirates \fg{}, \textit{Scientific American}, mai 1999 (attention : la chronique traite intégralement le cas des cinq pirates)
\item Voir aussi : Optimisation et théorie des jeux (\url{applied-optimization.html})
\end{refs}

\bigskip

\begin{problem}[{Les deux tiers de la moyenne --- Lycée}]
\textbf{PUZ-89.} \textbf{Problème.} Chacun de $ n \ge 3 $ joueurs écrit simultanément et indépendamment un
nombre réel de $ [0, 100] $. Le prix va à celui qui est le plus proche des
$ \tfrac{2}{3} $ de la moyenne des $ n $ nombres, les égalités étant partagées à parts
égales. Chaque joueur veut le prix et sait que tous les autres le veulent aussi.
\begin{enumerate}
\item[(a)] Déterminez, avec démonstration, tous les équilibres de Nash en stratégies pures, puis
tous les équilibres de Nash en stratégies mixtes.
\item[(b)] Démontrez que le jeu se résout par élimination itérée des stratégies strictement
dominées. Déterminez exactement quelle part de l'intervalle survit après $ r $ tours
d'élimination, et combien de tours sont nécessaires pour atteindre l'ensemble
d'équilibre.
\item[(c)] Modélisez des raisonneurs imparfaits. Dites qu'un joueur est de \textit{niveau $ 0 $}
s'il tire son nombre uniformément dans $ [0,100] $, et de \textit{niveau $ k $} s'il joue
une meilleure réponse à une population entièrement composée de joueurs de niveau
$ (k-1) $. Calculez explicitement le choix de niveau $ k $ en fonction de $ k $ et de
$ n $, et déterminez ce que la moyenne d'une population observée de choix vous apprend sur
la distribution des niveaux qui la composent.
\item[(d)] Remplacez $ \tfrac{2}{3} $ par une constante $ p > 0 $. Déterminez, avec
démonstration, l'ensemble des équilibres de Nash pour chaque tel $ p $, et décidez pour
quels $ p $ le jeu se résout encore par élimination itérée des stratégies strictement
dominées. Traitez $ p = 1 $ et $ p > 1 $ séparément et expliquez pourquoi ces cas sont
réellement différents.
\end{enumerate}

\end{problem}

\noindent Alain Ledoux a inventé ce jeu en 1981 comme épreuve de départage pour le
magazine français \textit{Jeux et Stratégie}, en demandant à quelque quatre mille lecteurs de
deviner un entier entre un et un milliard ; l'article de Rosemarie Nagel paru en 1995 dans
l'\textit{American Economic Review} en a fait l'instrument de laboratoire standard pour mesurer
ce que les économistes appellent la profondeur de raisonnement, et le quotidien danois
\textit{Politiken} l'a lancé en 2005 avec 19 196 participations. Son ancêtre est le chapitre 12
de la \textit{Théorie générale} de Keynes (1936), où l'investissement professionnel est comparé
à un concours de journal dans lequel il faut désigner non pas les plus jolis visages, mais
ceux que l'on croit que les autres concurrents désigneront --- \og{} et il en est, je crois, qui
pratiquent le quatrième degré, le cinquième et au-delà \fg{}. L'analyse d'équilibre des questions
(a) et (b) est un exercice court. Ce qui vaut à ce jeu sa célébrité, c'est que presque
personne ne joue l'équilibre, et que la distance entre les deux n'est pas du bruit mais une
quantité mesurable : la question (c) vous demande de construire l'instrument qui la
mesure.

\begin{refs}
\item Contexte : p-Concours de beauté (deviner les 2/3 de la moyenne) --- Wikipédia (\url{https://fr.wikipedia.org/wiki/P-Concours_de_beaut%C3%A9})
\item Contexte : Concours de beauté de Keynes --- Wikipédia (\url{https://fr.wikipedia.org/wiki/Concours_de_beaut%C3%A9_de_Keynes})
\item Article : Rosemarie Nagel, \og{} Unraveling in guessing games: an experimental study \fg{}, \textit{American Economic Review} 85 (1995), 1313--1326
\item Voir aussi : Optimisation et théorie des jeux (\url{applied-optimization.html})
\end{refs}

\bigskip

\begin{problem}[{Le jeu de Sim --- Lycée}]
\textbf{PUZ-90.} \textbf{Problème.} Six points sont tracés en position générale, et les $ 15 $ segments
qui les joignent deux à deux sont laissés incolores. Deux joueurs alternent, Rouge
commençant ; un coup colorie un segment incolore dans la couleur du joueur qui joue. Un
joueur qui complète un triangle dont les trois côtés sont de sa propre couleur perd
immédiatement.
\begin{enumerate}
\item[(a)] Démontrez que la partie ne peut jamais être nulle. C'est-à-dire : démontrez que tout
coloriage des arêtes du graphe complet $ K_6 $ en deux couleurs contient un triangle dont
les trois arêtes sont de la même couleur ; puis exhibez un coloriage de $ K_5 $ en deux
couleurs sans un tel triangle, de sorte que le nombre de Ramsey $ R(3,3) $ vaut exactement
$ 6 $.
\item[(b)] Déduisez de la question (a) et du théorème de Zermelo que l'un des deux joueurs possède
une stratégie qui gagne contre toute défense. Déterminez lequel, et décrivez une analyse
exhaustive qui trancherait la question ; le point a été réglé pour la première fois par
ordinateur en 1974, et une stratégie assez simple pour être menée par une personne n'a été
publiée qu'en 2020.
\item[(c)] Jouez la même partie sur cinq points. Démontrez qu'une partie nulle est désormais
possible, et déterminez, avec démonstration, l'issue de la partie à cinq points sous jeu
parfait.
\item[(d)] Généralisez. Énoncez la version jouée sur $ K_n $ dans laquelle un joueur perd en
complétant un $ K_r $ monochrome, déterminez pour quels couples $ (n, r) $ la partie ne
peut pas être nulle, puis expliquez soigneusement pourquoi savoir qu'une partie ne peut pas
être nulle ne vous apprend absolument rien sur celui qui la gagne.
\end{enumerate}

\end{problem}

\noindent Gustavus Simmons --- cryptographe aux Sandia National Laboratories, plus
connu pour avoir fondé la théorie des codes d'authentification --- a publié \og{} The game of SIM \fg{}
dans le \textit{Journal of Recreational Mathematics} en 1969. Ce qu'il avait bâti est la plus
petite conséquence jouable du théorème de Ramsey : la propriété d'absence de nulle de la
question (a) est précisément le fait mondain que, parmi six personnes quelconques, trois se
connaissent mutuellement ou trois sont mutuellement étrangères. Frank Ramsey a démontré le
théorème général dans un article de 1930 sur un problème de décision en logique formelle, où
il n'était qu'un lemme ; il est mort la même année, à vingt-six ans, et la discipline
combinatoire qui porte aujourd'hui son nom est née tout entière de ce lemme. La question (d)
est la morale honnête. La théorie de Ramsey garantit que de la structure doit apparaître, et
le garantit sans la produire --- Erdős avait coutume de dire que si des extraterrestres
exigeaient $ R(5,5) $ nous devrions mobiliser les ordinateurs du monde et l'obtenir, et
que s'ils exigeaient $ R(6,6) $ nous devrions attaquer les extraterrestres --- et le même
divorce entre existence et construction traverse PUZ-56 et PUZ-92 sur cette page.

\begin{refs}
\item Contexte : Sim (jeu de papier-crayon) --- Wikipédia (\url{https://fr.wikipedia.org/wiki/Sim_%28jeu%29})
\item Contexte : Théorème de Ramsey --- Wikipédia (\url{https://fr.wikipedia.org/wiki/Th%C3%A9or%C3%A8me_de_Ramsey})
\item Article : G. J. Simmons, \og{} The game of SIM \fg{}, \textit{Journal of Recreational Mathematics} 2 (1969), no 2
\item Voir aussi : Combinatoire et théorie des graphes (\url{pure-combinatorics.html})
\end{refs}

\bigskip

\begin{problem}[{La majorité qui tourne en rond --- Lycée, short}]
\textbf{PUZ-91.} \textbf{Problème.} Trois électeurs classent trois candidats : le premier
préfère $ A \succ B \succ C $, le deuxième $ B \succ C \succ A $, le troisième
$ C \succ A \succ B $. Vérifiez qu'une majorité stricte préfère $ A $ à $ B $, qu'une
majorité stricte préfère $ B $ à $ C $, et qu'une majorité stricte préfère $ C $ à
$ A $ ; déterminez ensuite, avec démonstration, le plus petit nombre de candidats et le
plus petit nombre d'électeurs pour lesquels un tel cycle de majorités deux à deux peut se
produire.
\end{problem}

\noindent Le marquis de Condorcet a publié ceci dans son \textit{Essai sur l'application
de l'analyse à la probabilité des décisions rendues à la pluralité des voix} en 1785, un
livre écrit pour donner aux élections un fondement mathématique ; il était mathématicien,
abolitionniste de la première heure et partisan du vote des femmes, et il est mort en prison
en 1794, sous la Terreur. Ramon Llull avait trouvé le même cycle au treizième siècle en
concevant des procédures d'élection des dignitaires de l'Église, dans des manuscrits restés
perdus jusqu'au vingt et unième. Le contenu, c'est que \og{} la volonté de la majorité \fg{} peut ne
définir aucun ordre : la préférence majoritaire peut ne pas être transitive, de sorte qu'il
peut n'exister aucun candidat que l'électorat préfère collectivement à tous les autres. Ce
n'est pas non plus un accident rare --- sous le modèle standard où le classement de chaque
électeur est tiré uniformément et indépendamment, la probabilité d'un cycle à trois candidats
tend vers environ neuf pour cent à mesure que l'électorat grandit, et croît avec le nombre de
candidats. Tout ce qui est dans PUZ-98 naît de cet unique exemple.

\begin{refs}
\item Contexte : Paradoxe de Condorcet --- Wikipédia (\url{https://fr.wikipedia.org/wiki/Paradoxe_de_Condorcet})
\item Contexte : Théorie du choix social --- Wikipédia (\url{https://fr.wikipedia.org/wiki/Th%C3%A9orie_du_choix_social})
\item Voir aussi : PUZ-98 (\url{#PUZ-98}) ci-dessous, et Optimisation et théorie des jeux (\url{applied-optimization.html})
\end{refs}

\bigskip


\section*{Licence}

\begin{problem}[{Cent prisonniers et cent boîtes --- Licence}]
\textbf{PUZ-09.} \textbf{Problème.} Cent prisonniers sont numérotés de $ 1 $ à $ 100 $. Dans une pièce
se trouvent cent boîtes contenant, dans un ordre uniformément aléatoire, des papiers
portant les cent numéros. Chaque prisonnier entre seul, peut ouvrir au plus cinquante
boîtes, et doit trouver le papier portant son propre numéro. Les boîtes sont remises dans
leur état initial entre deux visites et aucune communication n'est permise ensuite. Les
cent doivent tous réussir, faute de quoi tous sont exécutés. Les prisonniers peuvent
convenir d'une stratégie à l'avance.
\begin{enumerate}
\item[(a)] Montrez que si chaque prisonnier ouvre cinquante boîtes uniformément au hasard, la
probabilité que tous réussissent vaut $ 2^{-100} $ --- astronomiquement petite.
\item[(b)] Il existe une stratégie qui réussit avec une probabilité supérieure à $ 0{,}30 $.
Trouvez-la. L'indice auquel vous avez droit : la disposition des papiers est une
permutation, et une permutation se décompose en cycles.
\item[(c)] Démontrez que la probabilité de succès de cette stratégie est égale à la probabilité
qu'une permutation uniformément aléatoire de $ 2n $ éléments n'ait aucun cycle de
longueur supérieure à $ n $, et montrez que cela tend vers
$ 1 - \ln 2 \approx 0{,}3069 $ quand $ n \to \infty $.
\end{enumerate}

\end{problem}

\noindent C'est le casse-tête le plus stupéfiant de ce recueil : la réponse naïve
et la réponse correcte diffèrent d'un facteur d'environ $ 10^{29} $, et la première
réaction de presque tout le monde est que la stratégie de la question (b) ne peut pas être
légitime, puisque aucun prisonnier ne communique et que chacun n'ouvre toujours que la
moitié des boîtes. La résolution tient à ce que les succès individuels des prisonniers sont
rendus \textit{massivement corrélés} en indexant la recherche de chacun sur la même
permutation sous-jacente --- ils ne peuvent pas améliorer leurs propres chances, mais ils
peuvent s'arranger pour échouer ensemble plutôt qu'indépendamment. Anna Gál et Peter Bro
Miltersen ont introduit le problème en 2003 dans un article sur les structures de données ;
Peter Winkler l'a popularisé. Que $ 1 - \ln 2 $ apparaisse à la fin, à partir d'un calcul
de somme harmonique sur les longs cycles, est le genre de chose qui fait croire aux
mathématiques.

\begin{refs}
\item Contexte : Problème des 100 prisonniers --- Wikipédia (\url{https://fr.wikipedia.org/wiki/Probl%C3%A8me_des_100_prisonniers})
\item Contexte : Permutations et notation en cycles --- Wikipédia (\url{https://fr.wikipedia.org/wiki/Permutation})
\item Lecture : Peter Winkler, \textit{Mathematical Puzzles: A Connoisseur's Collection}
\item Voir aussi : Probabilités (\url{applied-probability.html}), Algèbre abstraite (\url{pure-algebra.html})
\end{refs}

\bigskip

\begin{problem}[{Des chapeaux en file --- Licence}]
\textbf{PUZ-10.} \textbf{Problème.} Cent prisonniers se tiennent en file, chacun coiffé d'un chapeau rouge
ou bleu. Chacun voit tous les chapeaux devant lui mais aucun de ceux qui sont derrière, ni
le sien. En partant du fond, chacun doit annoncer une couleur à voix haute ; une annonce
fausse est fatale. Ils peuvent convenir d'une stratégie à l'avance, et chaque annonce est
entendue de tous.
\begin{enumerate}
\item[(a)] Trouvez une stratégie garantissant qu'au moins $ 99 $ survivent, et démontrez que la
garantie vaut pour toute attribution possible des chapeaux --- ce n'est pas une affirmation
probabiliste.
\item[(b)] Démontrez qu'aucune stratégie ne peut garantir les $ 100 $ : le dernier prisonnier de
la file n'a aucune information sur son propre chapeau, de sorte qu'un adversaire peut
toujours le mettre en défaut.
\item[(c)] Supposons maintenant que les chapeaux existent en $ k $ couleurs. Adaptez la stratégie
pour qu'au moins $ 99 $ survivent encore, et identifiez la structure algébrique qui fait
le travail.
\end{enumerate}

\end{problem}

\noindent L'idée du bit de parité de la question (a) est un petit miracle de
conception d'information : une annonce unique, choisie pour encoder la parité de tous les
chapeaux qui précèdent, suffit à chaque prisonnier suivant pour déduire exactement sa propre
couleur, puisque chacun peut suivre ce qui a été dit et ce qu'il voit. La question (c)
révèle la structure --- la parité est l'addition dans $ \mathbb{Z}/2 $, et pour $ k $
couleurs on utilise $ \mathbb{Z}/k $, ce qui transforme un casse-tête en somme de
contrôle. C'est le principe même du bit de parité dans une transmission de données, raison
pour laquelle ce casse-tête est un incontournable des cours de théorie des codes.

\begin{refs}
\item Contexte : Casse-tête des chapeaux --- Wikipédia (en anglais) (\url{https://en.wikipedia.org/wiki/Hat_puzzle})
\item Contexte : Bit de parité --- Wikipédia (\url{https://fr.wikipedia.org/wiki/Somme_de_contr%C3%B4le%23Exemple_%3A_bit_de_parit%C3%A9})
\item Voir aussi : Cryptographie et théorie de l'information (\url{applied-crypto.html})
\end{refs}

\bigskip

\begin{problem}[{Les insulaires aux yeux bleus --- Licence}]
\textbf{PUZ-11.} \textbf{Problème.} Sur une île vivent $ 100 $ personnes aux yeux bleus et $ 100 $ aux
yeux bruns. Personne ne connaît la couleur de ses propres yeux, il n'y a ni miroir ni
reflet, et le sujet ne peut jamais être abordé. Quiconque déduit sa propre couleur doit
partir par le bac du soir même. Un jour, un visiteur, s'adressant à tous à la fois, fait
remarquer : \og{} L'un d'entre vous au moins a les yeux bleus. \fg{}
\begin{enumerate}
\item[(a)] Démontrez que, la centième nuit, les cent insulaires aux yeux bleus partent tous.
\item[(b)] Le visiteur a énoncé quelque chose que tous les insulaires savaient déjà par
observation directe. Expliquez précisément quelle information nouvelle a été ajoutée, en
utilisant la notion de \textit{connaissance commune} : distinguez \og{} tout le monde sait
$ P $ \fg{}, \og{} tout le monde sait que tout le monde sait $ P $ \fg{}, et la conjonction infinie
de tous ces énoncés.
\item[(c)] Traitez intégralement les cas $ n = 1, 2, 3 $, et montrez comment la récurrence est
entraînée par le raisonnement de chaque insulaire sur ce que les autres auraient fait les
nuits précédentes.
\end{enumerate}

\end{problem}

\noindent Ce casse-tête est célèbre pour provoquer de véritables désaccords entre
mathématiciens, ce qui est rare. La sensation de paradoxe vient de l'écart entre
connaissance mutuelle et connaissance commune : les insulaires savaient tous qu'il existait
une personne aux yeux bleus, et tous savaient que tous le savaient, et ainsi de suite
jusqu'à n'importe quelle profondeur finie --- mais l'annonce a rendu l'énoncé de connaissance
commune, refermant une régression infinie qu'aucune quantité d'observation privée ne pouvait
refermer. La connaissance commune a été formalisée par David Lewis en 1969 et par Robert
Aumann, dont le théorème d'accord montre que deux agents rationnels ayant des a priori
communs ne peuvent pas \og{} convenir d'être en désaccord \fg{}. Terence Tao a écrit un traitement
soigneux du casse-tête et des objections qu'on lui oppose.

\begin{refs}
\item Contexte : Connaissance commune (logique) --- Wikipédia (\url{https://fr.wikipedia.org/wiki/Connaissance_commune})
\item Contexte : Théorème d'accord d'Aumann --- Wikipédia (en anglais) (\url{https://en.wikipedia.org/wiki/Aumann%27s_agreement_theorem})
\item Lecture : Terence Tao, \og{} The blue-eyed islanders puzzle \fg{} (en anglais) (\url{https://terrytao.wordpress.com/2008/02/05/the-blue-eyed-islanders-puzzle/})
\item Voir aussi : Logique, théorie des ensembles et fondements (\url{pure-foundations.html})
\end{refs}

\bigskip

\begin{problem}[{Les prisonniers et l'ampoule --- Licence, short}]
\textbf{PUZ-12.} \textbf{Problème.} Cent prisonniers sont conduits, un à la fois, dans un
ordre choisi par un adversaire et se répétant indéfiniment, dans une pièce contenant une
unique ampoule qu'ils peuvent allumer ou éteindre. À tout moment, un prisonnier peut
déclarer que les cent sont passés dans la pièce ; si la déclaration est correcte ils sont
libérés, sinon tous sont exécutés. Concevez une stratégie garantissant une déclaration
correcte à terme, et démontrez qu'elle se termine.
\end{problem}

\noindent L'ampoule unique est un bit de mémoire partagée, écrit et lu de façon
asynchrone par cent processus sans horloge --- raison pour laquelle ce casse-tête figure dans
les cours de calcul réparti comme protocole de comptage artisanal. La solution standard
désigne un prisonnier comme compteur ; la question de suivi intéressante est de savoir
combien de visites il faut en moyenne, et si un protocole plus astucieux réduit ce temps
espéré.

\begin{refs}
\item Contexte : Casse-tête des prisonniers et des chapeaux --- Wikipédia (en anglais) (\url{https://en.wikipedia.org/wiki/Prisoners_and_hats_puzzle})
\item Voir aussi : Mathématiques discrètes et informatique (\url{applied-discrete-cs.html})
\end{refs}

\bigskip

\begin{problem}[{Partager un gâteau sans envie --- Licence, short}]
\textbf{PUZ-13.} \textbf{Problème.} Deux personnes partagent un gâteau par \og{} je coupe, tu
choisis \fg{} ; démontrez que ce procédé est sans envie, c'est-à-dire qu'aucune des deux ne
préfère la part de l'autre selon sa propre évaluation. Trouvez ensuite et vérifiez une
procédure sans envie pour trois personnes.
\end{problem}

\noindent Le cas à deux personnes est une démonstration d'une ligne et une leçon de
théorie des mécanismes : celui qui coupe est incité à partager équitablement puisqu'il
reçoit la part qui reste. À trois personnes, c'est étonnamment plus difficile --- la procédure
de Selfridge--Conway a été découverte indépendamment vers 1960 et demande cinq coupes --- et la
procédure bornée pour $ n $ personnes en général est restée ouverte jusqu'en 2016, ce qui
fait longtemps pour un problème de gâteau.

\begin{refs}
\item Contexte : Partage équitable (découpe sans envie) --- Wikipédia (\url{https://fr.wikipedia.org/wiki/Partage_%C3%A9quitable})
\item Voir aussi : Optimisation et théorie des jeux (\url{applied-optimization.html})
\end{refs}

\bigskip

\begin{problem}[{Le chameau et les bananes --- Licence}]
\textbf{PUZ-23.} \textbf{Problème.} Un marchand dispose de $ 3000 $ bananes à un entrepôt et d'un marché situé à $ 1000 $ kilomètres de là, de l'autre côté du désert. Son chameau porte au plus $ 1000 $ bananes à la fois et mange une banane par kilomètre parcouru, dans un sens comme dans l'autre. Les bananes peuvent être déchargées et stockées en n'importe quel point du trajet, et le chameau peut faire autant de voyages que le marchand le souhaite.
\begin{enumerate}
\item[(a)] Trouvez un ordonnancement de voyages et de dépôts qui livre au marché le plus grand nombre possible de bananes.
\item[(b)] Démontrez que votre ordonnancement est optimal : montrez qu'aucun agencement de voyages, de dépôts et de chargements partiels n'en livre davantage.
\item[(c)] Le parent continu de ce casse-tête est le problème de la jeep. Une jeep a un réservoir d'une unité de carburant, dispose de $ n $ unités au total à la base, peut mettre du carburant en cache n'importe où dans le désert, consomme du carburant à taux constant par unité de distance, et doit revenir à la base après chaque voyage sauf le dernier. Déterminez, avec démonstration, la plus grande distance qu'elle peut atteindre lors du voyage final.
\end{enumerate}

\end{problem}

\noindent Le problème 52 des \textit{Propositiones ad Acuendos Juvenes} d'Alcuin, vers l'an 800, met déjà en scène un chameau transportant du grain sur trente lieues et mangeant en chemin ; Luca Pacioli en discute une version dans \textit{De viribus quantitatis} vers 1500. L'analyse moderne est due à N. J. Fine, dont l'article de 1947 dans le \textit{Monthly} a donné la réponse exacte au problème continu et à ses variantes. La même arithmétique régit le ravitaillement en vol et toute expédition qui doit transporter le carburant qu'elle brûle, raison pour laquelle le problème circule aussi sous le nom de problème de l'exploration. Ce qui rend la question (b) difficile, c'est que l'espace des ordonnancements est immense --- vous pouvez fractionner les chargements, placer des dépôts n'importe où et faire un nombre quelconque de navettes --- de sorte qu'une démonstration d'optimalité doit trouver une quantité qu'aucun ordonnancement ne peut améliorer, plutôt que de comparer les ordonnancements un à un.

\begin{refs}
\item Contexte : Problème de la traversée du désert (problème de la jeep) --- Wikipédia (\url{https://fr.wikipedia.org/wiki/Probl%C3%A8me_de_la_travers%C3%A9e_du_d%C3%A9sert})
\item Contexte : Propositiones ad acuendos juvenes --- Wikipédia (\url{https://fr.wikipedia.org/wiki/Propositiones_ad_acuendos_juvenes})
\item Article : N. J. Fine, \og{} The jeep problem \fg{}, \textit{American Mathematical Monthly} 54 (1947), 24--31
\item Voir aussi : Optimisation et théorie des jeux (\url{applied-optimization.html})
\end{refs}

\bigskip

\begin{problem}[{Deux œufs et cent étages --- Licence}]
\textbf{PUZ-24.} \textbf{Problème.} Un immeuble compte $ 100 $ étages, et il existe un étage seuil $ f $ tel qu'un œuf lâché depuis l'étage $ f $ ou au-dessus se casse, tandis qu'un œuf lâché depuis un étage inférieur ne se casse pas. Un œuf qui survit à une chute est intact et peut resservir ; un œuf cassé est perdu. Vous devez déterminer $ f $ exactement, et vous êtes jugé sur le nombre de lâchers dans le pire des cas.
\begin{enumerate}
\item[(a)] Avec un seul œuf, démontrez que $ 100 $ lâchers peuvent être nécessaires et qu'aucune stratégie ne fait mieux dans le pire des cas.
\item[(b)] Avec deux œufs, trouvez une stratégie qui minimise le nombre de lâchers dans le pire des cas, et démontrez qu'aucune stratégie ne fait mieux.
\item[(c)] Avec $ k $ œufs et un budget de $ t $ lâchers, déterminez, avec démonstration, le plus grand nombre d'étages que l'on peut résoudre. Déduisez (b) de votre formule, et écrivez la récurrence qui exprime le problème général comme un programme dynamique.
\end{enumerate}

\end{problem}

\noindent La recherche dichotomique est indisponible ici, et la raison mérite d'être énoncée précisément : un essai qui renvoie \og{} cassé \fg{} vous coûte une ressource, de sorte que les deux branches de l'arbre de décision ne sont pas symétriques, et que l'arbre optimal est donc déséquilibré d'une manière que contrôle le nombre d'œufs restants. C'est le point d'entrée standard, en classe, vers la programmation dynamique et vers la distinction entre optimiser le pire des cas et optimiser le cas moyen, qui donnent ici des stratégies différentes ; Moshe Sniedovich a rédigé le casse-tête avec soin à des fins pédagogiques en 2003, après l'avoir utilisé devant des auditoires d'étudiants à Brunswick et à Hong Kong. Sa provenance relève du folklore --- il a circulé pendant des années comme question d'entretien avant que quiconque prenne la peine de le publier --- mais la réponse générale à $ k $ œufs de la question (c) est assez nette pour valoir la peine d'être devinée sur de petits cas puis démontrée en règle.

\begin{refs}
\item Contexte : Programmation dynamique : le casse-tête des œufs lâchés --- Wikipédia (\url{https://fr.wikipedia.org/wiki/Programmation_dynamique})
\item Contexte : Recherche dichotomique --- Wikipédia (\url{https://fr.wikipedia.org/wiki/Recherche_dichotomique})
\item Article : Moshe Sniedovich, \og{} OR/MS Games: 4. The Joy of Egg-Dropping in Braunschweig and Hong Kong \fg{}, \textit{INFORMS Transactions on Education} 4 (2003)
\item Voir aussi : Mathématiques discrètes et informatique (\url{applied-discrete-cs.html})
\end{refs}

\bigskip

\begin{problem}[{La fourmi sur l'élastique --- Licence, short}]
\textbf{PUZ-25.} \textbf{Problème.} Une fourmi part d'une extrémité d'un élastique de $ 1 $ kilomètre de long et rampe le long de celui-ci à $ 1 $ centimètre par seconde relativement à l'élastique ; à la fin de chaque seconde, l'élastique tout entier est instantanément et uniformément étiré de $ 1 $ kilomètre supplémentaire, entraînant la fourmi avec la matière sur laquelle elle se tient. Déterminez, avec démonstration, si la fourmi atteint jamais l'extrémité opposée et, si oui, estimez le temps que cela lui prend.
\end{problem}

\noindent Martin Gardner a publié ceci dans \textit{Scientific American} et l'a repris dans \textit{aha! Gotcha} en 1982, et cela reste la démonstration la plus nette qu'une intuition de \og{} se laisser distancer sans espoir \fg{} peut être exactement fausse. Le piège est que la distance absolue de la fourmi au point de départ est la mauvaise chose à surveiller : l'étirement éloigne le but, mais il porte aussi la fourmi en avant, et choisir une variable dans laquelle l'étirement se simplifie transforme le problème en une question de convergence de série. La réponse numérique, une fois obtenue, est une bonne leçon sur la différence entre \og{} à terme \fg{} et \og{} en un temps physiquement raisonnable \fg{}.

\begin{refs}
\item Contexte : Problème de la fourmi sur un élastique --- Wikipédia (\url{https://fr.wikipedia.org/wiki/Probl%C3%A8me_de_la_fourmi_sur_un_%C3%A9lastique})
\item Contexte : Série harmonique --- Wikipédia (\url{https://fr.wikipedia.org/wiki/S%C3%A9rie_harmonique})
\item Lecture : Martin Gardner, \textit{aha! Gotcha: Paradoxes to Puzzle and Delight} (1982)
\end{refs}

\bigskip

\begin{problem}[{Les quatre pièces du mercier --- Licence}]
\textbf{PUZ-26.} \textbf{Problème.} En 1902, Henry Dudeney demanda comment découper un triangle équilatéral en pièces qui se réassemblent, sans chevauchement ni trou, en un carré de même aire.
\begin{enumerate}
\item[(a)] Trouvez une telle dissection en quatre pièces. Démontrez qu'elle est exacte : calculez sous forme close les positions des points de coupe et les longueurs des arêtes obtenues, et vérifiez que la figure réassemblée est véritablement un carré et non simplement une figure qui s'en approche.
\item[(b)] Les quatre pièces de Dudeney peuvent être reliées en chaîne par trois charnières, de sorte qu'en faisant pivoter la chaîne dans un sens on referme le triangle et dans l'autre le carré. Démontrez que les points de charnière de votre dissection ont cette propriété, ou ajustez la dissection jusqu'à ce qu'ils l'aient.
\item[(c)] Menez le même programme pour un hexagone régulier : trouvez une dissection d'un hexagone régulier en un carré avec le moins de pièces possible, et démontrez que votre dissection est exacte.
\end{enumerate}

\end{problem}

\noindent Dudeney a posé le problème dans sa rubrique de casse-tête de journal en avril 1902 ; une solution en quatre pièces est arrivée de C. W. McElroy, de Manchester, et Dudeney l'a imprimée quinze jours plus tard, puis de nouveau dans \textit{The Canterbury Puzzles} en 1907, où il décrit un modèle articulé réalisé en acajou et en laiton. L'arrière-plan théorique est le théorème de Wallace--Bolyai--Gerwien : deux polygones quelconques de même aire admettent \textit{une} dissection commune, de sorte que l'existence n'est jamais en question et que toute la difficulté est la minimalité --- sur laquelle, en général, on ne sait presque rien. La question de savoir si les quatre pièces de Dudeney pouvaient être ramenées à trois est restée ouverte cent vingt-deux ans, jusqu'à ce qu'Erik Demaine, Tonan Kamata et Ryuhei Uehara démontrent, dans une prépublication de décembre 2024, qu'il n'existe aucune dissection commune d'un triangle équilatéral et d'un carré en trois pièces ou moins lorsque les pièces ne peuvent pas être retournées. La question (b) n'est pas décorative : Erik Demaine et ses coauteurs ont montré en 2007 que deux polygones quelconques de même aire admettent une dissection commune \textit{articulée}, théorème que le modèle en acajou de Dudeney anticipait d'un siècle.

\begin{refs}
\item Contexte : Puzzle de dissection --- Wikipédia (\url{https://fr.wikipedia.org/wiki/Puzzle_de_dissection})
\item Contexte : Dissection articulée --- Wikipédia (en anglais) (\url{https://en.wikipedia.org/wiki/Hinged_dissection})
\item Article : Erik D. Demaine, Tonan Kamata \& Ryuhei Uehara, \og{} Dudeney's Dissection is Optimal \fg{}, arXiv:2412.03865 (\url{https://arxiv.org/abs/2412.03865}) (2024)
\item Lecture : Greg N. Frederickson, \textit{Dissections: Plane \& Fancy} (1997)
\end{refs}

\bigskip

\begin{problem}[{La cubature du cube --- Licence, short}]
\textbf{PUZ-27.} \textbf{Problème.} Un carré peut être découpé en un nombre fini de carrés plus petits, de tailles deux à deux distinctes : la plus petite dissection de ce type, trouvée par A. J. W. Duijvestijn en 1978, utilise $ 21 $ carrés et a pour côté $ 112 $. Démontrez qu'aucun cube ne peut être découpé en un nombre fini de cubes plus petits de tailles deux à deux distinctes.
\end{problem}

\noindent Le problème bidimensionnel a été résolu à la fin des années 1930 au Trinity College de Cambridge par R. L. Brooks, C. A. B. Smith, A. H. Stone et W. T. Tutte, quatre étudiants qui ont remarqué qu'une dissection d'un rectangle en carrés est le même objet qu'un réseau électrique planaire, les lois de Kirchhoff jouant le rôle de la géométrie ; Roland Sprague a publié le premier carré parfaitement carrelé en 1939. L'énoncé tridimensionnel est faux, et c'est cet échec qui est la partie intéressante : une analogie qui paraît irrésistible dans le plan peut s'effondrer pour des raisons qui tiennent en un paragraphe, et le paragraphe en question est l'un des arguments courts les plus satisfaisants de la géométrie combinatoire.

\begin{refs}
\item Contexte : Quadrature du carré --- Wikipédia (\url{https://fr.wikipedia.org/wiki/Quadrature_du_carr%C3%A9})
\item Article : R. L. Brooks, C. A. B. Smith, A. H. Stone \& W. T. Tutte, \og{} The dissection of rectangles into squares \fg{}, \textit{Duke Mathematical Journal} 7 (1940), 312--340
\item Lecture : Martin Gardner, \og{} Squaring the square \fg{}, dans \textit{The 2nd Scientific American Book of Mathematical Puzzles and Diversions} (1961)
\end{refs}

\bigskip

\begin{problem}[{L'énigme logique la plus difficile du monde --- Licence}]
\textbf{PUZ-35.} \textbf{Problème.} Trois dieux, appelés A, B et C, sont --- dans un ordre que vous ignorez ---
Vrai, Faux et Aléatoire. Vrai répond toujours véridiquement, Faux répond toujours
faussement, et Aléatoire répond comme au gré d'une pièce équilibrée cachée dans sa tête :
chacune de ses réponses est véridique ou fausse au hasard. Vous devez déterminer qui est
qui en posant trois questions à réponse oui-ou-non, chacune adressée à exactement un dieu.
Vous pouvez choisir une question ultérieure à la lumière des réponses précédentes, et vous
pouvez poser plusieurs questions au même dieu. Les dieux comprennent votre langue mais
répondent dans la leur, où les deux mots pour oui et non sont \og{} da \fg{} et \og{} ja \fg{} --- et vous ne
savez pas lequel signifie quoi.
\begin{enumerate}
\item[(a)] Trouvez trois questions qui identifient toujours les trois dieux, et démontrez que votre
procédure fonctionne pour tout arrangement des dieux et tout comportement possible
d'Aléatoire.
\item[(b)] Démontrez le lemme clé sur les questions enchâssées : pour toute question oui-ou-non
$ Q $, Vrai et Faux répondent tous deux \og{} si je vous demandais $ Q $, diriez-vous
ja ? \fg{} par le même mot, et ce mot ne dépend que de la vérité de $ Q $ --- non de savoir
lequel de \og{} da \fg{} et \og{} ja \fg{} signifie oui. Montrez qu'il s'agit de la même simplification que
celle démontrée pour les valets en PUZ-07(c), et identifiez le groupe d'ordre deux qui agit
derrière les deux.
\item[(c)] Boolos accordait trois questions. Démontrez d'abord que trois sont nécessaires si chaque
question doit pouvoir être répondue par le dieu à qui elle est adressée : il y a six
arrangements à distinguer et seulement $ 2^{2} = 4 $ couples de réponses possibles. Lisez
ensuite l'article de 2008 de Rabern et Rabern et expliquez comment le fait d'admettre une
question à laquelle un dieu \textit{ne peut pas} répondre de manière cohérente modifie ce
décompte.
\end{enumerate}

\end{problem}

\noindent George Boolos a publié ceci dans \textit{The Harvard Review of Philosophy}
en 1996 sous le titre qui lui est resté, attribuant l'énigme à Raymond Smullyan et à John
McCarthy la cruelle addition que vous ne savez pas quel mot signifie oui. Boolos a aussi
consigné les règles du jeu --- Aléatoire est réellement une pièce équilibrée, les questions
peuvent être adaptatives, et un dieu répondra à toute question qu'il peut --- car sans elles
l'énigme est mal posée plutôt que difficile. La littérature qui a suivi est d'une qualité
inhabituelle : Roberts en 2001, puis Brian et Landon Rabern en 2008, ont montré qu'un unique
dispositif contrefactuel dompte à la fois le mensonge et la langue, et les Rabern ont ensuite
montré que deux questions suffisent si l'on accepte de poser quelque chose
d'autoréférentiel. La question (b) est le point où l'énigme cesse d'être un tour de passe-passe
pour devenir de l'algèbre : un dieu honnête et un dieu menteur diffèrent par composition avec
la négation, et la négation est une involution.

\begin{refs}
\item Contexte : L'Énigme la plus difficile du monde --- Wikipédia (\url{https://fr.wikipedia.org/wiki/L%27%C3%89nigme_la_plus_difficile_du_monde})
\item Article : George Boolos, \og{} The Hardest Logic Puzzle Ever \fg{}, \textit{The Harvard Review of Philosophy} 6 (1996), 62--65
\item Article : Brian Rabern \& Landon Rabern, \og{} A simple solution to the hardest logic puzzle ever \fg{}, \textit{Analysis} 68 (2008), 105--112
\item Voir aussi : PUZ-07 ci-dessus, et Logique, théorie des ensembles et fondements (\url{pure-foundations.html})
\end{refs}

\bigskip

\begin{problem}[{La pendaison inattendue --- Licence}]
\textbf{PUZ-36.} \textbf{Problème.} Un samedi, un juge dit à un prisonnier : vous serez pendu à midi l'un
des cinq jours de la semaine prochaine, et le matin de la pendaison vous ne saurez pas que
c'est pour ce jour-là. Le prisonnier raisonne que la sentence ne peut pas être exécutée.
Elle ne peut pas tomber le vendredi, car si le jeudi midi était passé il le saurait le
vendredi matin ; elle ne peut donc pas tomber le jeudi non plus, puisque le vendredi est
déjà exclu ; et ainsi de suite en remontant jusqu'au lundi. Il conclut qu'il ne peut pas
être pendu comme décrété. Il est pendu le mercredi, et c'est une surprise complète.
\begin{enumerate}
\item[(a)] Formalisez l'annonce du juge. Notez $ H_d $ pour \og{} la pendaison a lieu à midi le jour
$ d $ \fg{} et $ K_d(\varphi) $ pour \og{} le prisonnier sait le matin du jour $ d $ que
$ \varphi $ \fg{}, et exprimez le décret par une phrase unique. Énoncez précisément quelle
prémisse le raisonnement rétrograde du prisonnier utilise à chaque étape, et lesquelles de
ces prémisses découlent du seul décret.
\item[(b)] Montrez qu'une version à deux jours produit déjà le paradoxe, et analysez-la
complètement. Sous une lecture naturelle, le décret formalisé renvoie à la connaissance
qu'a le prisonnier du décret lui-même ; rendez cette autoréférence explicite et comparez-la
à la phrase du menteur.
\item[(c)] Deux familles d'analyse existent : les logiques, qui cherchent à localiser une véritable
incohérence dans le décret (Shaw, Fitch), et les épistémologiques, qui localisent une
prémisse fausse dans le raisonnement du prisonnier (Quine, Binkley, Chow). Lisez le panorama
de Chow, puis engagez-vous : énoncez une version précise du décret et démontrez soit qu'elle
est satisfiable, soit qu'elle ne l'est pas. Il n'y a pas de consensus, de sorte que toute la
valeur réside ici dans la précision de votre énoncé.
\end{enumerate}

\end{problem}

\noindent Le paradoxe est parvenu à l'imprimé pour la première fois en 1948, dans
la note \og{} Pragmatic paradoxes \fg{} de D. J. O'Connor parue dans \textit{Mind} ; la version que
tout le monde raconte aujourd'hui est celle de Michael Scriven, publiée dans la même revue en
1951, et la chronique de Martin Gardner dans \textit{Scientific American} en mars 1963 l'a
rendue célèbre. Frederic Fitch en a donné un traitement formel en 1964, et le panorama de
Timothy Chow paru en 1998 dans l'\textit{American Mathematical Monthly} reste la meilleure
carte du désaccord, lequel est réel et non résolu : philosophes et logiciens ne s'accordent
pas sur ce qui cloche, seulement sur le fait que quelque chose cloche. Ce paradoxe a sa place
sur une page de casse-tête parce qu'il est la démonstration la plus nette qui soit que \og{} vous
ne saurez pas \fg{} n'est pas une proposition sur le monde mais une proposition sur un sujet
connaissant, et que l'autoréférence à travers un opérateur de connaissance se comporte tout
autrement que l'autoréférence à travers la vérité.

\begin{refs}
\item Contexte : Paradoxe de l'interrogation surprise --- Wikipédia (\url{https://fr.wikipedia.org/wiki/Paradoxe_de_l%27interrogation_surprise})
\item Article : Timothy Y. Chow, \og{} The surprise examination or unexpected hanging paradox \fg{}, \textit{American Mathematical Monthly} 105 (1998), 41--51
\item Lecture : Martin Gardner, \textit{The Unexpected Hanging and Other Mathematical Diversions} (1969)
\item Voir aussi : Logique, théorie des ensembles et fondements (\url{pure-foundations.html})
\end{refs}

\bigskip

\begin{problem}[{Un bit équitable tiré d'une pièce truquée --- Licence, short}]
\textbf{PUZ-37.} \textbf{Problème.} Une pièce tombe sur pile avec probabilité $ p $, où
$ p $ est inconnu et $ 0 < p < 1 $, et les lancers successifs sont indépendants. En
n'utilisant que des lancers de cette pièce, donnez une procédure qui se termine avec
probabilité $ 1 $ et produit $ 0 $ ou $ 1 $ avec probabilité exactement
$ \tfrac{1}{2} $ chacun, quel que soit $ p $, et démontrez qu'elle est exacte.
\end{problem}

\noindent John von Neumann a décrit la réponse dans une note de 1951, \og{} Various
techniques used in connection with random digits \fg{}, écrite alors qu'il se souciait d'obtenir
des chiffres aléatoires utilisables à partir de dispositifs physiques imparfaits pour les
premiers ordinateurs. Ce qui frappe, c'est que $ p $ n'est jamais estimé et n'a même pas
besoin d'être connu approximativement : l'exactitude vient de ce qu'on trouve un événement
dont la probabilité est symétrique par échange des deux faces. Deux questions de suivi
récompensent l'effort. Combien de lancers votre procédure consomme-t-elle en moyenne, et
comment cela se compare-t-il à l'entropie de la source,
$ H(p) = -p\log_2 p - (1-p)\log_2(1-p) $ bits par lancer ? Yuval Peres a montré en 1992
qu'itérer l'idée de von Neumann atteint ce taux.

\begin{refs}
\item Contexte : Pièce équilibrée --- Wikipédia (en anglais) (\url{https://en.wikipedia.org/wiki/Fair_coin}), section sur l'obtention de résultats équitables à partir d'une pièce biaisée
\item Contexte : Extracteur d'aléa --- Wikipédia (en anglais) (\url{https://en.wikipedia.org/wiki/Randomness_extractor})
\item Article : Yuval Peres, \og{} Iterating von Neumann's procedure for extracting random bits \fg{}, \textit{Annals of Statistics} 20 (1992), 590--597
\item Voir aussi : Probabilités (\url{applied-probability.html})
\end{refs}

\bigskip

\begin{problem}[{Les deux généraux --- Licence, short}]
\textbf{PUZ-38.} \textbf{Problème.} Deux généraux campés sur des collines opposées doivent
attaquer à la même heure ou être tous deux anéantis ; ils ne peuvent communiquer que par des
messagers traversant la vallée qui les sépare, et tout messager peut être capturé, auquel
cas l'expéditeur n'apprend jamais si le message est arrivé. Démontrez qu'aucun protocole
échangeant un nombre fini de messages ne peut garantir que les deux généraux attaquent
ensemble.
\end{problem}

\noindent Akkoyunlu, Ekanadham et Huber ont publié le problème et sa démonstration
d'impossibilité en 1975, dans un article sur la conception des communications réseau ; Jim
Gray lui a donné le nom sous lequel on le connaît dans ses \og{} Notes on Data Base Operating
Systems \fg{} de 1978. C'est le premier théorème d'impossibilité du calcul réparti, et c'est la
raison pour laquelle aucun protocole --- y compris tous ceux que votre ordinateur exécute en ce
moment --- ne garantit réellement l'accord sur une liaison non fiable. Ce qui ne peut pas être
atteint est précisément la connaissance commune de PUZ-11 : chaque général a besoin d'un
accusé de réception, puis d'un accusé de réception de l'accusé de réception, et la chaîne ne
se referme jamais. Les systèmes réels se contentent d'un accord avec forte probabilité ; la
suite naturelle est donc de se demander ce qu'un protocole randomisé peut et ne peut pas
promettre.

\begin{refs}
\item Contexte : Problème des deux généraux --- Wikipédia (\url{https://fr.wikipedia.org/wiki/Probl%C3%A8me_des_deux_g%C3%A9n%C3%A9raux})
\item Contexte : Problème du consensus (informatique) --- Wikipédia (\url{https://fr.wikipedia.org/wiki/Probl%C3%A8me_du_consensus})
\item Voir aussi : PUZ-11 ci-dessus, et Mathématiques discrètes et informatique (\url{applied-discrete-cs.html})
\end{refs}

\bigskip

\begin{problem}[{Le paradoxe de Bertrand : trois réponses pour une corde --- Licence}]
\textbf{PUZ-49.} \textbf{Problème.} Un triangle équilatéral est inscrit dans un cercle de rayon $ 1 $. On
trace une corde du cercle \og{} au hasard \fg{} ; quelle est la probabilité qu'elle soit plus longue
qu'un côté du triangle ?
\begin{enumerate}
\item[(a)] Calculez la réponse sous chacune de trois constructions, et démontrez chaque calcul :
(i) les deux extrémités sont indépendantes et uniformes sur la circonférence ; (ii) une
direction est choisie uniformément, puis la distance de la corde au centre est uniforme sur
$ [0,1] $ ; (iii) le milieu de la corde est uniforme sur le disque. Montrez que les trois
réponses sont
\[ \tfrac{1}{3}, \qquad \tfrac{1}{2}, \qquad \tfrac{1}{4} . \]
\item[(b)] Calculez, pour chacune des trois constructions, la loi de la distance $ D $ de la
corde au centre. Vérifiez que les trois lois sont véritablement différentes, et redérivez
les trois réponses de la question (a) à partir de la seule condition $ D < 1/2 $ --- de
sorte que le désaccord soit localisé en un point et ne soit pas un artefact de calcul
bâclé.
\item[(c)] Montrez que, des trois, seule (ii) produit une loi invariante par translation et
homothétie du cercle dans le plan. C'est la raison qu'invoque E. T. Jaynes pour la préférer ;
énoncez son critère précisément et démontrez l'affirmation d'unicité.
\item[(d)] Démontrez que (ii) est aussi ce que l'on obtient à partir de l'unique mesure sur les
droites du plan invariante par tous les déplacements, restreinte aux droites rencontrant le
disque --- la mesure cinématique de la géométrie intégrale. Argumentez ensuite soigneusement
pourquoi cela ne montre \textit{pas} que (i) et (iii) sont fausses.
\end{enumerate}

\end{problem}

\noindent Joseph Bertrand a exposé ceci dans son \textit{Calcul des probabilités} de
1889 pour faire valoir un point qu'il jugeait évident et que la plupart des lecteurs trouvent
scandaleux : \og{} au hasard \fg{} n'est pas une spécification. Tant que vous n'avez pas dit quel
mécanisme produit la corde --- ou, ce qui revient au même, quelle mesure sur l'espace des
cordes vous avez en tête --- la question n'a pas de réponse, et les trois constructions
ci-dessus sont trois questions différentes sous les mêmes habits. L'article de Jaynes de
1973, \og{} The Well-Posed Problem \fg{}, soutenait qu'une version physique bien posée singularise
bel et bien (ii), parce qu'une expérience réelle consistant à jeter des pailles sur le cercle
ne peut savoir ni où est le cercle ni quelle est sa taille ; Alon Drory a depuis soutenu que
le principe même de Jaynes peut être orienté vers n'importe laquelle des trois, et la dispute
philosophique n'est pas close. Les mathématiques, en revanche, sont entièrement closes, et
la morale --- spécifiez la mesure --- est l'habitude la plus utile qui soit en probabilités
appliquées.

\begin{refs}
\item Contexte : Paradoxe de Bertrand (probabilités) --- Wikipédia (\url{https://fr.wikipedia.org/wiki/Paradoxe_de_Bertrand})
\item Contexte : Principe d'indifférence --- Wikipédia (en anglais) (\url{https://en.wikipedia.org/wiki/Principle_of_indifference})
\item Article : E. T. Jaynes, \og{} The well-posed problem \fg{}, \textit{Foundations of Physics} 3 (1973)
\item Lecture : Daniel A. Klain \& Gian-Carlo Rota, \textit{Introduction to Geometric Probability}
\end{refs}

\bigskip

\begin{problem}[{Le jeu de Wythoff et le nombre d'or --- Licence}]
\textbf{PUZ-50.} \textbf{Problème.} Le jeu de Wythoff se joue avec deux tas de jetons. Un coup retire soit un
nombre strictement positif quelconque de jetons d'un seul tas, soit le \textit{même} nombre
strictement positif dans les deux tas. Le joueur qui prend le dernier jeton gagne.
\begin{enumerate}
\item[(a)] À la main, trouvez toutes les positions perdantes $ (a, b) $ avec
$ a \le b \le 20 $ --- une position est perdante lorsque le joueur qui doit jouer perd sous
jeu parfait --- et énumérez-les par ordre croissant sous la forme
$ (a_1, b_1), (a_2, b_2), \dots $.
\item[(b)] Démontrez que ces positions sont exactement données par
\[ a_k = \lfloor k\varphi \rfloor, \qquad b_k = \lfloor k\varphi^2 \rfloor,
 \qquad \varphi = \frac{1+\sqrt{5}}{2}, \quad k \ge 1, \]
et que $ b_k = a_k + k $.
\item[(c)] Démontrez le théorème de Rayleigh (dit aussi de Beatty) : si $ \alpha $ et
$ \beta $ sont irrationnels, supérieurs à $ 1 $, et vérifient
$ \frac{1}{\alpha} + \frac{1}{\beta} = 1 $, alors les suites
$ \lfloor k\alpha \rfloor $ et $ \lfloor k\beta \rfloor $, $ k \ge 1 $, contiennent
ensemble chaque entier strictement positif exactement une fois.
\item[(d)] Expliquez pourquoi la question (c) est exactement le fait combinatoire qu'un ensemble de
positions perdantes de ce jeu doit vérifier, et vérifiez que $ \varphi $ et
$ \varphi^2 $ en satisfont les hypothèses.
\end{enumerate}

\end{problem}

\noindent W. A. Wythoff a publié son analyse dans \textit{Nieuw Archief voor Wiskunde}
en 1907 ; il ajoutait une règle au Nim et découvrait, à l'étonnement général, que la réponse
est le nombre d'or. Rien dans l'énoncé ne mentionne $ \sqrt{5} $ ; il apparaît parce que
les positions perdantes doivent former deux suites complémentaires avec un motif de
différences fixe, et le théorème de partition de Rayleigh --- que Lord Rayleigh avait signalé
dans la deuxième édition de \textit{The Theory of Sound} en 1894 et que Samuel Beatty a
redécouvert comme problème du \textit{Monthly} en 1926 --- dit que le seul couple d'irrationnels
capable de faire le travail est $ \varphi, \varphi^2 $. C'est l'un des plus purs exemples
en mathématiques d'une constante continue extorquée à une question discrète. Martin Gardner
a rapporté que le jeu se pratiquait en Chine sous le nom de \textit{jiǎn shízǐ}, \og{} ramasser des
pierres \fg{}, même si, comme pour le Nim, l'histoire est incertaine.

\begin{refs}
\item Contexte : Jeu de Wythoff --- Wikipédia (\url{https://fr.wikipedia.org/wiki/Jeu_de_Wythoff})
\item Contexte : Suite de Beatty --- Wikipédia (\url{https://fr.wikipedia.org/wiki/Th%C3%A9or%C3%A8me_de_Beatty})
\item Contexte : Nombre d'or --- Wikipédia (\url{https://fr.wikipedia.org/wiki/Nombre_d%27or})
\item Lecture : Elwyn Berlekamp, John Conway \& Richard Guy, \textit{Winning Ways for Your Mathematical Plays}, vol. 1
\end{refs}

\bigskip

\begin{problem}[{Le problème du scrutin --- Licence, short}]
\textbf{PUZ-51.} \textbf{Problème.} Lors d'une élection, le candidat A reçoit $ p $ voix
et le candidat B en reçoit $ q $, avec $ p > q $, et les $ p + q $ bulletins sont
dépouillés dans un ordre uniformément aléatoire. Démontrez que la probabilité que A soit
strictement en tête de B pendant tout le dépouillement vaut
\[ \frac{p-q}{p+q} . \]
\end{problem}

\noindent W. A. Whitworth a publié le résultat en 1878 et Joseph Bertrand l'a
redécouvert en 1887, en remarquant qu'une réponse aussi simple devait avoir une démonstration
directe ; Désiré André en a fourni une dans l'année. Tous les manuels appellent aujourd'hui
cela la méthode de réflexion --- et Marc Renault a montré en 2008, en revenant au texte
français, qu'André n'avait employé aucune réflexion, mais une bijection explicite entre deux
familles d'ordres de dépouillement défavorables. Quel que soit le nom qu'on lui donne, le
dispositif est l'ancêtre du principe de réflexion pour les marches aléatoires et pour le
mouvement brownien, ainsi que de la démonstration standard que les nombres de Catalan
comptent les chemins de réseau restant faiblement d'un même côté de la diagonale.

\begin{refs}
\item Contexte : Problème du scrutin --- Wikipédia (\url{https://fr.wikipedia.org/wiki/Probl%C3%A8me_du_scrutin})
\item Contexte : Nombre de Catalan --- Wikipédia (\url{https://fr.wikipedia.org/wiki/Nombre_de_Catalan})
\item Article : Marc Renault, \og{} Lost (and found) in translation: André's actual method and its application to the generalized ballot problem \fg{}, \textit{American Mathematical Monthly} 115 (2008), 358--363
\item Voir aussi : Combinatoire et théorie des graphes (\url{pure-combinatorics.html}) pour les nombres de Catalan, Probabilités (\url{applied-probability.html}) pour les marches aléatoires
\item Voir aussi : PRB-36 (\url{applied-probability.html#PRB-36}) --- la page Probabilités, où il est développé avec le principe de réflexion.
\end{refs}

\bigskip

\begin{problem}[{Tout jeu impartial est un tas de Nim --- Licence}]
\textbf{PUZ-52.} \textbf{Problème.} Un jeu \textit{impartial} est un jeu à deux joueurs sans hasard ni
information cachée dans lequel, depuis toute position, les deux joueurs disposent du même
ensemble de coups licites, aucune partie infinie n'est possible, et le joueur qui ne peut
plus jouer perd. Pour une position $ G $ d'options $ G_1, \dots, G_m $, définissez la
valeur de Grundy
\[ \mathcal{G}(G) = \operatorname{mex}\{\mathcal{G}(G_1), \dots, \mathcal{G}(G_m)\}, \]
où $ \operatorname{mex} S $ est le plus petit entier positif ou nul n'appartenant pas à
$ S $.
\begin{enumerate}
\item[(a)] Démontrez que $ \mathcal{G}(G) = 0 $ si et seulement si le joueur qui doit jouer
depuis $ G $ perd sous jeu parfait, et qu'un unique tas de Nim de $ n $ jetons a pour
valeur de Grundy $ n $.
\item[(b)] Démontrez le théorème de Sprague--Grundy : pour la somme disjonctive $ G + H $, dans
laquelle un coup consiste à jouer dans exactement l'une des deux composantes,
\[ \mathcal{G}(G + H) = \mathcal{G}(G) \oplus \mathcal{G}(H), \]
la somme de Nim de PUZ-48. Concluez que toute position impartiale est équivalente, à toutes
fins de jeu, à un unique tas de Nim.
\item[(c)] Appliquez-le. Tabulez les valeurs de Grundy du jeu de soustraction dans lequel un coup
retire $ 1 $, $ 2 $ ou $ 3 $ jetons d'un tas, et celles du jeu de Grundy, dans lequel
un coup remplace un tas par deux tas de tailles \textit{différentes} ; donnez les vingt
premières valeurs de chacun. Servez-vous ensuite de (b) pour décider, avec démonstration,
qui gagne une position composée d'un tas de chaque jeu joués simultanément.
\item[(d)] Expliquez pourquoi le théorème échoue sous la convention misère, où le joueur qui ne
peut plus jouer gagne : montrez par un exemple que la valeur de Grundy d'une composante ne
détermine plus sa contribution à une somme.
\end{enumerate}

\end{problem}

\noindent R. P. Sprague (dans le \textit{Tohoku Mathematical Journal}, 1936) et
P. M. Grundy (dans \textit{Eureka}, 1939) ont démontré cela indépendamment, une trentaine
d'années après que Bouton eut résolu le Nim, et c'est le théorème qui a transformé un jeu
résolu en une théorie. L'affirmation est saisissante à la manière d'un théorème de
classification : il existe une infinité de jeux impartiaux, ils ne se ressemblent en rien,
et chacun d'eux est un tas de Nim. La question (d) marque la frontière du pouvoir de la
théorie --- le jeu misère est resté en désordre pendant soixante-dix ans, jusqu'à ce que les
quotients misère de Plambeck et Siegel lui donnent une structure dans les années 2000, et il
demeure bien plus difficile que le jeu normal. Le jeu de Grundy de la question (c) rappelle
combien peu de choses sont comprises dans le détail : savoir si sa suite de valeurs de Grundy
devient jamais périodique est un problème non résolu, conjecturé vrai par Berlekamp, Conway
et Guy, et vérifié par Achim Flammenkamp jusqu'à une longueur énorme sans être tranché.

\begin{refs}
\item Contexte : Théorème de Sprague-Grundy --- Wikipédia (\url{https://fr.wikipedia.org/wiki/Th%C3%A9or%C3%A8me_de_Sprague-Grundy})
\item Contexte : Jeu de Grundy --- Wikipédia (\url{https://fr.wikipedia.org/wiki/Jeu_de_Grundy})
\item Contexte : Mex (mathématiques) --- Wikipédia (en anglais) (\url{https://en.wikipedia.org/wiki/Mex_(mathematics)})
\item Lecture : Aaron N. Siegel, \textit{Combinatorial Game Theory}, AMS Graduate Studies in Mathematics 146 (2013)
\end{refs}

\bigskip

\begin{problem}[{Pourquoi le ciel est bleu, et pourquoi le couchant ne l'est pas --- Licence}]
\textbf{PUZ-64.} \textbf{Problème.} Une lumière de longueur d'onde $ \lambda $ traverse un gaz dilué de
$ n $ molécules par unité de volume. Chaque molécule est bien plus petite que
$ \lambda $, n'absorbe pas, et répond au champ incident comme un dipôle induit
$ \vec{p} = \alpha \vec{E} $, où, en unités de Gauss, $ \alpha $ a les dimensions d'un
volume.
\begin{enumerate}
\item[(a)] Calculez la puissance moyennée dans le temps rayonnée par un dipôle oscillant, divisez
par l'intensité incidente, et établissez la section efficace de Rayleigh
\[ \sigma(\lambda) \;=\; \frac{128\,\pi^{5}\,\alpha^{2}}{3\,\lambda^{4}} . \]
Énumérez chaque hypothèse utilisée par la dérivation, et dites où chacune échouerait.
\item[(b)] Obtenez l'exposant $ -4 $ une seconde fois, par la seule analyse dimensionnelle, en
n'utilisant que $ \alpha $ et $ \lambda $. Démontrez qu'aucune autre puissance n'est
disponible pour un diffuseur non absorbant bien plus petit que la longueur d'onde, et
expliquez ce qui se dérègle lorsque le diffuseur est comparable à $ \lambda $.
\item[(c)] À l'aide de $ \sigma $ et de la densité numérique de l'air, calculez l'épaisseur
optique de l'atmosphère à $ \lambda = 450\,\mathrm{nm} $ et à
$ \lambda = 650\,\mathrm{nm} $, d'abord le long d'un trajet vertical puis le long d'un
trajet à $ 89^\circ $ du zénith. À partir des deux spectres transmis, rendez
quantitativement compte de la couleur du ciel de jour et de celle d'un soleil bas.
\item[(d)] Le violet diffuse plus fortement que le bleu, et pourtant le ciel est bleu. Rendez
quantitativement compte de l'écart, en utilisant l'éclairement spectral solaire au sommet de
l'atmosphère ainsi que les courbes de réponse des trois types de cônes humains, et montrez
que les deux contributions sont d'ampleur comparable.
\end{enumerate}

\end{problem}

\noindent John Tyndall a montré en 1869 que la lumière diffusée latéralement hors
d'un faisceau par une fine suspension est bleue et fortement polarisée, ce qui vaut à l'effet
de porter son nom en chimie des colloïdes. Lord Rayleigh a établi la loi en
$ \lambda^{-4} $ en 1871 et a d'abord attribué la couleur du ciel à la poussière ; en 1899
il s'est corrigé et a soutenu que les molécules de l'air suffisent à elles seules --- calcul
qui, mené à rebours à partir de la luminosité mesurée du ciel, donne le nombre d'Avogadro,
de sorte que compter les atomes et admirer le temps qu'il fait se révèlent être la même
activité. Einstein a refondu tout le sujet en 1910 en termes de fluctuations thermodynamiques
de densité, ce qui explique que le ciel et l'opalescence critique d'un fluide près de son
point critique soient un seul et même phénomène. La question (d) est celle que les manuels
sautent, et aucun argument d'échelle ne peut la trancher : elle exige de vrais nombres tirés
d'un spectre solaire et d'un jeu de courbes de réponse des cônes, et la réponse honnête
nomme deux effets plutôt qu'un.

\begin{refs}
\item Contexte : Diffusion Rayleigh --- Wikipédia (\url{https://fr.wikipedia.org/wiki/Diffusion_Rayleigh})
\item Contexte : Couleur du ciel --- Wikipédia (\url{https://fr.wikipedia.org/wiki/Couleur_du_ciel})
\item Article : Lord Rayleigh, \og{} On the light from the sky, its polarisation and colour \fg{}, \textit{Philosophical Magazine} 41 (1871) ; et \og{} On the transmission of light through an atmosphere containing small particles in suspension, and on the origin of the blue of the sky \fg{}, \textit{Philosophical Magazine} 47 (1899)
\item Lecture : J. D. Jackson, \textit{Classical Electrodynamics}, 3e éd., chap. 10 (diffusion et diffraction)
\end{refs}

\bigskip

\begin{problem}[{Les feuilles de thé d'Einstein --- Licence}]
\textbf{PUZ-65.} \textbf{Problème.} Des feuilles de thé reposent au fond plat d'une tasse cylindrique de
rayon $ R $, remplie jusqu'à une profondeur $ H $ d'eau de viscosité cinématique
$ \nu $. L'eau est remuée jusqu'à tourner presque exactement comme un solide, à la vitesse
angulaire $ \Omega $, puis la cuillère est retirée. Les feuilles migrent alors le long du
fond.
\begin{enumerate}
\item[(a)] Calculez le champ de pression d'un fluide incompressible en rotation solide dans un
cylindre, y compris la forme de la surface libre, et vérifiez qu'à toute profondeur le
gradient de pression radial est exactement celui qu'il faut pour maintenir le fluide sur sa
trajectoire circulaire.
\item[(b)] Au fond, la condition d'adhérence annule la vitesse azimutale sur une couche mince.
Montrez que le gradient de pression radial calculé en (a) est imposé à cette couche
essentiellement inchangé depuis l'intérieur, et écrivez les équations de couche limite qui
gouvernent l'écoulement radial qui en résulte.
\item[(c)] Justifiez l'estimation $ \delta \sim \sqrt{\nu/\Omega} $ de l'épaisseur de la couche,
calculez le débit volumique qu'elle transporte, et déduisez-en l'échelle de temps sur
laquelle tout le contenu de la tasse est retourné. Comparez cette échelle de temps au temps
visqueux de ralentissement $ R^{2}/\nu $ et à ce que vous observez dans une vraie tasse de
thé.
\item[(d)] Déterminez où les feuilles viennent se poser, démontrez que la conclusion ne dépend pas
du sens dans lequel on remue, puis dites ce qui change si la tasse est conique plutôt que
cylindrique, ou si les feuilles sont assez lourdes pour que leur vitesse de sédimentation
soit comparable à l'écoulement trouvé en (c).
\end{enumerate}

\end{problem}

\noindent James Thomson --- le frère aîné de Kelvin --- a donné l'explication
essentielle en 1857, à propos des méandres de rivières plutôt que de tasses de thé, et
Boussinesq a traité théoriquement l'écoulement secondaire en 1868. Einstein y est revenu dans
une note de deux pages parue dans \textit{Die Naturwissenschaften} en mars 1926, intitulée
\og{} Die Ursache der Mäanderbildung der Flussläufe und des sogenannten Baerschen Gesetzes \fg{} --- la
cause de la formation des méandres des cours d'eau et de la loi dite de Baer. La tasse de thé
est le modèle jouet ; le vrai sujet est de savoir pourquoi une rivière creuse sa rive
extérieure et divague dans sa plaine d'inondation, et pourquoi, comme l'avait affirmé Karl
Ernst von Baer, les rives droites des fleuves du Nord sont les plus escarpées. La même
circulation, entraînée par le désaccord entre un gradient de pression intérieur et une paroi
à adhérence, est la couche d'Ekman de l'océanographie et de la météorologie, découverte par
Vagn Walfrid Ekman en 1905 alors qu'il expliquait pourquoi la glace dérivante ne se déplace
pas dans le sens du vent. Les brasseurs s'en servent délibérément : un \og{} tourbillon \fg{} obtenu
en remuant rassemble le houblon usé en un cône bien net au centre de la cuve, et les
centrifugeuses de laboratoire exploitent le même écoulement secondaire pour séparer le
sang.

\begin{refs}
\item Contexte : Paradoxe des feuilles de thé --- Wikipédia (\url{https://fr.wikipedia.org/wiki/Paradoxe_des_feuilles_de_th%C3%A9})
\item Contexte : Couche d'Ekman --- Wikipédia (\url{https://fr.wikipedia.org/wiki/Transport_d%27Ekman})
\item Article : A. Einstein, \og{} Die Ursache der Mäanderbildung der Flussläufe und des sogenannten Baerschen Gesetzes \fg{}, \textit{Die Naturwissenschaften} 14 (1926), 223--224
\item Lecture : H. P. Greenspan, \textit{The Theory of Rotating Fluids} (1968) ; ou G. K. Batchelor, \textit{An Introduction to Fluid Dynamics}
\end{refs}

\bigskip

\begin{problem}[{Dans quel sens la baignoire se vide --- Licence, short}]
\textbf{PUZ-66.} \textbf{Problème.} De l'eau s'écoule par un trou de rayon $ a $ percé
dans le fond plat d'un bassin circulaire de rayon $ R $ et de profondeur $ H $, à la
latitude $ \varphi $ ; notez $ f = 2\Omega_\oplus \sin\varphi $ le paramètre de Coriolis
et $ \mathrm{Ro} = U/(fL) $ le nombre de Rossby. Déterminez la plus grande circulation
résiduelle que le bassin peut conserver si la rotation de la Terre doit contrôler le sens du
tourbillon, énoncez le critère obtenu sous forme d'inégalité explicite en $ a $, $ R $,
$ H $, $ \varphi $ et le temps de repos, et comparez-le à ce que le bassin de Shapiro au
MIT en 1962 et la reprise de Trefethen à Sydney en 1965 ont réellement dû organiser.
\end{problem}

\noindent Ascher Shapiro a publié une note d'une page dans \textit{Nature} en 1962 :
un bassin de deux mètres de diamètre, rempli par une arrivée radiale, couvert de plastique et
laissé au repos une journée entière, se vidait dans le sens antihoraire à Cambridge,
Massachusetts, à chaque fois. Trois ans plus tard, Lloyd Trefethen et quatre collègues de
l'université de Sydney ont fait la version australe et ont rapporté le sens horaire. Les deux
résultats sont réels, et aucun des deux ne tranche l'affirmation populaire sur les baignoires
et les toilettes, car cette affirmation porte sur des \textit{ordres de grandeur} tandis que
les expériences décrivent ce qui subsiste une fois retiré, à grand-peine, tout ordre de
grandeur concurrent --- une journée de repos, un bassin couvert, une arrivée symétrique.
Décider si une baignoire ordinaire est seulement proche de ce régime est précisément le
calcul demandé ci-dessus, et il vaut la peine d'être fait, car l'habitude qu'il enseigne est
la plus transférable de la physique : un effet peut être à la fois réel et sans importance,
et seul un rapport sans dimension vous dit lequel. Le même nombre de Rossby décide si un
système météorologique sent la planète tourner ; pour un cyclone il est petit, ce qui
explique que les ouragans, eux, obéissent bien à l'hémisphère.

\begin{refs}
\item Contexte : Force de Coriolis --- Wikipédia (\url{https://fr.wikipedia.org/wiki/Force_de_Coriolis})
\item Contexte : Nombre de Rossby --- Wikipédia (\url{https://fr.wikipedia.org/wiki/Nombre_de_Rossby})
\item Article : A. H. Shapiro, \og{} Bath-tub vortex \fg{}, \textit{Nature} 196 (1962), 1080--1081
\item Article : L. M. Trefethen, R. W. Bilger, P. T. Fink, R. E. Luxton \& R. I. Tanner, \og{} The bath-tub vortex in the southern hemisphere \fg{}, \textit{Nature} 207 (1965), 1084--1085
\end{refs}

\bigskip

\begin{problem}[{L'échelle, la grange et les jumeaux --- Licence}]
\textbf{PUZ-67.} \textbf{Problème.} Une échelle de longueur propre $ L $ se déplace selon sa propre
direction à la vitesse constante $ v $ à travers une grange de longueur propre
$ \ell < L $ munie d'une porte avant et d'une porte arrière. Dans le référentiel de la
grange, les deux portes sont fermées simultanément à l'instant où l'échelle est entièrement
à l'intérieur, et rouvertes un instant plus tard.
\begin{enumerate}
\item[(a)] Calculez, dans le référentiel de la grange et dans celui de l'échelle, l'instant et le
lieu où chaque porte se ferme. Mettez explicitement en évidence la dépendance au référentiel
du mot \og{} simultanément \fg{}, et déterminez la condition sur $ v $ sous laquelle la description
du référentiel de la grange est seulement possible.
\item[(b)] Dans le référentiel de l'échelle, la grange est plus courte qu'au repos, de sorte que
l'échelle dépasse la grange de plus qu'avant. Montrez que les deux référentiels ont
néanmoins raison, en calculant l'intervalle invariant entre les deux événements de fermeture
des portes et en le classant.
\item[(c)] Supposez que la porte arrière se coince et ne s'ouvre pas. Démontrez que l'échelle ne
peut pas être traitée comme rigide, établissez une borne supérieure sur la vitesse à
laquelle la nouvelle du choc peut se propager le long d'elle, et déduisez-en que \og{} solide
rigide \fg{} n'est pas un concept relativiste.
\item[(d)] Un jumeau reste sur Terre ; l'autre voyage vers une étoile située à la distance $ D $
dans le référentiel terrestre à la vitesse $ v $, fait instantanément demi-tour, et
revient. Calculez les deux temps propres. Localisez ensuite, avec démonstration, l'erreur
exacte dans l'argument naïf de symétrie selon lequel chaque jumeau devrait trouver l'autre
plus jeune, en suivant ce qu'il advient de la surface de simultanéité du voyageur au moment
du demi-tour.
\end{enumerate}

\end{problem}

\noindent La contraction des longueurs est entrée deux fois en physique : comme
rétrécissement physique dans l'éther, proposé par FitzGerald en 1889 et par Lorentz en 1892
pour sauver le résultat nul de Michelson--Morley, puis en 1905 comme énoncé d'Einstein sur ce
que signifie mesurer. La version de la grange et de l'échelle est un dispositif pédagogique
des années 1960, développé longuement dans les manuels de Wolfgang Rindler, et tout son
intérêt est que \og{} longueur \fg{} abrège \og{} les positions des deux extrémités au même instant \fg{} ---
formule qui cesse discrètement d'être indépendante du référentiel. Le jumeau voyageur a été
publié par Paul Langevin en 1911, sous la forme d'une histoire de voyageur expédié dans un
boulet de canon, bien qu'Einstein eût noté l'effet d'horloge dès 1905. Il vaut la peine de
savoir que la résolution n'est pas évidente, même pour des professionnels : Herbert Dingle a
passé la fin des années 1950 et les années 1960 à soutenir dans les pages de \textit{Nature}
que le résultat des jumeaux était contradictoire, et l'échange est devenu acrimonieux avant
de s'achever. Hafele et Keating ont réglé la question empirique en 1971 en faisant voler des
horloges au césium autour du monde dans les deux sens et en les comparant aux horloges de
l'observatoire naval des États-Unis.

\begin{refs}
\item Contexte : Paradoxe de l'échelle (paradoxe du train) --- Wikipédia (\url{https://fr.wikipedia.org/wiki/Paradoxe_du_train})
\item Contexte : Paradoxe des jumeaux --- Wikipédia (\url{https://fr.wikipedia.org/wiki/Paradoxe_des_jumeaux})
\item Cours : MIT OCW 8.20 --- Introduction to Special Relativity (\url{https://ocw.mit.edu/courses/8-20-introduction-to-special-relativity-january-iap-2021/})
\item Lecture : E. F. Taylor \& J. A. Wheeler, \textit{Spacetime Physics}, 2e éd. --- voir aussi Physique mathématique (\url{applied-physics.html}) pour l'établissement de la transformation de Lorentz
\end{refs}

\bigskip

\begin{problem}[{Le pari des anniversaires, et comment le gagner honnêtement --- Licence}]
\textbf{PUZ-79.} \textbf{Problème.} Prenez $ n $ personnes dont les dates d'anniversaire sont
indépendantes et uniformes sur $ N = 365 $ jours ; ignorez les jours bissextiles et les
jumeaux.
\begin{enumerate}
\item[(a)] Démontrez que $ n = 23 $ est la plus petite taille de groupe pour laquelle un
anniversaire commun est plus probable qu'improbable. Démontrez ensuite que le seuil est
d'ordre $ \sqrt{N} $ avec la bonne constante : montrez que le plus petit tel $ n $ vaut
$ \big(\sqrt{2\ln 2} + o(1)\big)\sqrt{N} $ quand $ N \to \infty $, en démontrant des
bornes concordantes dans les deux sens plutôt qu'une seule inégalité.
\item[(b)] Soit $ W $ le nombre de personnes que vous devez interroger, une à une, jusqu'à ce
qu'un anniversaire se répète. Démontrez que
\[ \mathbb{E}[W] \;=\; 1 + \sum_{k=1}^{N} \frac{N!}{(N-k)!\,N^{k}}, \]
et établissez le développement asymptotique
$ \mathbb{E}[W] = \sqrt{\pi N/2} + \tfrac{2}{3} + o(1) $. Expliquez pourquoi le temps
d'attente espéré et le seuil médian de la question (a) diffèrent à ce point, et lequel des
deux est celui dont dépend réellement le tour de société.
\item[(c)] Abandonnez l'uniformité. Soient les anniversaires indépendants de loi arbitraire
$ p = (p_1, \dots, p_N) $ sur les $ N $ jours. Démontrez que la probabilité que les
$ n $ anniversaires soient tous distincts est maximale exactement lorsque $ p $ est
uniforme, de sorte que l'irrégularité d'un vrai calendrier ne peut que vous aider.
Identifiez le principe général --- concavité de Schur, ou majorisation --- dont votre
démonstration est un cas particulier.
\item[(d)] Faites-en un numéro. Pour une salle de $ 23 $ personnes, calculez la probabilité que
votre annonce réussisse ; calculez ensuite la probabilité que deux anniversaires tombent à
moins d'une semaine l'un de l'autre, et déterminez la plus petite taille de salle pour
laquelle \textit{ce} pari-là est favorable. Rendez quantitativement compte de l'ampleur de
l'écart entre les deux nombres.
\end{enumerate}

\end{problem}

\noindent Harold Davenport semble avoir connu le problème vers 1927 sans le publier,
disant qu'il ne pouvait croire qu'il n'eût pas été énoncé auparavant ; la première
publication est celle de Richard von Mises, en 1939. C'est un tour de société depuis lors, et
le tour est honnête --- avec vingt-trois personnes vous êtes légèrement favori et avec trente
vous l'êtes largement. Les mathématiques qui comptent sont la loi en $ \sqrt{N} $ de la
question (a), car c'est elle qui rend la réponse contre-intuitive : l'intuition compare
$ n $ à $ N $ quand elle devrait comparer $ \binom{n}{2} $ à $ N $. Le même calcul
est l'attaque des anniversaires en cryptographie, où $ N $ est le nombre de valeurs de
hachage possibles et $ 1{,}18\sqrt{N} $ le budget de l'attaquant, et c'est pourquoi une
fonction de hachage a besoin du double de la longueur de sortie que l'on exigerait
naïvement. La question (c) est celle qui empêche le tour d'être une escroquerie : un
sceptique objectera que les anniversaires ne sont pas uniformes, et la réponse honnête est
un théorème disant que l'objection joue dans l'autre sens. L'analyse par Donald Knuth du
temps d'attente espéré de la question (b), via la fonction $ Q $ de Ramanujan, est le
traitement standard de la version relative au hachage.

\begin{refs}
\item Contexte : Paradoxe des anniversaires --- Wikipédia (\url{https://fr.wikipedia.org/wiki/Paradoxe_des_anniversaires})
\item Contexte : Attaque des anniversaires --- Wikipédia (\url{https://fr.wikipedia.org/wiki/Attaque_des_anniversaires})
\item Article : D. M. Bloom, \og{} A birthday problem \fg{}, \textit{American Mathematical Monthly} 80 (1973)
\item Voir aussi : Probabilités (\url{applied-probability.html}) pour le calcul de base, et Cryptographie et théorie de l'information (\url{applied-crypto.html}) pour l'attaque
\end{refs}

\bigskip

\begin{problem}[{Le problème du vestiaire, sans inclusion-exclusion --- Licence}]
\textbf{PUZ-80.} \textbf{Problème.} Un préposé prend les chapeaux de $ n $ invités et les rend dans un
ordre uniformément aléatoire. Notez $ D_n $ le nombre de permutations de $ n $ objets
sans point fixe, de sorte que $ D_0 = 1 $ et $ D_1 = 0 $.
\begin{enumerate}
\item[(a)] Démontrez la récurrence
\[ D_n \;=\; (n-1)\big(D_{n-1} + D_{n-2}\big), \qquad n \ge 2, \]
par un argument direct sur la destination du premier chapeau --- et non en manipulant une
formule close que vous connaîtriez déjà.
\item[(b)] Démontrez que $ D_n = n D_{n-1} + (-1)^{n} $, et faites-le en exhibant une
quasi-bijection explicite : deux ensembles finis décrits combinatoirement, accompagnés d'un
appariement entre eux qui laisse exactement un élément non apparié, dont la parité rend
compte du signe. Une involution qui renverse le signe est l'outil ; une dérivation algébrique
à partir de la formule d'inclusion-exclusion ne compte pas.
\item[(c)] Soit $ F_n $ le nombre d'invités qui reçoivent leur propre chapeau. Démontrez que le
$ r $-ième moment factoriel vérifie $ \mathbb{E}\big[(F_n)_r\big] = 1 $ pour tout
$ n \ge r $, où $ (x)_r = x(x-1)\cdots(x-r+1) $. Déduisez-en que $ F_n $ converge en
loi vers une variable aléatoire de Poisson de moyenne $ 1 $, et expliquez ce que la
constance de tous les moments factoriels dit de la dépendance entre les événements
\og{} l'invité $ i $ reçoit son propre chapeau \fg{}.
\item[(d)] Le jeu de Montmort lui-même. Au \textit{treize}, les treize cartes d'une couleur sont
retournées une à une pendant que le donneur compte \og{} un, deux, \dots{}, treize \fg{} ; une
\textit{rencontre} se produit chaque fois que le compte coïncide avec la carte retournée.
Calculez la probabilité que le donneur réalise au moins une rencontre. Faites ensuite de
même pour un jeu complet de cinquante-deux cartes compté quatre fois de un à treize, et
déterminez, avec démonstration, lequel des deux jeux favorise le plus le donneur.
\end{enumerate}

\end{problem}

\noindent Pierre Rémond de Montmort a publié le problème des rencontres dans son
\textit{Essay d'analyse sur les jeux de hazard} de 1708, l'a développé dans la seconde édition
de 1713 grâce à une longue correspondance avec Nicolas Bernoulli, et on l'appelle
habituellement le premier problème de probabilités combinatoires. Tous les manuels
établissent aujourd'hui $ D_n $ par inclusion-exclusion, ce qui est correct, rapide et ne
vous apprend rien ; ce problème existe pour vous faire obtenir les mêmes faits de trois
autres manières, car une récurrence, une bijection et un calcul de moments exposent chacun
une part différente de la structure. La question (c) est la lecture moderne : les points
fixes d'une permutation aléatoire sont rares et presque indépendants, de sorte que leur
nombre est asymptotiquement de Poisson, et l'égalité à $ 1 $ de tous les moments
factoriels en est un indice aussi net que la combinatoire puisse en offrir. Que la réponse
se stabilise près de $ 1/e $ dès $ n = 7 $ stupéfie les joueurs de cartes depuis trois
siècles.

\begin{refs}
\item Contexte : Problème des rencontres (dérangements) --- Wikipédia (\url{https://fr.wikipedia.org/wiki/Probl%C3%A8me_des_rencontres})
\item Contexte : Nombres de rencontres --- Wikipédia (\url{https://fr.wikipedia.org/wiki/D%C3%A9rangement_partiel})
\item Lecture : William Feller, \textit{An Introduction to Probability Theory and Its Applications}, vol. 1, chap. IV (le problème des rencontres)
\item Voir aussi : Combinatoire et théorie des graphes (\url{pure-combinatorics.html}) pour la dérivation par inclusion-exclusion, et Probabilités (\url{applied-probability.html})
\end{refs}

\bigskip

\begin{problem}[{Les nuggets et le nombre de Frobenius --- Licence}]
\textbf{PUZ-81.} \textbf{Problème.} Fixez des entiers strictement positifs $ a_1, \dots, a_k $ de plus
grand commun diviseur $ 1 $. Dites qu'un entier $ n $ est \textit{représentable} si
$ n = x_1 a_1 + \cdots + x_k a_k $ pour certains entiers $ x_i $ positifs ou nuls. Le
\textit{nombre de Frobenius} $ g(a_1, \dots, a_k) $ est le plus grand entier non
représentable.
\begin{enumerate}
\item[(a)] Pour $ k = 2 $, avec $ a $ et $ b $ premiers entre eux, démontrez que
\[ g(a,b) = ab - a - b, \]
qu'exactement $ (a-1)(b-1)/2 $ entiers strictement positifs sont non représentables, et
que $ n $ est représentable si et seulement si $ ab - a - b - n $ ne l'est pas.
\item[(b)] Les nuggets de poulet se vendaient en boîtes de $ 6 $, $ 9 $ et $ 20 $.
Déterminez, avec démonstration, le plus grand nombre de nuggets que l'on ne peut pas acheter
exactement. Déterminez-le ensuite de nouveau après l'introduction d'une boîte de $ 4 $, et
expliquez pourquoi ajouter une dénomination peut modifier la réponse aussi violemment.
\item[(c)] Démontrez qu'aucune formule de la forme de celle de la question (a) ne peut valoir pour
trois dénominations ou plus : exhibez des triplets pour lesquels le nombre de Frobenius se
comporte irrégulièrement quand les dénominations varient régulièrement, et démontrez la
meilleure borne supérieure générale sur $ g $ que vous puissiez, en fonction de
$ a_1, \dots, a_k $. Énoncez ensuite précisément ce que l'on sait du calcul de $ g $ :
pour un nombre fixé de dénominations il existe un algorithme polynomial en le nombre de
chiffres des $ a_i $, tandis que le problème où $ k $ fait partie de l'entrée est
NP-difficile.
\item[(d)] Reformulez le tout. Les entiers représentables forment un \textit{demi-groupe numérique}
$ S \subseteq \mathbb{N} $, c'est-à-dire une partie stable par addition, contenant
$ 0 $, et de complémentaire fini. Définissez le \textit{genre} $ g(S) $ comme le nombre de
lacunes et $ F(S) $ comme le nombre de Frobenius, démontrez que l'on a toujours
$ g(S) \ge (F(S)+1)/2 $, et caractérisez les demi-groupes qui atteignent l'égalité.
Identifiez lesquels des exemples à deux générateurs de la question (a) sont de ce type.
\end{enumerate}

\end{problem}

\noindent James Joseph Sylvester a publié la formule à deux pièces en 1882, sous
forme de question et de réponse dans l'\textit{Educational Times} ; si Frobenius est resté dans
le nom, c'est qu'il avait l'habitude de mentionner le problème général dans ses cours sans
rien publier dessus, et le problème général est réellement difficile --- le seul cas
$ k = 3 $ a engendré une petite bibliothèque. La version des nuggets est la forme que
prend le casse-tête dans les écoles américaines, et c'est une bonne forme parce qu'elle rend
clair ce qui change quand on ajoute une dénomination. La question (d) est le basculement qui
transforme un casse-tête en discipline : dès que l'on regarde l'ensemble entier des nombres
représentables plutôt que la plus grande lacune, on regarde un demi-groupe numérique, et les
questions qui se présentent alors --- combien de demi-groupes ont un genre donné, lesquels sont
symétriques, à quoi ressemblent leurs lacunes --- se relient directement à la théorie des
courbes algébriques singulières, où l'ensemble des lacunes d'un point d'une courbe est la
suite de lacunes de Weierstrass. La question du dénombrement, à savoir combien de
demi-groupes numériques ont le genre $ n $, n'a été montrée croître selon une récurrence de
type Fibonacci qu'au cours des vingt dernières années.

\begin{refs}
\item Contexte : Problème des pièces de monnaie --- Wikipédia (\url{https://fr.wikipedia.org/wiki/Probl%C3%A8me_des_pi%C3%A8ces_de_monnaie})
\item Contexte : Demi-groupe numérique --- Wikipédia (\url{https://fr.wikipedia.org/wiki/Demi-groupe_num%C3%A9rique})
\item Lecture : J. L. Ramírez Alfonsín, \textit{The Diophantine Frobenius Problem} (Oxford University Press, 2005)
\item Voir aussi : Théorie des nombres (\url{pure-number-theory.html}), Mathématiques discrètes et informatique (\url{applied-discrete-cs.html})
\end{refs}

\bigskip

\begin{problem}[{Les cubes de Langford --- Licence, short}]
\textbf{PUZ-82.} \textbf{Problème.} Un \textit{appariement de Langford} d'ordre $ n $ est
une disposition en ligne des $ 2n $ nombres $ 1,1,2,2,\dots,n,n $ telle que, pour chaque
$ k $, les deux exemplaires de $ k $ aient exactement $ k $ entrées entre eux.
Déterminez, avec démonstration, exactement quels $ n $ admettent un appariement de
Langford, en donnant une construction pour ceux qui en admettent un et un argument
d'impossibilité pour les autres.
\end{problem}

\noindent C. Dudley Langford a posé le problème en 1958 après avoir regardé son
jeune fils empiler des cubes colorés et remarqué que l'enfant avait, par accident, produit la
disposition pour $ n = 3 $. La moitié \og{} impossibilité \fg{} est un argument de comptage du
genre qu'il vaut la peine d'avoir sous la main : additionnez les positions des $ 2n $
entrées de deux manières différentes et comparez les résultats modulo un petit nombre. La
moitié \og{} construction \fg{} est tout autre et exige une famille explicite de dispositions
couvrant les classes résiduelles survivantes. Ce que le problème ne vous donne \textit{pas},
c'est un dénombrement : le nombre d'appariements de Langford d'ordre $ n $ n'a aucune
formule connue, les valeurs (consignées comme suite A014552 dans l'OEIS) n'ont été obtenues
que par de gros calculs, et Donald Knuth utilise le problème comme cas de test standard des
algorithmes de couverture exacte dans \textit{The Art of Computer Programming}. C'est une
silhouette familière en combinatoire --- existence réglée par un argument élémentaire,
énumération grande ouverte --- et c'est le même écart qui sépare, en
PUZ-56 (\url{#PUZ-56}), le savoir qu'une stratégie gagnante existe de tout moyen de la
nommer.

\begin{refs}
\item Contexte : Appariement de Langford (suite de Skolem) --- Wikipédia (\url{https://fr.wikipedia.org/wiki/Suite_de_Skolem})
\item Contexte : OEIS A014552 --- le nombre d'appariements de Langford (\url{https://oeis.org/A014552})
\item Lecture : Donald E. Knuth, \textit{The Art of Computer Programming}, volume 4 (algorithmes combinatoires ; couverture exacte et liens dansants)
\item Voir aussi : Combinatoire et théorie des graphes (\url{pure-combinatorics.html})
\end{refs}

\bigskip

\begin{problem}[{Le Hex, le jeu qui ne peut pas être nul --- Licence}]
\textbf{PUZ-92.} \textbf{Problème.} Le Hex se joue sur un losange de $ n \times n $ cases hexagonales. Deux
joueurs, Noir et Blanc, colorient tour à tour une case vide dans leur propre couleur. Noir
possède une paire de côtés opposés du losange et Blanc l'autre paire ; Noir gagne en formant
une chaîne connexe de cases noires reliant ses deux côtés, Blanc en formant une chaîne
connexe de cases blanches reliant les siens.
\begin{enumerate}
\item[(a)] Démontrez le théorème du Hex : dans tout coloriage du plateau entier, exactement l'un des
deux joueurs possède une chaîne gagnante. Les deux moitiés demandent une démonstration ---
qu'au moins un joueur relie, et qu'ils ne peuvent pas relier tous les deux.
\item[(b)] Démontrez par vol de stratégie que le premier joueur possède une stratégie gagnante sur
tout plateau $ n \times n $. Isolez les deux propriétés du Hex qu'utilise votre argument ---
qu'une pierre supplémentaire de sa propre couleur n'est jamais un handicap, et que (a) exclut
la partie nulle --- puis énoncez exactement ce que l'argument ne livre pas.
\item[(c)] Démontrez que le théorème du Hex équivaut au théorème du point fixe de Brouwer en
dimension deux : déduisez Brouwer de (a), et déduisez (a) de Brouwer.
\item[(d)] Affrontez l'écart entre (b) et la pratique. Déterminez et démontrez un premier coup
gagnant sur le plateau $ 3 \times 3 $ ; passez ensuite en revue ce que l'on sait pour
$ n $ plus grand, y compris le théorème de Reisch de 1981 selon lequel décider d'une
position de Hex arbitraire est PSPACE-complet, et dites ce que ce théorème implique et
n'implique pas au sujet des plateaux sur lesquels on joue réellement.
\end{enumerate}

\end{problem}

\noindent Piet Hein a présenté le Hex à l'institut Niels-Bohr en 1942, et le
quotidien copenhaguois \textit{Politiken} l'a publié sous le nom de Polygon ; John Nash l'a
réinventé à Princeton vers 1948, où il se jouait sur le carrelage hexagonal de la salle
commune et s'appelait Nash, ou moins aimablement John. Martin Gardner l'a fait connaître en
Amérique dans \textit{Scientific American} en juillet 1957. Sa place définitive en
mathématiques repose sur les questions (b) et (c). L'argument du vol de stratégie est le
spécimen le plus pur d'une démonstration qui établit une existence sans rien fournir --- le
même état inconfortable du savoir que le Chomp en PUZ-56 --- et David Gale a montré en 1979 que
l'anodin énoncé de connexité de la question (a) est un théorème de point fixe déguisé en jeu
de plateau, de sorte qu'un passe-temps d'enfant équivaut logiquement au résultat qui fonde
l'existence des équilibres de Nash et de l'équilibre économique général. On dit que Nash
aimait ce jeu en partie pour cette raison. Les ordinateurs ont résolu exactement les plateaux
jusqu'à $ 9 \times 9 $, et pas au-delà.

\begin{refs}
\item Contexte : Hex (jeu de plateau) --- Wikipédia (\url{https://fr.wikipedia.org/wiki/Hex})
\item Contexte : Argument du vol de stratégie --- Wikipédia (en anglais) (\url{https://en.wikipedia.org/wiki/Strategy-stealing_argument})
\item Article : David Gale, \og{} The game of Hex and the Brouwer fixed-point theorem \fg{}, \textit{American Mathematical Monthly} 86 (1979), 818--827
\item Voir aussi : PUZ-56 (\url{#PUZ-56}) ci-dessus, et Topologie (\url{pure-topology.html}) pour le théorème de Brouwer
\end{refs}

\bigskip

\begin{problem}[{Chapeaux, passes et code de Hamming --- Licence}]
\textbf{PUZ-93.} \textbf{Problème.} Chacun de $ n $ joueurs reçoit un chapeau, rouge ou bleu,
indépendamment et avec probabilité $ \tfrac{1}{2} $ pour chaque couleur. Chaque joueur voit
tous les autres chapeaux mais pas le sien. Puis, simultanément et sans aucune communication,
chaque joueur soit annonce une couleur, soit passe. L'équipe gagne si au moins un joueur
annonce une couleur et si tout joueur qui annonce une couleur annonce celle de son propre
chapeau. Une stratégie est le choix, fait à l'avance et connu de tous, d'une fonction
associant aux $ n-1 $ chapeaux qu'un joueur voit l'une des trois réponses \og{} rouge \fg{},
\og{} bleu \fg{}, \og{} passe \fg{} --- une telle fonction pour chaque joueur.
\begin{enumerate}
\item[(a)] Démontrez que pour $ n = 3 $ une certaine stratégie gagne avec probabilité
$ \tfrac{3}{4} $, et démontrez qu'aucune stratégie ne fait mieux.
\item[(b)] Démontrez la borne générale : aucune stratégie à $ n $ joueurs ne gagne avec une
probabilité dépassant $ \dfrac{n}{n+1} $.
\item[(c)] Fixez une stratégie et soit $ L \subseteq \{0,1\}^{n} $ son ensemble de configurations
de chapeaux perdantes. Démontrez que toute configuration est à distance de Hamming au plus
$ 1 $ d'un élément de $ L $ --- autrement dit, $ L $ est un code couvrant binaire de
longueur $ n $ et de rayon $ 1 $. Démontrez également la construction réciproque : à
partir de n'importe quel tel code couvrant $ C $, construisez une stratégie dont l'ensemble
perdant est exactement $ C $.
\item[(d)] Déduisez-en que la plus grande probabilité de victoire vaut exactement
$ 1 - K(n,1)/2^{n} $, où $ K(n,1) $ est la plus petite taille d'un code couvrant binaire
de longueur $ n $ et de rayon $ 1 $. Démontrez que la borne de la question (b) est
atteinte précisément lorsque $ n = 2^{m} - 1 $, en identifiant le code optimal au code de
Hamming de cette longueur, puis indiquez honnêtement pour quels $ n $ la valeur de
$ K(n,1) $ est actuellement connue.
\end{enumerate}

\end{problem}

\noindent Todd Ebert a introduit cette version du jeu des chapeaux dans sa thèse de
doctorat de 1998 à l'université de Californie à Santa Barbara, au cours de travaux sur le
calcul avec conseil limité ; elle a touché un public général grâce à un article du \textit{New
York Times} en 2001 et a été pendant plusieurs années un incontournable des entretiens
techniques. Hendrik Lenstra et Gadiel Seroussi ont présenté \og{} On hats and other covers \fg{} au
symposium international de l'IEEE sur la théorie de l'information à Lausanne en 2002, où la
correspondance avec les codes couvrants de la question (c) était exposée en toute généralité,
y compris le cas à plusieurs couleurs. La raison pour laquelle ce casse-tête est stupéfiant
mérite d'être énoncée avec soin : aucun joueur pris isolément n'a la moindre information sur
son propre chapeau, et aucun ne peut faire mieux qu'un tirage à pile ou face, de sorte que
chaque parcelle de l'avantage de l'équipe doit venir de l'agencement conjoint de qui parle et
de qui se tait, et non du savoir de quiconque. Comparez PUZ-10, où les joueurs parlent à tour
de rôle et s'entendent, et PUZ-14, où ils sont en nombre infini. La question (d) est le
gain : une question récréative se révèle être exactement un problème de théorie des codes
encore ouvert en général, de sorte que le casse-tête et le problème de recherche sont
littéralement le même problème.

\begin{refs}
\item Contexte : Casse-tête des chapeaux --- Wikipédia (en anglais) (\url{https://en.wikipedia.org/wiki/Hat_puzzle}), section sur la version d'Ebert et les codes de Hamming
\item Contexte : Code couvrant --- Wikipédia (en anglais) (\url{https://en.wikipedia.org/wiki/Covering_code})
\item Article : H. W. Lenstra \& G. Seroussi, \og{} On hats and other covers \fg{}, IEEE International Symposium on Information Theory, Lausanne, 2002 --- arXiv:cs/0509045 (\url{https://arxiv.org/abs/cs/0509045})
\item Voir aussi : PUZ-10 (\url{#PUZ-10}) et PUZ-14 (\url{#PUZ-14}) ci-dessus, et Cryptographie et théorie de l'information (\url{applied-crypto.html})
\end{refs}

\bigskip

\begin{problem}[{Les boîtes d'allumettes de Banach --- Licence, short}]
\textbf{PUZ-94.} \textbf{Problème.} Un mathématicien porte une boîte d'allumettes dans
chaque poche, chacune contenant $ N $ allumettes au départ ; chaque fois qu'il veut une
allumette, il choisit une poche par un tirage à pile ou face équitable et indépendant, et
retire une allumette de la boîte qu'il y trouve. Déterminez, avec démonstration, la
probabilité qu'au premier moment où il tend la main vers une boîte et la trouve vide, l'autre
boîte contienne exactement $ k $ allumettes, et démontrez que le nombre espéré d'allumettes
alors restantes vaut $ (2N+1)\binom{2N}{N}2^{-2N} - 1 $, soit asymptotiquement
$ 2\sqrt{N/\pi} - 1 $.
\end{problem}

\noindent Hugo Steinhaus a posé ce problème dans un discours prononcé lors d'un
banquet en l'honneur de Stefan Banach, en plaisantant sur le tabagisme de Banach ; le
problème porte le nom de Banach depuis lors, bien que celui-ci ne l'ait jamais posé.
Steinhaus et Banach étaient au centre de l'école de Lwów, qui faisait une grande part de son
travail au Café écossais et inscrivait ses problèmes ouverts, dotés de prix --- une bouteille
de vin, une oie vivante ---, dans le cahier connu aujourd'hui sous le nom de Livre écossais.
La loi qui en sort est la loi binomiale négative, et le $ \sqrt{N} $ de la réponse est le
même $ \sqrt{N} $ qui régit le déplacement typique d'une marche aléatoire simple après
$ N $ pas : le compte des allumettes dans les deux poches effectue exactement une telle
marche, ce qui explique que la surprise de la réponse soit une surprise sur les marches
aléatoires plutôt que sur le tabac. Feller en donne le traitement classique.

\begin{refs}
\item Contexte : Problème des boîtes d'allumettes de Banach --- Wikipédia (en anglais) (\url{https://en.wikipedia.org/wiki/Banach%27s_matchbox_problem})
\item Contexte : Loi binomiale négative --- Wikipédia (\url{https://fr.wikipedia.org/wiki/Loi_binomiale_n%C3%A9gative})
\item Lecture : William Feller, \textit{An Introduction to Probability Theory and Its Applications}, vol. 1, chap. VI \S{}8
\item Voir aussi : Probabilités (\url{applied-probability.html}) pour les marches aléatoires
\end{refs}

\bigskip

\begin{problem}[{Les mariages stables, et qui gagne à mentir --- Licence}]
\textbf{PUZ-95.} \textbf{Problème.} Il y a $ n $ hôpitaux et $ n $ candidats. Chaque hôpital dispose
d'un classement strict des $ n $ candidats et chaque candidat d'un classement strict des
$ n $ hôpitaux. Un \textit{appariement} est une bijection entre les deux ensembles. Un
appariement est \textit{instable} si un hôpital et un candidat qui ne sont pas appariés
l'un à l'autre se préfèrent mutuellement à leur partenaire attribué ; sinon il est
\textit{stable}.
\begin{enumerate}
\item[(a)] Démontrez qu'un appariement stable existe toujours, en décrivant la procédure
d'acceptation différée --- les membres non appariés d'un côté font des propositions par ordre
de préférence, l'autre côté retient provisoirement la meilleure offre reçue jusque-là et
rejette le reste --- et en démontrant qu'elle s'arrête après au plus $ n^{2} $ propositions
sur un appariement stable.
\item[(b)] Démontrez que les appariements stables forment un treillis : étant donné $ M $ et
$ M' $ stables, définissez $ M \vee M' $ en attribuant à chaque candidat celui des deux
partenaires qu'il préfère, démontrez que c'est encore un appariement et qu'il est stable, et
identifiez les éléments maximal et minimal. Démontrez que la procédure de (a) produit
toujours l'élément optimal pour le côté qui propose et simultanément le pire pour l'autre
côté, quel que soit l'ordre dans lequel les propositions sont faites.
\item[(c)] Démontrez que, lorsque ce sont les candidats qui proposent, aucun candidat ne peut
obtenir un meilleur partenaire en soumettant un classement mensonger. Exhibez ensuite une
petite instance dans laquelle un hôpital gagne bel et bien à soumettre un classement
mensonger, et démontrez qu'aucun mécanisme retournant toujours un appariement stable ne peut
être immunisé contre une telle manipulation des deux côtés à la fois.
\item[(d)] Abandonnez la structure bipartie : $ 2n $ personnes classent chacune les $ 2n - 1 $
autres, et un appariement les apparie désormais arbitrairement. Construisez une instance à
quatre personnes n'admettant aucun appariement stable, démontrez-le, et énoncez ce que
l'algorithme d'Irving de 1985 réalise pour cette version.
\end{enumerate}

\end{problem}

\noindent David Gale et Lloyd Shapley ont publié \og{} College admissions and the
stability of marriage \fg{} dans l'\textit{American Mathematical Monthly} en 1962, en y incluant
une défense de leur article comme œuvre mathématique malgré l'absence de nombres, de formules
et de figures --- argument qui mérite d'être lu à côté du théorème. Le sel de la discipline est
que le National Resident Matching Program affectait les diplômés en médecine américains aux
hôpitaux par une procédure essentiellement équivalente depuis 1952, obtenue par négociation
et par tâtonnement plutôt que par démonstration ; Alvin Roth a identifié le lien en 1984 et a
plus tard reconçu l'appariement pour tenir compte des couples mariés, ce qui brise la
structure de treillis de la question (b) et fait échouer l'existence. Roth et Shapley ont
partagé le prix Nobel d'économie 2012 ; Gale était mort en 2008. Des variantes de
l'algorithme affectent aujourd'hui les enfants aux écoles à New York et à Boston. La question
(c) est la raison pour laquelle les économistes s'y intéressent : la stabilité est une
propriété d'un résultat et l'immunité au mensonge une propriété d'un mécanisme, et le
théorème dit que vous ne pouvez pas avoir les deux pour les deux côtés --- un résultat
véritablement négatif au fondement de l'ingénierie des marchés.

\begin{refs}
\item Contexte : Algorithme de Gale et Shapley --- Wikipédia (\url{https://fr.wikipedia.org/wiki/Algorithme_de_Gale_et_Shapley})
\item Contexte : Problème des mariages stables --- Wikipédia (\url{https://fr.wikipedia.org/wiki/Probl%C3%A8me_des_mariages_stables})
\item Article : D. Gale \& L. S. Shapley, \og{} College admissions and the stability of marriage \fg{}, \textit{American Mathematical Monthly} 69 (1962), no 1
\item Lecture : Dan Gusfield \& Robert W. Irving, \textit{The Stable Marriage Problem: Structure and Algorithms} (MIT Press, 1989) --- voir aussi Optimisation et théorie des jeux (\url{applied-optimization.html}) et Mathématiques discrètes et informatique (\url{applied-discrete-cs.html})
\end{refs}

\bigskip

\begin{problem}[{L'urne d'Ellsberg --- Licence}]
\textbf{PUZ-96.} \textbf{Problème.} Une urne contient $ 90 $ boules. Exactement $ 30 $ sont rouges ; les
$ 60 $ autres sont noires ou jaunes, dans une proportion qu'on ne vous dit pas et que vous
ne pouvez pas apprendre. Une boule sera tirée au hasard. Considérez quatre paris, chacun
rapportant une unité si l'événement indiqué se produit et rien sinon :
\[ f_1 : \text{rouge}, \qquad f_2 : \text{noire}, \qquad f_3 : \text{rouge ou jaune}, \qquad f_4 : \text{noire ou jaune}. \]
La plupart des gens à qui l'on demande de choisir déclarent une préférence stricte pour
$ f_1 $ sur $ f_2 $ doublée d'une préférence stricte pour $ f_4 $ sur $ f_3 $.
\begin{enumerate}
\item[(a)] Démontrez qu'aucune mesure de probabilité sur
$ \{\text{rouge}, \text{noire}, \text{jaune}\} $ et aucune fonction d'utilité, quelles
qu'elles soient, ne rendent ces deux préférences strictes compatibles avec la maximisation de
l'utilité espérée. Énoncez ensuite précisément lequel des postulats de Savage ce schéma
viole, et démontrez qu'il le viole.
\item[(b)] Démontrez que ce schéma ne s'explique pas par l'aversion au risque : montrez que, pour
un a priori unique et \textit{toute} utilité $ u $ strictement croissante, si concave
soit-elle, les deux préférences demeurent conjointement impossibles.
\item[(c)] Démontrez que ce schéma est compatible avec la règle maximin de Gilboa et Schmeidler,
dans laquelle un pari $ f $ est évalué par
$ \min_{P \in \mathcal{P}} \mathbb{E}_{P}[u(f)] $ pour un ensemble convexe fermé
$ \mathcal{P} $ de mesures de probabilité. Exhibez un tel $ \mathcal{P} $, puis
caractérisez tous les ensembles $ \mathcal{P} $ qui reproduisent ce schéma.
\item[(d)] Concevez une suite d'échanges qui transforme une personne affichant ce schéma en pompe à
billets, puis déterminez exactement ce qu'il faut supposer en outre de cette personne --- sur
les suites de paris qu'elle est prête à accepter, et sur la manière dont elle les traite
comme un tout --- pour que la pompe fonctionne. Cela fait, dites si vous jugez ce schéma
irrationnel, et défendez votre verdict.
\end{enumerate}

\end{problem}

\noindent Daniel Ellsberg a publié \og{} Risk, ambiguity, and the Savage axioms \fg{} dans
le \textit{Quarterly Journal of Economics} en 1961, à partir de travaux de doctorat à Harvard
menés alors qu'il était analyste à la RAND ; une décennie plus tard, il copiait et divulguait
les documents du Pentagone. Le paradoxe vise une affirmation précise, non les probabilités en
général. Les axiomes de Savage, en 1954, impliquent qu'un agent cohérent se comporte comme
s'il disposait d'une unique loi de probabilité subjective, et l'urne montre que les gens
distinguent nettement une chance qu'ils connaissent d'une chance qu'ils ne connaissent pas ---
la distinction que Frank Knight avait tracée en 1921 entre risque et incertitude, et que
Keynes avait faite la même année dans \textit{A Treatise on Probability} à l'aide d'une urne de
composition inconnue. La réponse de la discipline n'a pas été de convaincre les sujets
d'erreur, mais de bâtir des théories de la décision où les croyances sont des ensembles de
mesures ou des fonctions d'ensemble non additives : l'utilité espérée de Choquet de
Schmeidler et l'utilité espérée maximin de Gilboa et Schmeidler, toutes deux de 1989, en sont
les versions standard. Placez cela à côté de PUZ-49 : le paradoxe de Bertrand demande quelle
mesure vous avez en tête, celui d'Ellsberg demande s'il en existe seulement une.

\begin{refs}
\item Contexte : Paradoxe d'Ellsberg --- Wikipédia (\url{https://fr.wikipedia.org/wiki/Paradoxe_d%27Ellsberg})
\item Contexte : Incertitude knightienne --- Wikipédia (\url{https://fr.wikipedia.org/wiki/Th%C3%A9orie_du_risque})
\item Article : Daniel Ellsberg, \og{} Risk, ambiguity, and the Savage axioms \fg{}, \textit{Quarterly Journal of Economics} 75 (1961), 643--669
\item Article : I. Gilboa \& D. Schmeidler, \og{} Maxmin expected utility with non-unique prior \fg{}, \textit{Journal of Mathematical Economics} 18 (1989), 141--153 --- voir aussi PUZ-49 (\url{#PUZ-49}) ci-dessus
\end{refs}

\bigskip


\section*{Master}

\begin{problem}[{Une infinité de chapeaux, et l'axiome du choix --- Master}]
\textbf{PUZ-14.} \textbf{Problème.} Une infinité dénombrable de prisonniers, indexés par $ \mathbb{N} $,
portent chacun un chapeau étiqueté d'un nombre réel. Chaque prisonnier voit tous les
chapeaux sauf le sien. Simultanément, sans communication d'aucune sorte, chacun devine son
propre nombre.
\begin{enumerate}
\item[(a)] À l'aide de l'axiome du choix, construisez une stratégie sous laquelle seul un nombre
fini de prisonniers se trompent. Considérez la relation sur les suites \og{} coïncident sauf en
un nombre fini de places \fg{}, vérifiez que c'est une relation d'équivalence, et choisissez un
représentant dans chaque classe.
\item[(b)] Expliquez pourquoi cela ne contredit pas l'intuition évidente qu'un prisonnier n'a
aucune information sur son propre chapeau --- remarquez en particulier que la stratégie n'est
pas mesurable et ne donne à aucun prisonnier pris isolément une chance meilleure que le
hasard.
\item[(c)] Montrez qu'une forme de choix est véritablement nécessaire : aucune telle stratégie
n'est définissable dans ZF seul, raison pour laquelle ce casse-tête sert d'argument contre
l'acceptation non critique de l'axiome du choix.
\end{enumerate}

\end{problem}

\noindent C'est le casse-tête qui fait passer l'axiome du choix d'une commodité
technique à quelque chose de troublant. La conclusion n'est pas seulement surprenante, elle
semble violer le sens commun informationnel, et pourtant la dérivation est courte et
correcte. Ce qu'elle met à nu, c'est qu'une fonction de choix sur une infinité non
dénombrable de classes d'équivalence n'est ni un algorithme ni un objet mesurable --- c'est une
pure affirmation d'existence, et la \og{} stratégie \fg{} ne peut être exécutée par personne.
Comparez le paradoxe de Banach--Tarski, troublant exactement de la même manière et pour
exactement la même raison.

\begin{refs}
\item Contexte : Casse-tête des chapeaux --- Wikipédia (en anglais) (\url{https://en.wikipedia.org/wiki/Hat_puzzle})
\item Contexte : Axiome du choix --- Wikipédia (\url{https://fr.wikipedia.org/wiki/Axiome_du_choix})
\item Contexte : Paradoxe de Banach-Tarski --- Wikipédia (\url{https://fr.wikipedia.org/wiki/Paradoxe_de_Banach-Tarski})
\item Voir aussi : Logique, théorie des ensembles et fondements (\url{pure-foundations.html}), Analyse réelle (\url{pure-real-analysis.html})
\end{refs}

\bigskip

\begin{problem}[{Les soldats de Conway --- Master, short}]
\textbf{PUZ-15.} \textbf{Problème.} Des pions occupent tous les points du réseau situés
sous une droite horizontale, et un coup consiste pour un pion à en sauter un autre
orthogonalement vers une case vide, en retirant la pièce sautée. Démontrez qu'aucune suite
finie de coups ne place un pion cinq rangées au-dessus de la droite.
\end{problem}

\noindent La démonstration, trouvée par John Conway en 1961, est l'un des arguments
les plus admirés des mathématiques récréatives : attribuez à chaque case un poids
$ \varphi^{d} $, où $ \varphi $ est l'inverse du nombre d'or et $ d $ la distance à la
cible, choisi de sorte qu'aucun coup n'augmente jamais le poids total. L'impossibilité
découle alors d'une somme géométrique convergente. C'est l'exemple définitif d'argument par
monovariant, et il explique pourquoi quatre rangées sont atteignables et cinq ne le seront
jamais.

\begin{refs}
\item Contexte : Les soldats de Conway --- Wikipédia (en anglais) (\url{https://en.wikipedia.org/wiki/Conway%27s_Soldiers})
\item Lecture : Elwyn Berlekamp, John Conway \& Richard Guy, \textit{Winning Ways for Your Mathematical Plays}
\item Voir aussi : L'art de la démonstration (\url{proofs.html})
\end{refs}

\bigskip

\begin{problem}[{L'équidécomposabilité dans le plan et dans l'espace --- Master}]
\textbf{PUZ-28.} \textbf{Problème.} Deux polytopes sont \textit{équidécomposables} si l'on peut découper l'un en un nombre fini de morceaux polytopaux qui se réassemblent en l'autre, en n'utilisant que des translations et des rotations.
\begin{enumerate}
\item[(a)] Démontrez le théorème de Wallace--Bolyai--Gerwien : deux polygones du plan sont équidécomposables si et seulement s'ils ont la même aire.
\item[(b)] Pour un polyèdre $ P $ d'arêtes $ e $ de longueur $ \ell(e) $ et d'angles dièdres $ \theta(e) $, l'invariant de Dehn est \[ D(P) = \sum_{e} \ell(e) \otimes \theta(e) \in \mathbb{R} \otimes_{\mathbb{Z}} (\mathbb{R}/2\pi\mathbb{Z}). \] Démontrez que $ D $ est additif et inchangé par dissection : si $ P $ est découpé en polyèdres $ P_1, \ldots, P_m $, alors $ D(P) = D(P_1) + \cdots + D(P_m) $.
\item[(c)] Démontrez que $ \arccos(1/3) $ n'est pas un multiple rationnel de $ \pi $, et déduisez-en qu'un tétraèdre régulier et un cube de même volume ne sont pas équidécomposables.
\end{enumerate}

\end{problem}

\noindent Hilbert en a fait le troisième des vingt-trois problèmes qu'il a posés à Paris en 1900, et ce fut le premier d'entre eux à être réglé : son étudiant Max Dehn y a répondu négativement en 1901, moins d'un an après la conférence. La question avait un vrai poids historique, car la géométrie plane élémentaire peut définir l'aire par le seul découpage-recollage --- c'est le contenu de la question (a), démontré par William Wallace en 1807 et indépendamment par Farkas Bolyai et P. Gerwien dans les années 1830 --- et Hilbert voulait savoir si la géométrie dans l'espace pouvait définir le volume de la même façon, sans passage à la limite. Elle ne le peut pas, et l'obstruction de la question (b) est la première apparition en mathématiques d'un invariant à valeurs dans un produit tensoriel, idée qui traverse aujourd'hui la K-théorie algébrique. Jean-Pierre Sydler a démontré en 1965 que le volume et l'invariant de Dehn forment ensemble un invariant complet, et Børge Jessen a donné la formulation tensorielle moderne en 1968 ; la question correspondante en géométrie sphérique et hyperbolique est encore ouverte.

\begin{refs}
\item Contexte : Troisième problème de Hilbert --- Wikipédia (\url{https://fr.wikipedia.org/wiki/Troisi%C3%A8me_probl%C3%A8me_de_Hilbert})
\item Contexte : Invariant de Dehn --- Wikipédia (en anglais) (\url{https://en.wikipedia.org/wiki/Dehn_invariant})
\item Contexte : Théorème de Wallace-Bolyai-Gerwien --- Wikipédia (\url{https://fr.wikipedia.org/wiki/Th%C3%A9or%C3%A8me_de_Wallace-Bolyai-Gerwien})
\item Lecture : Martin Aigner \& Günter M. Ziegler, \textit{Raisonnements divins} (\textit{Proofs from THE BOOK}), chapitre \og{} Le troisième problème de Hilbert : la décomposition des polyèdres \fg{}
\end{refs}

\bigskip

\begin{problem}[{La sauterelle et les points interdits --- Master, short}]
\textbf{PUZ-29.} \textbf{Problème.} Soient $ a_1, a_2, \ldots, a_n $ des entiers strictement positifs distincts de somme $ s $, et soit $ M $ un ensemble de $ n - 1 $ entiers strictement positifs avec $ s \notin M $. Une sauterelle part de $ 0 $ et effectue $ n $ sauts vers la droite, de longueurs $ a_1, \ldots, a_n $ dans un certain ordre ; démontrez que l'ordre peut être choisi de sorte que la sauterelle ne se pose jamais sur un point de $ M $.
\end{problem}

\noindent C'est le problème 6 de la 50e Olympiade internationale de mathématiques, tenue à Brême en 2009, et il figure parmi les plus difficiles jamais posés à cette compétition : sur quelque cinq cent cinquante concurrents, une poignée seulement a produit une solution complète. Terence Tao l'a mis en ligne sur son blog trois jours plus tard comme première expérience \og{} mini-polymath \fg{}, une collaboration ouverte en ligne qui a abouti à une démonstration en quelques heures, et la transcription est un rare document public sur la manière dont un problème difficile est effectivement enfoncé. Toute démonstration connue est élémentaire et inductive, ce qui fait précisément sa difficulté : rien de plus que la récurrence ordinaire n'est nécessaire, et pourtant l'énoncé correct sur lequel faire porter la récurrence est remarquablement difficile à trouver.

\begin{refs}
\item Contexte : Olympiades internationales de mathématiques --- Wikipédia (\url{https://fr.wikipedia.org/wiki/Olympiades_internationales_de_math%C3%A9matiques})
\item Lecture : Terence Tao, \og{} IMO 2009 Q6 as a mini-polymath project \fg{} (en anglais) (\url{https://terrytao.wordpress.com/2009/07/20/imo-2009-q6-as-a-mini-polymath-project/})
\item Voir aussi : Combinatoire et théorie des graphes (\url{pure-combinatorics.html}), L'art de la démonstration (\url{proofs.html})
\end{refs}

\bigskip

\begin{problem}[{Les enfants aux fronts sales --- Master}]
\textbf{PUZ-39.} \textbf{Problème.} Sur $ n $ enfants qui ont joué dehors, exactement $ k \ge 1 $ ont le
front boueux. Chaque enfant voit tous les autres fronts mais pas le sien ; qu'ils voient tous,
et qu'ils raisonnent tous correctement, est de connaissance commune. Leur père dit, de sorte
que tous l'entendent et que tous entendent que tous l'entendent : \og{} L'un d'entre vous au moins
a le front boueux. \fg{} Puis il demande de façon répétée : \og{} Savez-vous si vous êtes boueux ?
Répondez, tous ensemble, oui ou non. \fg{}
\begin{enumerate}
\item[(a)] Démontrez que chaque enfant répond \og{} non \fg{} à chacun des $ k-1 $ premiers tours, et
qu'au tour $ k $ ce sont exactement les $ k $ enfants boueux qui répondent \og{} oui \fg{}.
\item[(b)] Construisez le modèle. Prenez pour états les $ 2^{n} $ attributions de boue, et faites
valoir la relation d'indiscernabilité $ \sim_i $ de l'enfant $ i $ entre deux états qui
coïncident sur toutes les coordonnées sauf éventuellement la $ i $-ième. Vérifiez que
chaque $ \sim_i $ est une relation d'équivalence, de sorte que le modèle valide S5 ;
définissez $ E\varphi $ comme \og{} tout enfant sait $ \varphi $ \fg{} et $ C\varphi $ via la
clôture réflexive-transitive de $ \bigcup_i \sim_i $ ; et démontrez que l'annonce du père
agit sur le modèle en supprimant exactement les états qu'elle exclut.
\item[(c)] Soit $ \varphi $ l'énoncé \og{} au moins un enfant est boueux \fg{}. Lorsque exactement
$ k $ enfants sont boueux, démontrez que $ E^{\,k-1}\varphi $ est vrai avant que le père
ne parle mais que $ E^{\,k}\varphi $ ne l'est pas, et que $ C\varphi $ est vrai
immédiatement après. Déduisez-en que, pour $ k \ge 2 $, le père n'apprend aux enfants rien
qu'aucun d'eux ne voyait déjà, et pourtant modifie le modèle.
\end{enumerate}

\end{problem}

\noindent C'est le casse-tête des enfants boueux, raconté aussi comme celui des
sages et de leurs chapeaux et, moins poliment, comme celui des épouses infidèles ; c'est
l'exemple le plus utilisé de la littérature sur le raisonnement à propos de la connaissance.
Il est entré en informatique par Halpern et Moses dans les années 1980, qui soutenaient
qu'une grande quantité de protocoles de coordination des systèmes répartis sont, analysés
honnêtement, des tentatives d'établir une connaissance commune --- ce qui relie directement ce
problème à PUZ-38, où l'on démontre que la connaissance commune ne peut pas être atteinte. La
question (b) est le gain : la connaissance modélisée par une relation d'équivalence
d'indiscernabilité, S5 comme logique modale obtenue, et une annonce publique comme
suppression d'états. Une fois que vous aurez construit le modèle à la main pour $ n = 3 $,
vous ne trouverez plus jamais PUZ-11 mystérieux, car vous pourrez montrer du doigt les états
qui disparaissent.

\begin{refs}
\item Contexte : Casse-tête à récurrence --- Wikipédia (en anglais) (\url{https://en.wikipedia.org/wiki/Induction_puzzles})
\item Contexte : Sémantique de Kripke --- Wikipédia (\url{https://fr.wikipedia.org/wiki/S%C3%A9mantique_de_Kripke})
\item Lecture : Ronald Fagin, Joseph Halpern, Yoram Moses \& Moshe Vardi, \textit{Reasoning About Knowledge} (MIT Press, 1995), chapitres 1--2
\item Voir aussi : PUZ-11 et PUZ-38 ci-dessus, et Logique, théorie des ensembles et fondements (\url{pure-foundations.html})
\end{refs}

\bigskip

\begin{problem}[{La somme et le produit --- Master}]
\textbf{PUZ-40.} \textbf{Problème.} Deux entiers $ x $ et $ y $ vérifient $ 1 < x < y $ et
$ x + y \le 100 $. On dit au mathématicien S la somme $ x+y $ et au mathématicien P le
produit $ xy $ ; tous deux connaissent ces contraintes et tous deux raisonnent parfaitement.
Ils parlent tour à tour. S : \og{} P ne peut pas connaître $ x $ et $ y $. \fg{} P : \og{} Maintenant
je connais $ x $ et $ y $. \fg{} S : \og{} Maintenant je les connais aussi. \fg{}
\begin{enumerate}
\item[(a)] Déterminez $ x $ et $ y $, avec démonstration. Il y a $ 2352 $ couples
admissibles, de sorte que l'argument doit être organisé plutôt qu'exhaustif : chacune des
trois déclarations est une condition exacte, et vous devriez énoncer les trois avant
d'éliminer quoi que ce soit.
\item[(b)] Montrez que la déclaration initiale de S est une condition portant sur $ s = x+y $
seul, et caractérisez exactement quelles sommes $ s \le 100 $ permettent à S de la faire
de façon véridique.
\item[(c)] La borne $ 100 $ n'est pas décorative. Pour une borne générale $ N $ à la place de
$ 100 $, le dialogue est une question finie ; déterminez pour quels $ N $ il détermine
encore un couple unique, et décrivez comment se comporte l'ensemble des couples survivants
quand $ N $ grandit. Dites clairement quelles parties de votre réponse sont démontrées et
lesquelles sont calculées.
\end{enumerate}

\end{problem}

\noindent Hans Freudenthal a publié ceci en 1969 dans \textit{Nieuw Archief voor
Wiskunde}, et Martin Gardner l'a présenté dans \textit{Scientific American} en décembre 1979
sous le nom qui lui est resté : le casse-tête impossible. Ce nom consigne la première
réaction de chacun, à savoir que les données sont manifestement insuffisantes --- S ne connaît
qu'une somme, P qu'un produit, et ni l'un ni l'autre n'énonce jamais un nombre. Ce qui est
transmis à la place, c'est de l'ignorance, trois fois de suite, et chaque déclaration réduit
l'ensemble des candidats. Structurellement, c'est PUZ-31 de nouveau, mais déroulé sur
$ 2352 $ couples au lieu de dix dates, raison pour laquelle il ne peut pas se faire de tête
et pour laquelle il récompense qui l'organise comme un calcul assorti d'une démonstration. La
question (b) est enchevêtrée avec la conjecture de Goldbach --- les sommes exprimables comme
somme de deux nombres premiers se comportent autrement que les autres ---, ce qui explique en
partie que la borne soit un modeste $ 100 $.

\begin{refs}
\item Contexte : Casse-tête de la somme et du produit --- Wikipédia (\url{https://fr.wikipedia.org/wiki/Casse-t%C3%AAte_de_la_somme_et_du_produit})
\item Contexte : Connaissance commune (logique) --- Wikipédia (\url{https://fr.wikipedia.org/wiki/Connaissance_commune})
\item Lecture : Martin Gardner, \og{} Mathematical Games \fg{}, \textit{Scientific American}, décembre 1979
\item Voir aussi : PUZ-31 ci-dessus, et Théorie des nombres (\url{pure-number-theory.html})
\end{refs}

\bigskip

\begin{problem}[{Les deux enveloppes, après avoir regardé --- Master}]
\textbf{PUZ-41.} \textbf{Problème.} Deux enveloppes sont posées devant vous ; l'une contient deux fois plus
d'argent que l'autre. Vous en choisissez une, l'ouvrez, et y trouvez $ a $. On vous propose
d'échanger. L'argument tentant est le suivant : l'autre enveloppe contient $ 2a $ ou
$ a/2 $, chacun avec probabilité $ \tfrac{1}{2} $, de sorte que son contenu espéré vaut
\[ \tfrac{1}{2}(2a) + \tfrac{1}{2}\!\left(\tfrac{a}{2}\right) = \tfrac{5a}{4} > a, \]
et vous devriez donc échanger --- puis, par le même raisonnement, ré-échanger à jamais.
\begin{enumerate}
\item[(a)] Construisez un espace probabilisé dans lequel la plus petite somme $ X $ suit une loi
donnée, une enveloppe est choisie au hasard, $ A $ est la somme observée et $ B $ la
somme contenue dans l'autre enveloppe. Identifiez exactement quelle étape de l'argument
tentant est injustifiée, et dites ce que $ \tfrac{1}{2} $ est supposé être.
\item[(b)] Démontrez que si $ \mathbb{E}[X] < \infty $ alors il est impossible d'avoir
$ \mathbb{E}[B \mid A = a] > a $ pour tout $ a $ du support de $ A $. Ainsi aucun a
priori propre à moyenne finie ne rend l'échange favorable à chaque observation.
\item[(c)] Posez maintenant $ \Pr[X = 2^{n}] = \dfrac{2^{n}}{3^{\,n+1}} $ pour
$ n = 0, 1, 2, \ldots $ Vérifiez que c'est bien une loi de probabilité, démontrez que
$ \mathbb{E}[B \mid A = a] > a $ pour tout $ a $ observable, et calculez
$ \mathbb{E}[X] $. Expliquez précisément pourquoi cela ne fournit pas une pompe à billets,
et quel principe de la théorie de la décision échoue.
\end{enumerate}

\end{problem}

\noindent Le paradoxe descend du paradoxe des cravates de Maurice Kraitchik, en
1943, par l'intermédiaire du jeu de portefeuilles de Martin Gardner de 1982 ; la forme aux
enveloppes est due à Barry Nalebuff à la fin des années 1980, et la loi de la question (c)
est celle de John Broome, en 1995. C'est la démonstration la plus nette, en probabilités,
qu'une espérance conditionnelle est une fonction et non un nombre, et que \og{} chacun avec
probabilité un demi \fg{} est une hypothèse de modélisation qu'il faut payer par une véritable
loi. La partie réellement troublante est la (c) : il existe bel et bien des lois honnêtes et
propres sous lesquelles l'échange est favorable pour toute somme que vous pourriez voir. Le
prix en est une moyenne infinie, après quoi le principe de dominance --- si $ B $ l'emporte
sur $ A $ dans tous les états, préférez $ B $ --- cesse tout simplement d'être valide, et
c'est là la leçon à emporter.

\begin{refs}
\item Contexte : Paradoxe des deux enveloppes --- Wikipédia (\url{https://fr.wikipedia.org/wiki/Paradoxe_des_deux_enveloppes})
\item Article : John Broome, \og{} The two-envelope paradox \fg{}, \textit{Analysis} 55 (1995)
\item Lecture : David Williams, \textit{Probability with Martingales} (Cambridge, 1991), pour l'espérance conditionnelle traitée comme il faut
\item Voir aussi : Probabilités (\url{applied-probability.html}), Optimisation et théorie des jeux (\url{applied-optimization.html})
\item Voir aussi : PUZ-45 (\url{#PUZ-45}), la version scellée de ce paradoxe, où l'argument tentant est plus facile à énoncer et tout aussi faux.
\end{refs}

\bigskip

\begin{problem}[{Le Mastermind en cinq essais --- Master, short}]
\textbf{PUZ-42.} \textbf{Problème.} Au Mastermind, le codeur choisit une suite secrète de
quatre pions parmi six couleurs, les répétitions étant autorisées, et après chaque essai on
indique au décodeur combien de pions sont de la bonne couleur à la bonne place et combien
sont de la bonne couleur à la mauvaise place. Démontrez que le décodeur dispose d'une
stratégie qui identifie le secret en au plus cinq essais, quel qu'il soit.
\end{problem}

\noindent Donald Knuth a réglé la question dans \og{} The computer as master mind \fg{} en
1976, au moyen d'une recherche minimax sur les $ 6^4 = 1296 $ secrets : sa stratégie n'a
jamais besoin d'un sixième essai et en demande $ 4{,}478 $ en moyenne. Koyama et Lai ont
montré en 1993 que la meilleure moyenne possible est $ 5625/1296 \approx 4{,}340 $,
atteinte par une stratégie qui a parfois besoin de six essais --- de sorte que le pire cas et le
cas moyen sont optimisés par des jeux différents, ce qui mérite d'être remarqué. Que cinq
soit aussi nécessaire a été établi par analyse exhaustive, non par comptage : avec au plus
quatorze réponses distinctes par essai, la borne informationnelle grossière
$ 14^{q} \ge 1296 $ ne donne que $ q \ge 3 $. Le jeu général, à $ n $ positions et
$ k $ couleurs, fait l'objet de PUZ-43 ci-dessous.

\begin{refs}
\item Contexte : Mastermind --- Wikipédia (\url{https://fr.wikipedia.org/wiki/Mastermind})
\item Article : Donald E. Knuth, \og{} The computer as master mind \fg{}, \textit{Journal of Recreational Mathematics} 9 (1976/77), 1--6
\item Voir aussi : PUZ-43 ci-dessous
\end{refs}

\bigskip

\begin{problem}[{Hackenbush, et les nombres cachés dans les jeux --- Master}]
\textbf{PUZ-53.} \textbf{Problème.} Une position de Hackenbush bleu-rouge est un graphe fini dont les arêtes
sont coloriées en bleu ou en rouge, certains de ses sommets étant rattachés au sol. Gauche
supprime une arête bleue, Droite supprime une arête rouge, et toute arête laissée sans chemin
vers le sol est effacée avec elle ; un joueur qui ne peut plus jouer perd.
\begin{enumerate}
\item[(a)] Déterminez la valeur de toute tige verticale finie --- un unique chemin d'arêtes coloriées
s'élevant du sol --- et démontrez que les valeurs obtenues sont exactement les rationnels
dyadiques. Énoncez la règle qui lit la valeur sur la suite des couleurs, et démontrez-la par
récurrence sur le nombre d'arêtes.
\item[(b)] Démontrez que les valeurs du Hackenbush bleu-rouge forment un groupe abélien totalement
ordonné pour la somme disjonctive, avec $ -H $ obtenu à partir de $ H $ en échangeant les
deux couleurs, et que $ G = 0 $ exactement lorsque le second joueur gagne $ G $.
Identifiez $ G \ge H $ à un énoncé sur celui qui gagne $ G - H $ quand Droite commence.
\item[(c)] Le Hackenbush vert, où toute arête peut être supprimée par l'un ou l'autre joueur, est
impartial. Démontrez le principe des deux-points : les branches se rencontrant en un sommet
peuvent être remplacées par une seule tige dont la longueur est la somme de Nim de leurs
valeurs de Grundy. Énoncez le principe de fusion pour les positions contenant des cycles, et
utilisez les deux, avec PUZ-52, pour évaluer une position verte contenant une boucle.
\item[(d)] Montrez que les chaînes bleu-rouge infinies réalisent tout nombre réel, et expliquez en
quel sens les nombres surréels de Conway sont le Hackenbush porté dans le transfini.
\end{enumerate}

\end{problem}

\noindent John Conway a inventé le Hackenbush, et c'est la porte par laquelle il est
entré dans les nombres surréels : dans \textit{On Numbers and Games} (1976), la même définition
récursive produit les entiers, les rationnels dyadiques, les réels, les ordinaux et des
infinitésimaux comme $ \ast $ et $ \uparrow $, tous comme valeurs de positions. Le saut
conceptuel qu'exige la question (b) est qu'un jeu n'a pas seulement à être gagné ou perdu --- il
peut avoir un \textit{nombre}, et les jeux peuvent être additionnés, opposés et comparés comme
des nombres, de sorte qu'une position compliquée s'analyse en évaluant ses morceaux séparément
et en les sommant. Donald Knuth a écrit la nouvelle \textit{Surreal Numbers} en 1974 pour
expliquer la construction, et elle reste l'une des très rares œuvres mathématiques publiées
d'abord comme fiction. Les principes des deux-points et de fusion de la question (c) montrent
la théorie impartiale de PUZ-52 logée à l'intérieur de la théorie partisane, plus riche.

\begin{refs}
\item Contexte : Hackenbush --- Wikipédia (\url{https://fr.wikipedia.org/wiki/Hackenbush})
\item Contexte : Nombre surréel --- Wikipédia (\url{https://fr.wikipedia.org/wiki/Nombre_surr%C3%A9el})
\item Lecture : John H. Conway, \textit{On Numbers and Games} (2e éd., A K Peters, 2001)
\item Lecture : Elwyn Berlekamp, John Conway \& Richard Guy, \textit{Winning Ways for Your Mathematical Plays}, vol. 1
\end{refs}

\bigskip

\begin{problem}[{Combien miser : le critère de Kelly --- Master}]
\textbf{PUZ-54.} \textbf{Problème.} Vous pariez de façon répétée sur une proposition favorable. À chaque tour
vous misez une fraction fixe $ f $ de votre richesse courante, gagnant $ b $ fois la mise
avec probabilité $ p $ et perdant la mise avec probabilité $ q = 1 - p $, indépendamment
d'un tour à l'autre, avec $ bp > q $. La richesse après $ n $ tours vaut
\[ W_n = W_0 \prod_{i=1}^{n} (1 + f R_i), \qquad R_i \in \{b, -1\} . \]
\begin{enumerate}
\item[(a)] Démontrez que $ f = 1 $ maximise $ E[W_n] $ pour tout $ n $, et que cela envoie
également $ W_n \to 0 $ presque sûrement. Concluez que maximiser la richesse espérée est
le mauvais objectif, et expliquez précisément comment une espérance peut croître sans borne
tandis que presque toute trajectoire individuelle s'effondre à néant.
\item[(b)] Pour $ f \in [0,1) $, démontrez que le taux de croissance exponentiel
\[ g(f) = \lim_{n \to \infty} \frac{1}{n} \log \frac{W_n}{W_0} \]
existe presque sûrement et vaut $ E[\log(1 + fR_1)] $. Montrez que $ g $ est strictement
concave, déterminez la fraction maximisante $ f^{*} $, et démontrez que
$ g(f^{*}) > 0 $ précisément lorsque le pari est favorable.
\item[(c)] Démontrez l'optimalité asymptotique : pour tout $ f \in [0,1) $ fixé avec
$ f \ne f^{*} $, le rapport de la richesse du parieur en $ f^{*} $ à celle du parieur en
$ f $ tend vers l'infini presque sûrement, et exponentiellement vite. Montrez ensuite que
cette optimalité est asymptotique et rien de plus : exhibez un critère à horizon fini --- par
exemple maximiser $ P(W_n > c) $ pour $ n $ fixé et une cible $ c $ fixée --- sous lequel
une fraction autre que $ f^{*} $ fait strictement mieux.
\item[(d)] Généralisez à une course à $ m $ issues mutuellement exclusives de probabilités
$ p_i $ et de cotes $ o_i $, où vous devez répartir toute votre richesse entre les issues.
Établissez la répartition optimale et montrez que le taux de croissance maximal est une
différence de deux entropies --- ce qui explique que l'article de Kelly soit un article sur le
débit d'information plutôt que sur le jeu d'argent.
\end{enumerate}

\end{problem}

\noindent John L. Kelly Jr. a écrit \og{} A new interpretation of information rate \fg{} aux
Bell Labs en 1956, dans le bâtiment de Shannon et dans la langue de Shannon : il imaginait un
parieur recevant des tuyaux par un canal bruité et a démontré que le taux maximal auquel un
capital peut être amené à croître est égal au débit d'arrivée de l'information. Leo Breiman a
démontré l'optimalité asymptotique en 1961. Ed Thorp a porté le critère aux tables de
blackjack du Nevada, puis à Wall Street, et Paul Samuelson a objecté pendant des décennies ---
dans un article resté célèbre pour être écrit entièrement en mots d'une syllabe --- que
maximiser le logarithme espéré n'est qu'une fonction d'utilité parmi d'autres et n'a aucun
titre à être \textit{l'}objectif rationnel. La question (c) est l'endroit où le débat se règle
honnêtement : le pari de Kelly remporte un concours asymptotique précis et démontrable, et en
perd d'autres.

\begin{refs}
\item Contexte : Formule de Kelly --- Wikipédia (\url{https://fr.wikipedia.org/wiki/Formule_de_Kelly})
\item Article : J. L. Kelly Jr., \og{} A new interpretation of information rate \fg{}, \textit{Bell System Technical Journal} 35 (1956), 917--926
\item Article : Leo Breiman, \og{} Optimal gambling systems for favorable games \fg{}, \textit{Proc. Fourth Berkeley Symposium on Mathematical Statistics and Probability} (1961)
\item Lecture : Thomas M. Cover \& Joy A. Thomas, \textit{Elements of Information Theory}, chap. 6 (Gambling and Data Compression)
\end{refs}

\bigskip

\begin{problem}[{Deux réponses sur un même grand cercle --- Master, short}]
\textbf{PUZ-55.} \textbf{Problème.} Un point est choisi uniformément sur la sphère unité ;
notez $ \varphi \in [-\pi/2, \pi/2] $ sa latitude et $ \lambda \in [-\pi, \pi) $ sa
longitude. Démontrez que conditionner par $ \varphi = 0 $ laisse $ \lambda $ uniforme,
tandis que conditionner par $ \lambda = 0 $ donne à $ \varphi $ la densité
$ \tfrac{1}{2}\cos\varphi $ --- alors même qu'une rotation de la sphère amène la première
courbe sur un grand cercle passant par les pôles --- puis rendez compte précisément de l'écart,
en termes de probabilité conditionnelle régulière.
\end{problem}

\noindent C'est le paradoxe de Borel--Kolmogorov, et Kolmogorov l'a placé dans les
\textit{Grundbegriffe} de 1933 comme la raison pour laquelle ses axiomes définissent la
probabilité conditionnelle de la manière dont ils le font : une loi conditionnelle se
définit relativement à une tribu, ou $ \sigma $-algèbre, jamais relativement à un
événement négligeable isolé.
Latitude et longitude plongent le même cercle dans deux familles différentes de voisinages
qui rétrécissent --- de fins anneaux dans un cas, de fins fuseaux dans l'autre --- et les lois
limites diffèrent parce que les questions diffèrent. C'est le paradoxe de Bertrand (PUZ-49)
sous forme infinitésimale, et c'est la raison honnête pour laquelle la théorie de la mesure
définit l'espérance conditionnelle par une propriété de Radon--Nikodym plutôt que par un
quotient.

\begin{refs}
\item Contexte : Paradoxe de Borel--Kolmogorov --- Wikipédia (\url{https://fr.wikipedia.org/wiki/Paradoxe_de_Borel})
\item Contexte : Probabilité conditionnelle régulière --- Wikipédia (en anglais) (\url{https://en.wikipedia.org/wiki/Regular_conditional_probability})
\item Lecture : David Williams, \textit{Probability with Martingales}, chap. 9 (l'espérance conditionnelle)
\end{refs}

\bigskip

\begin{problem}[{Le chat qui tombe, et la rotation sans moment cinétique --- Master}]
\textbf{PUZ-68.} \textbf{Problème.} Modélisez un chat par deux cylindres rigides identiques et axisymétriques
joints en un point de la colonne vertébrale, capables de fléchir l'un par rapport à l'autre
d'un angle $ \theta $ et de se tordre autour de leurs propres axes longs d'un angle relatif
$ \psi $. L'animal est lâché au repos avec un moment cinétique total nul, et la seule force
extérieure est une pesanteur uniforme, qui n'exerce aucun couple autour du centre de masse.
Notez $ S $ l'espace des formes des configurations internes $ (\theta,\psi) $ et $ Q $
l'espace des configurations complet, de sorte que $ Q \to S $ est un
$ SO(3) $-fibré principal.
\begin{enumerate}
\item[(a)] Démontrez que la contrainte $ \vec{L} = 0 $ détermine une connexion principale sur
$ Q \to S $ : montrez qu'en chaque configuration elle exprime la vitesse angulaire globale
comme fonction linéaire des vitesses de forme, et calculez explicitement cette application
linéaire pour le modèle à deux cylindres.
\item[(b)] Calculez la courbure de la connexion et montrez qu'elle ne s'annule pas. Déduisez-en
qu'une boucle fermée dans l'espace des formes a une holonomie non triviale dans
$ SO(3) $, et déterminez la boucle qui maximise la rotation nette par unité de longueur
pour la métrique d'énergie cinétique sur $ S $.
\item[(c)] Démontrez que la rotation obtenue le long d'une boucle donnée est indépendante de la
vitesse à laquelle la boucle est parcourue. Déduisez-en que le phénomène ne peut être retrouvé
par aucun argument traitant la conservation du moment cinétique comme un simple énoncé
scalaire de comptabilité.
\item[(d)] Déterminez le nombre minimal de degrés de liberté de forme indépendants dont un corps
déformable a besoin pour atteindre une orientation arbitraire à moment cinétique nul, et
démontrez votre borne.
\end{enumerate}

\end{problem}

\noindent Étienne-Jules Marey a photographié un chat lâché en l'air à cadence rapide
en 1894 et a présenté les plaques à l'Académie des sciences à Paris, où l'accueil fut une
franche incrédulité : chacun dans la salle voyait bien qu'un corps sans moment cinétique ne
peut pas tourner, et en voici un qui tournait ; le chat devait donc avoir triché en prenant
appui sur les mains de l'expérimentateur. Le rédacteur de \textit{Nature} observa, à propos
d'une des plaques, que l'air de dignité offensée qu'affiche le chat à la fin de la première
série trahit un certain manque d'intérêt pour la recherche scientifique. Thomas Kane et M. P.
Scher ont rendu précis le modèle à deux cylindres en 1969, et Richard Montgomery lui a donné
sa forme moderne en 1993 en reconnaissant dans la condition de moment cinétique nul une
connexion, et dans le demi-tour net du chat une holonomie --- un problème de Yang--Mills sur
l'espace des formes. La même structure réapparaît partout où l'état interne d'un système est
cyclique alors que son état externe ne l'est pas : la phase de Berry en mécanique quantique,
la nage des micro-organismes à nombre de Reynolds nul, le contrôle d'attitude d'un satellite
à court d'ergols, et la vrille qu'un plongeur ajoute après avoir quitté le plongeoir. La
question (d) est celle à laquelle un ingénieur spatial a réellement besoin d'une
réponse.

\begin{refs}
\item Contexte : Problème du chat qui tombe --- Wikipédia (en anglais) (\url{https://en.wikipedia.org/wiki/Falling_cat_problem})
\item Contexte : Phase géométrique --- Wikipédia (\url{https://fr.wikipedia.org/wiki/Phase_g%C3%A9om%C3%A9trique})
\item Article : T. R. Kane \& M. P. Scher, \og{} A dynamical explanation of the falling cat phenomenon \fg{}, \textit{International Journal of Solids and Structures} 5 (1969), 663--670
\item Article : R. Montgomery, \og{} Gauge theory of the falling cat \fg{}, dans M. J. Enos (dir.), \textit{Dynamics and Control of Mechanical Systems}, American Mathematical Society (1993), 193--218
\end{refs}

\bigskip

\begin{problem}[{Le tourniquet inversé de Feynman --- Master}]
\textbf{PUZ-69.} \textbf{Problème.} Une tête d'arroseur en forme de S est immergée dans un grand bac et
montée sur un palier à faible frottement, libre de tourner autour d'un axe vertical. En
fonctionnement direct, l'eau est pompée à travers les bras et éjectée des buses au débit
volumique $ Q $ ; en fonctionnement inversé, l'eau est aspirée par les mêmes buses au même
débit $ Q $.
\begin{enumerate}
\item[(a)] Calculez le couple stationnaire en fonctionnement direct à partir du flux de moment
cinétique à travers une surface de contrôle enveloppant la tête, et montrez qu'il redonne la
réponse élémentaire de la \og{} réaction du jet \fg{}.
\item[(b)] Refaites le calcul par volume de contrôle en fonctionnement inversé. Montrez que, pour un
fluide parfait non visqueux à profils de vitesse uniformes, le flux de moment cinétique à
travers la surface de contrôle donne un couple stationnaire nul, puis énumérez chacune des
idéalisations faites par la dérivation --- traitez cette liste comme le véritable produit de la
question.
\item[(c)] Autorisez maintenant un profil non uniforme à l'écoulement dans un conduit courbe, de
sorte que le fluide porte la circulation secondaire qu'engendre la courbure. Déterminez le
signe du couple stationnaire qui en résulte, montrez que le flux de moment cinétique à
travers l'admission ne s'annule plus, et estimez son ordre de grandeur en fonction de la
courbure du conduit, de l'alésage, du débit $ Q $ et de la viscosité.
\item[(d)] Concevez une expérience capable de distinguer la prédiction de la question (b) de celle
de la question (c). Précisez le frottement du palier, le débit, le fluide de travail et la
résolution angulaire dont vous auriez besoin, et dites quel résultat nul falsifierait
(c).
\end{enumerate}

\end{problem}

\noindent Ernst Mach a posé la question dans \textit{Die Mechanik in ihrer
Entwicklung} (1883) et a rapporté n'avoir observé \og{} aucune rotation nette \fg{}. Feynman,
doctorant à Princeton au début des années 1940, en a débattu dans le salon de thé de la
bibliothèque Palmer, n'a pas su trancher, et a fini par brancher une bonbonne d'eau sur
l'alimentation en air du cyclotron ; la bouteille a explosé, l'eau a tout arrosé et tout le
monde avec, et l'expérience était terminée. Il raconte l'histoire dans \textit{Surely You're
Joking, Mr. Feynman!} (1985) avec les arguments des deux camps et, ostensiblement, sans le
résultat, ce qui a déclenché quatre décennies d'articles ; Edward Creutz, qui était présent,
a rapporté plus tard un léger frémissement initial, puis plus rien. Alejandro Jenkins a écrit
des traitements élémentaires soigneux en 2004 et 2011. La question a été tranchée
expérimentalement en janvier 2024 par Kaizhe Wang, Brennan Sprinkle, Mingxuan Zuo et Leif
Ristroph, à l'université de New York, qui ont construit un arroseur sur un palier à
frottement ultra-faible, mesuré directement le champ d'écoulement interne, et attribué un
petit couple stationnaire au comportement de l'écoulement dans les bras courbes. Leur article
est cité ci-dessous et --- vous voilà prévenu --- son titre vend la mèche : engagez-vous par écrit
sur une prédiction avant de l'ouvrir.

\begin{refs}
\item Contexte : Tourniquet de Feynman --- Wikipédia (\url{https://fr.wikipedia.org/wiki/Tourniquet_de_Feynman})
\item Article : A. Jenkins, \og{} An elementary treatment of the reverse sprinkler \fg{}, \textit{American Journal of Physics} 72 (2004), 1276--1282 --- arXiv:physics/0312087 (\url{https://arxiv.org/abs/physics/0312087})
\item Article : K. Wang, B. Sprinkle, M. Zuo \& L. Ristroph, \og{} Centrifugal flows drive reverse rotation of Feynman's sprinkler \fg{}, \textit{Physical Review Letters} 132 (2024), 044003
\item Lecture : R. P. Feynman, \textit{Surely You're Joking, Mr. Feynman!} (1985) --- la provenance d'origine, et aucune solution
\end{refs}

\bigskip

\begin{problem}[{La toupie tippe-top se retourne --- Master, short}]
\textbf{PUZ-70.} \textbf{Problème.} Une toupie tippe-top est une sphère tronquée de rayon
$ R $ munie d'une courte tige, dont le centre de masse se trouve à une distance $ a $ du
centre géométrique, du côté de la tige ; lancée rapidement sur une table de coefficient de
frottement de glissement $ \mu $, elle se redresse, s'inverse, et finit par tourner sur sa
tige, son centre de masse \textit{plus haut} qu'au départ. En prenant le modèle standard --- une
sphère roulant et glissant sur un plan sous frottement sec au point de contact --- démontrez
que la quantité de Jellett $ \lambda = -\vec{L}\cdot\hat{a} $, où $ \hat{a} $ est le
vecteur unitaire allant du point de contact vers le centre de masse, est conservée, et
déterminez la condition exacte, portant sur $ \lambda $, sur $ a/R $ et sur les moments
d'inertie principaux, sous laquelle l'état droit est instable tandis que l'état inversé est
asymptotiquement stable.
\end{problem}

\noindent Helene Sperl, infirmière allemande, a breveté une \textit{Wendekreisel} en
1891 ; le jouet est devenu une obsession de physiciens dans les années 1950, et une
photographie prise lors de l'inauguration de l'Institut de physique de Lund en 1954 montre
Niels Bohr et Wolfgang Pauli penchés sur une table à en regarder une, arborant la même mine
défaite. Si c'est plus qu'un jouet, c'est à cause du bilan énergétique : le centre de masse
s'élève, la toupie fournit donc un travail contre la pesanteur, tandis que l'énergie totale
décroît de façon monotone sous l'effet du frottement --- le frottement fait donc quelque chose
de plus intéressant que simplement dissiper, et trouver quoi est tout le problème. Richard
Cohen a rendu quantitatif le compte rendu par frottement de glissement dans l'\textit{American
Journal of Physics} en 1977 ; Nawaf Bou-Rabee, Jerrold Marsden et Louis Romero ont refondu
en 2004 l'inversion en \textit{instabilité induite par dissipation}, où ajouter un peu
d'amortissement à un équilibre stabilisé gyroscopiquement en détruit la stabilité. Thomson et
Tait connaissaient l'effet dans les années 1870, et il a coûté de l'argent depuis :
Explorer 1, lancé en 1958 en rotation autour de son axe de moindre inertie, a dissipé de
l'énergie dans ses antennes souples et a basculé en rotation plate en quelques heures,
enseignant à toute une génération d'ingénieurs la règle de l'axe majeur.

\begin{refs}
\item Contexte : Toupie tippe-top --- Wikipédia (\url{https://fr.wikipedia.org/wiki/Toupie_tippe-top})
\item Contexte : Mécanique du solide --- Wikipédia (\url{https://fr.wikipedia.org/wiki/M%C3%A9canique_du_solide})
\item Article : R. J. Cohen, \og{} The tippe top revisited \fg{}, \textit{American Journal of Physics} 45 (1977), 12--17
\item Article : N. Bou-Rabee, J. E. Marsden \& L. A. Romero, \og{} Tippe top inversion as a dissipation-induced instability \fg{}, \textit{SIAM Journal on Applied Dynamical Systems} 3 (2004) --- exemplaire de l'auteur (\url{https://www.cds.caltech.edu/\textasciitilde{}marsden/bib/2004/04-BoMaRo2004/BoMaRo2004_tippe.pdf})
\end{refs}

\bigskip

\begin{problem}[{Un carré et un nombre impair de triangles --- Master}]
\textbf{PUZ-83.} \textbf{Problème.} Une \textit{équidissection} d'un polygone en $ m $ morceaux est une
décomposition en $ m $ triangles d'intérieurs disjoints, de même aire, dont la réunion est
le polygone. Les sommets des triangles peuvent se trouver n'importe où ; ils n'ont pas à se
raccorder arête contre arête.
\begin{enumerate}
\item[(a)] Montrez que le carré unité admet une équidissection en $ m $ triangles pour tout
$ m $ pair, et vérifiez directement, à la main, que $ m = 3 $ est impossible.
\item[(b)] Démontrez que la valuation $ 2 $-adique $ \nu_2 $ sur $ \mathbb{Q} $ se prolonge en
une valuation $ \nu : \mathbb{R}^{\times} \to \mathbb{Q} $ --- vous aurez besoin du lemme de
Zorn. Servez-vous-en pour colorier chaque point $ (x,y) $ du plan de l'une de trois
couleurs, selon la comparaison de $ \nu(x) $, $ \nu(y) $ et $ 0 $, et démontrez que ce
coloriage est invariant par translation selon tout vecteur dont les coordonnées ont une
valuation strictement positive.
\item[(c)] Démontrez la moitié combinatoire. Montrez que, le long du côté inférieur du carré, toute
subdivision en segments contient un nombre impair de segments dont les deux extrémités
portent les deux premières couleurs, et déduisez-en, par un argument à la manière du lemme de
Sperner, que toute dissection du carré en triangles contient un triangle dont les trois
sommets portent trois couleurs différentes.
\item[(d)] Achevez la démonstration du théorème de Monsky : il n'existe aucune équidissection d'un
carré en un nombre impair de triangles. Étendez-la ensuite : démontrez par un argument affine
que la même chose vaut pour tout parallélogramme, et pour tout polygone à symétrie
centrale.
\item[(e)] Démontrez l'une des généralisations en dimension supérieure --- le théorème de Mead de
1979, selon lequel une équidissection du cube de dimension $ d $ en $ m $ simplexes force
$ d! \mid m $, ou le théorème de Kasimatis de 1989, selon lequel un $ n $-gone régulier
avec $ n \ge 5 $ n'admet une équidissection en $ m $ triangles que si $ n \mid m $.
\end{enumerate}

\end{problem}

\noindent Fred Richman a lancé tout cela à l'université d'État du Nouveau-Mexique en
1965, en rédigeant un sujet d'examen de master : il voulait une question sur le découpage
d'un carré en triangles de même aire, a constaté qu'il savait le faire avec n'importe quel
nombre pair et ne savait pas le faire avec trois ni avec cinq, et ne voyait pas pourquoi.
John Thomas a réglé en 1968 le cas des sommets à coordonnées rationnelles de dénominateur
impair, et Paul Monsky a démontré l'énoncé complet dans l'\textit{American Mathematical
Monthly} en 1970. La démonstration est célèbre pour ses ingrédients plutôt que pour sa
longueur. Une moitié est le lemme de Sperner, énoncé de point fixe purement combinatoire ;
l'autre est une valuation sur les réels prolongeant la valuation $ 2 $-adique sur
$ \mathbb{Q} $, qui n'existe que par le lemme de Zorn et n'est donc pas un objet que
quiconque puisse écrire. Qu'une question complètement élémentaire sur le découpage d'un carré
exige l'axiome du choix dans son unique démonstration connue est le genre de fait qui change
la manière dont on pense aux questions élémentaires --- et aucune démonstration sans choix
n'est connue. Comparez PUZ-28 (\url{#PUZ-28}), où un invariant différent, à valeurs dans
un produit tensoriel, règle le troisième problème de Hilbert.

\begin{refs}
\item Contexte : Théorème de Monsky --- Wikipédia (\url{https://fr.wikipedia.org/wiki/Th%C3%A9or%C3%A8me_de_Monsky})
\item Contexte : Équidissection --- Wikipédia (\url{https://fr.wikipedia.org/wiki/%C3%89quidissection})
\item Contexte : Lemme de Sperner --- Wikipédia (\url{https://fr.wikipedia.org/wiki/Lemme_de_Sperner})
\item Article : Paul Monsky, \og{} On dividing a square into triangles \fg{}, \textit{American Mathematical Monthly} 77 (1970), 161--164
\item Lecture : Martin Aigner \& Günter M. Ziegler, \textit{Raisonnements divins} (\textit{Proofs from THE BOOK}), le chapitre sur un carré et un nombre impair de triangles
\end{refs}

\bigskip

\begin{problem}[{Sept chapeaux, une réponse chacun, et le droit de se taire --- Master, short}]
\textbf{PUZ-84.} \textbf{Problème.} Sept joueurs reçoivent chacun un chapeau rouge ou bleu,
les couleurs étant tirées indépendamment et uniformément au hasard ; chaque joueur voit les
six autres chapeaux mais pas le sien, et aucune communication n'est permise une fois les
chapeaux posés. À un signal, les sept, simultanément, soit annoncent une couleur, soit
passent, et l'équipe gagne si au moins un joueur annonce une couleur et si toute couleur
annoncée est celle du joueur qui l'annonce : déterminez, avec démonstration, la plus grande
probabilité de victoire qu'une stratégie convenue à l'avance permette d'atteindre, et
démontrez qu'aucune stratégie ne fait mieux.
\end{problem}

\noindent Todd Ebert a introduit cette version dans sa thèse de doctorat de 1998 à
l'université de Californie à Santa Barbara, et elle s'est répandue dans la communauté
mathématique avec une rapidité inhabituelle, parce que la réponse est bien meilleure que la
borne évidente ne le laisse croire --- aucun joueur pris isolément ne peut faire mieux que le
hasard sur son propre chapeau, et pourtant l'équipe le peut. C'est le même phénomène qu'en
PUZ-09 (\url{#PUZ-09}) : les probabilités de succès individuelles ne peuvent pas être
améliorées, mais leur \textit{corrélation} peut être fabriquée, de sorte que les échecs de
l'équipe s'entassent sur le moins de configurations de chapeaux possible. Compter ces
configurations transforme le casse-tête en une question sur l'efficacité avec laquelle
l'espace des $ 2^{7} $ coloriages peut être couvert par des boules, ce qui est l'objet des
codes couvrants, et les nombres $ 2^{k}-1 $ ne sont pas une coïncidence. Le problème des
pronostics sportifs --- combien de grilles au minimum garantissent au plus $ R $ résultats
faux --- est la même question sous un autre costume, et il est non résolu pour la plupart des
paramètres.

\begin{refs}
\item Contexte : Casse-tête des chapeaux --- Wikipédia (en anglais) (\url{https://en.wikipedia.org/wiki/Hat_puzzle}) (attention : donne la stratégie)
\item Contexte : Code couvrant --- Wikipédia (en anglais) (\url{https://en.wikipedia.org/wiki/Covering_code}), y compris le problème des pronostics sportifs
\item Contexte : Code de Hamming --- Wikipédia (\url{https://fr.wikipedia.org/wiki/Code_de_Hamming})
\item Lecture : Peter Winkler, \textit{Mathematical Puzzles: A Connoisseur's Collection}
\item Voir aussi : PUZ-10 (\url{#PUZ-10}) et PUZ-14 (\url{#PUZ-14}) ci-dessus, et Cryptographie et théorie de l'information (\url{applied-crypto.html})
\end{refs}

\bigskip

\begin{problem}[{Le lion et l'homme --- Master, short}]
\textbf{PUZ-97.} \textbf{Problème.} Un lion et un homme, chacun réduit à un point, se
déplacent le long de trajectoires continues à l'intérieur d'un disque fermé de rayon
$ 1 $ ; tous deux ont une vitesse au plus égale à $ 1 $, chacun voit l'autre à tout
instant, et le lion attrape l'homme si leurs positions coïncident jamais. Déterminez, avec
démonstration, si le lion peut forcer la capture en temps fini, s'il peut forcer la distance
entre eux à tendre vers zéro, et si l'homme peut se garantir de n'être jamais attrapé --- en
énonçant d'abord, et exactement, ce dont une stratégie a le droit de dépendre, car les
réponses dépendent de cette définition.
\end{problem}

\noindent Le problème est dû à Richard Rado et circulait à Cambridge dans les années
1930. Pendant une vingtaine d'années, l'opinion établie fut que le lion gagne, par l'argument
de la \og{} courbe de poursuite \fg{} où il se maintient sur le rayon joignant le centre à l'homme et
referme l'écart ; Littlewood consigne l'argument, puis consigne le renversement qu'en opéra
Abram Besicovitch en 1952, dans \textit{A Mathematician's Miscellany} --- l'un des meilleurs
passages courts de la littérature sur la manière dont un argument convaincant peut être faux.
La leçon moderne est la parenthèse de l'énoncé. Bollobás, Leader et Walters ont montré, dans
\og{} Lion and Man --- Can Both Win? \fg{}, que sous des formalisations naturelles mais différentes de
ce sur quoi chaque joueur a le droit de conditionner, et de qui est réputé jouer le premier
dans l'approximation discrète, on peut munir à la fois le lion et l'homme de stratégies
gagnantes --- de sorte que la question classique n'est pas bien posée tant que la structure
d'information n'a pas été écrite. C'est un danger général des jeux de poursuite, et c'est
pourquoi Rufus Isaacs, en écrivant \textit{Differential Games} à la RAND dans les années 1950,
a consacré une si grande part du livre à la formalisation avant de poursuivre quoi que ce
soit.

\begin{refs}
\item Contexte : Poursuite-évasion --- Wikipédia (en anglais) (\url{https://en.wikipedia.org/wiki/Pursuit%E2%80%93evasion})
\item Contexte : Problème du chauffeur homicide --- Wikipédia (en anglais) (\url{https://en.wikipedia.org/wiki/Homicidal_chauffeur_problem})
\item Lecture : J. E. Littlewood, \textit{A Mathematician's Miscellany} (1953) --- le passage sur le lion et l'homme
\item Article : B. Bollobás, I. Leader \& M. Walters, \og{} Lion and Man --- Can Both Win? \fg{}, arXiv:0909.2524 (\url{https://arxiv.org/abs/0909.2524})
\end{refs}

\bigskip

\begin{problem}[{Aucune façon équitable de compter les voix --- Master}]
\textbf{PUZ-98.} \textbf{Problème.} Soit $ A $ un ensemble fini d'options avec $ |A| \ge 3 $ et soit
$ N = \{1, \dots, n\} $ un ensemble fini d'électeurs. Un \textit{profil} attribue à chaque
électeur un ordre total strict sur $ A $ ; une \textit{fonction de bien-être social} $ F $
attribue à chaque profil un ordre total strict sur $ A $. On dit que $ F $ est
\textit{unanime} si, dès que tout électeur classe $ a $ au-dessus de $ b $, $ F $ le fait
aussi ; que $ F $ satisfait l'\textit{indépendance des options non pertinentes} si le
classement social relatif de $ a $ et $ b $ ne dépend que de la façon dont les électeurs
classent $ a $ par rapport à $ b $ ; et que $ F $ est \textit{dictatoriale} si l'ordre
d'un électeur fixé est toujours l'ordre social.
\begin{enumerate}
\item[(a)] Démontrez le théorème d'Arrow : toute fonction de bien-être social unanime satisfaisant
l'indépendance des options non pertinentes est dictatoriale. Donnez la démonstration
complète par la méthode de l'électeur pivot, y compris le lemme extrémal selon lequel une
option classée en tête ou en queue par tous les électeurs doit être classée en tête ou en
queue socialement.
\item[(b)] Démontrez que les trois hypothèses sont nécessaires, en construisant, pour chacune, une
fonction de bien-être social qui satisfait les deux autres et la viole.
\item[(c)] Démontrez que le théorème exige véritablement $ |A| \ge 3 $ : montrez que, pour deux
options et un nombre impair d'électeurs, la règle de la majorité simple satisfait toutes les
conditions, et démontrez le théorème de May qui la caractérise. Puis désignez la ligne exacte
de votre démonstration de (a) qui échoue lorsque $ |A| = 2 $.
\item[(d)] Démontrez le théorème de Gibbard--Satterthwaite --- toute règle qui sélectionne un unique
vainqueur parmi au moins trois options, qui est surjective et qui n'est pas dictatoriale peut
être manipulée par un électeur déclarant un classement mensonger --- et déduisez-le de (a).
Déterminez enfin laquelle des hypothèses d'Arrow une règle par notes, comme le vote par
valeurs, met en défaut, et démontrez votre réponse.
\end{enumerate}

\end{problem}

\noindent Kenneth Arrow a démontré ceci alors qu'il était doctorant et l'a publié
sous le titre \og{} A difficulty in the concept of social welfare \fg{} dans le \textit{Journal of
Political Economy} en 1950, avant de le développer l'année suivante en \textit{Social Choice
and Individual Values} ; il a reçu le prix Nobel d'économie en 1972, à cinquante et un
ans, et reste le plus jeune économiste à l'avoir obtenu. Le théorème est la forme mûre du
cycle de PUZ-91 : Condorcet a montré qu'une règle particulière peut ne pas produire d'ordre,
et Arrow a montré que cet échec est structurel plutôt qu'un défaut de la règle majoritaire.
Le lire correctement importe, et les contresens sont fréquents. Il porte sur des règles dont
l'entrée est un ensemble de classements ordinaux, de sorte que les échappatoires
habituelles --- admettre de l'information cardinale, restreindre le domaine aux préférences
unimodales comme le fit Duncan Black en 1948, ou autoriser le tirage au sort --- reviennent
toutes à nier une hypothèse plutôt qu'à trouver une faille. John Geanakoplos a publié en 2005
trois démonstrations d'une page chacune, dont la deuxième est la voie la plus efficace vers
la question (a) ; le paradoxe libéral d'Amartya Sen montre le peu qu'il faut pour fabriquer
des impossibilités de cette forme.

\begin{refs}
\item Contexte : Théorème d'impossibilité d'Arrow --- Wikipédia (\url{https://fr.wikipedia.org/wiki/Th%C3%A9or%C3%A8me_d%27impossibilit%C3%A9_d%27Arrow})
\item Contexte : Théorème de Gibbard-Satterthwaite --- Wikipédia (\url{https://fr.wikipedia.org/wiki/Th%C3%A9or%C3%A8me_de_Gibbard-Satterthwaite})
\item Article : John Geanakoplos, \og{} Three brief proofs of Arrow's impossibility theorem \fg{}, \textit{Economic Theory} 26 (2005), 211--215
\item Lecture : Kenneth J. Arrow, \textit{Social Choice and Individual Values} (2e éd., Yale, 1963) --- voir aussi PUZ-91 (\url{#PUZ-91}) ci-dessus
\end{refs}

\bigskip


\section*{Recherche}

\begin{problem}[{Un carré magique 3$\times$3 de carrés distincts --- Recherche}]
\textbf{PUZ-16.} \textbf{Problème.} Un carré magique est un tableau $ 3 \times 3 $ dont les trois lignes,
les trois colonnes et les deux diagonales ont toutes la même somme.
\begin{enumerate}
\item[(a)] Démontrez le théorème de structure de base : dans tout carré magique $ 3 \times 3 $, la
case centrale vaut un tiers de la somme magique, et les huit autres entrées s'apparient
autour d'elle.
\item[(b)] \textbf{Problème ouvert.} Existe-t-il un carré magique $ 3 \times 3 $ dont les neuf
entrées sont des carrés parfaits distincts ? On n'en a jamais trouvé aucun, et aucune
démonstration d'impossibilité n'est connue. Martin Gardner offrait \textbackslash{}\$100 pour un tel
carré.
\item[(c)] Trouvez les meilleurs résultats partiels que vous pourrez : le \og{} carré de Parker \fg{} est un
échec de peu resté célèbre, dans lequel une condition tombe en défaut, et l'on connaît des
carrés dont sept des neuf entrées sont des carrés. Expliquez le lien avec les points
rationnels des courbes elliptiques, qui explique la difficulté véritable de cette question
innocente.
\end{enumerate}

\end{problem}

\noindent C'est le problème ouvert le plus abordable du site : un enfant peut
vérifier un candidat, et personne sur Terre ne sait en produire un ni l'exclure. Le lien avec
les courbes elliptiques de la question (c) est ce qui l'élève --- les contraintes définissent
une surface dont les points rationnels sont gouvernés par la géométrie arithmétique, ce qui
place la question dans le même voisinage que la conjecture de Birch et Swinnerton-Dyer. Le
\og{} carré de Parker \fg{}, tentative ratée de Matt Parker, est devenu une mascotte d'Internet pour
l'échec enthousiaste, emblème qui convient bien à une page de casse-tête : c'est la tentative
qui compte.

\begin{refs}
\item Contexte : Carré magique de carrés --- Wikipédia (\url{https://fr.wikipedia.org/wiki/Carr%C3%A9_magique_de_carr%C3%A9s})
\item Contexte : Courbe elliptique --- Wikipédia (\url{https://fr.wikipedia.org/wiki/Courbe_elliptique})
\item Voir aussi : Problèmes ouverts (\url{open-problems.html}), Théorie des nombres (\url{pure-number-theory.html})
\end{refs}

\bigskip

\begin{problem}[{Le plus grand canapé qui passe le coin --- Recherche}]
\textbf{PUZ-30.} \textbf{Problème.} Un couloir de largeur unité tourne à angle droit. Un \textit{canapé} est une région plane connexe que l'on peut déplacer continûment --- par translation et rotation, en restant à tout instant à l'intérieur du couloir --- du corridor d'arrivée vers celui de sortie. La \textit{constante du canapé} est la borne supérieure des aires de tous les canapés.
\begin{enumerate}
\item[(a)] Démontrez les bornes élémentaires. Montrez que la forme de Hammersley --- deux quarts de disque unité flanquant un rectangle $ 1 \times 4/\pi $ auquel on a ôté un demi-disque de rayon $ 2/\pi $ --- est un canapé, et calculez son aire $ \pi/2 + 2/\pi \approx 2{,}2074 $. Démontrez ensuite qu'aucun canapé n'a une aire dépassant $ 2\sqrt{2} \approx 2{,}8284 $.
\item[(b)] \textbf{Revendiqué récemment.} Joseph Gerver a construit en 1992 un canapé d'aire $ \approx 2{,}2195 $ assemblé à partir de dix-huit arcs analytiques, et l'a conjecturé optimal. En novembre 2024, Jineon Baek a déposé une prépublication de 119 pages démontrant l'optimalité de la forme de Gerver ; en 2026, cet argument n'a pas achevé son évaluation par les pairs, de sorte que le problème est mieux décrit comme très probablement réglé mais pas encore formellement clos. Lisez la prépublication et vérifiez-en une portion autonome.
\item[(c)] \textbf{Problème ouvert.} Le problème du canapé \textit{ambidextre} demande la plus grande forme capable de négocier un virage à angle droit du même couloir dans les deux sens. Dan Romik a construit en 2018 un candidat d'aire $ \approx 1{,}6449 $, dont le bord est une courbe algébrique par morceaux donnée sous forme close. Aucune borne supérieure correspondante n'est connue. Déterminez la constante du canapé ambidextre.
\end{enumerate}

\end{problem}

\noindent Leo Moser a publié le problème en 1966, bien qu'il eût circulé auparavant, et sa célébrité tient à l'écart entre la facilité de son énoncé et l'acharnement de sa résistance. La difficulté est qu'une \og{} forme que l'on peut déplacer \fg{} est définie par une contrainte portant sur un espace de dimension infinie de programmes de rotation, de sorte que toute borne supérieure doit maîtriser simultanément tous les mouvements possibles ; Yoav Kallus et Dan Romik ont obtenu la borne $ 2{,}37 $ en 2018 par un argument assisté par ordinateur qui discrétise l'ensemble des angles de rotation et qui est, par conception, capable d'être poussé arbitrairement près de la vérité pourvu qu'on calcule assez. L'argument de Baek construit au contraire une fonctionnelle de majoration taillée sur la géométrie et procède sans calcul à grande échelle. La question (c) est là où les questions élémentaires ont migré : la version ambidextre a résisté à toutes les techniques qui avaient réussi sur la version manchote, et c'est actuellement l'endroit naturel où un nouveau venu peut attaquer ce cercle de problèmes.

\begin{refs}
\item Contexte : Problème du sofa --- Wikipédia (\url{https://fr.wikipedia.org/wiki/Probl%C3%A8me_du_sofa})
\item Article : Jineon Baek, \og{} Optimality of Gerver's sofa \fg{}, arXiv:2411.19826 (\url{https://arxiv.org/abs/2411.19826}) (2024) --- prépublication non encore expertisée
\item Article : Dan Romik, \og{} Differential equations and exact solutions in the moving sofa problem \fg{}, \textit{Experimental Mathematics} 27 (2018)
\item Voir aussi : Problèmes ouverts (\url{open-problems.html}), Géométrie (\url{pure-geometry.html})
\end{refs}

\bigskip

\begin{problem}[{Le Mastermind à nombreuses positions et nombreuses couleurs --- Recherche}]
\textbf{PUZ-43.} \textbf{Problème.} Généralisez le jeu de PUZ-42. Le secret est un mot de $ [k]^{n} $ ;
chaque essai est un mot de $ [k]^{n} $, et la réponse est le couple (nombre de positions
ayant la bonne couleur, nombre de bonnes couleurs à de mauvaises positions). Soit
$ f(n,k) $ le plus petit nombre d'essais qui suffise dans le pire des cas face à un codeur
adverse.
\begin{enumerate}
\item[(a)] Démontrez la borne informationnelle. Il y a au plus $ \binom{n+2}{2} $ réponses
possibles à un essai, d'où
\[ f(n,k) \;\ge\; \frac{n \log k}{\log \binom{n+2}{2}} \;=\; \Omega\!\left(\frac{n\log k}{\log n}\right). \]
Démontrez ensuite que, pour deux couleurs, le jeu équivaut à un problème de pesée de pièces,
et reconstituez le théorème d'Erdős et Rényi (1963) selon lequel $ \Theta(n/\log n) $ essais
sont nécessaires, $ \frac{2n}{\log_2 n}(1+o(1)) $ suffisant même en jeu non adaptatif.
\item[(b)] Reconstituez le paysage moderne et vérifiez les affirmations par vous-même : Chvátal
(1983) a démontré $ f(n,k) = \Theta(n \log k / \log n) $ pour
$ k \le n^{1-\varepsilon} $ ; Doerr, Doerr, Spöhel et Thomas (2016) ont obtenu
$ O(n \log\log n) $ essais pour $ k = n $ ; et Martinsson et Su ont refermé l'écart par
une stratégie randomisée utilisant $ O(n) $ essais pour $ k = n $, ce qui fixe l'ordre de
croissance de $ f(n,k) $ pour tous $ n $ et $ k $.
\item[(c)] \textbf{Ouvert.} Les constantes ne sont pas connues. Même dans le cas à deux couleurs, où
la réponse non adaptative a pour terme dominant $ 2n/\log_2 n $, on ignore si l'adaptativité
achète un meilleur terme dominant, et déterminer la constante optimale pour les stratégies
générales est la question ouverte centrale du sujet. Que la complexité en requêtes
\textit{déterministe} pour $ k = n $ soit également linéaire n'est pas démontré --- Martinsson
et Su le conjecturent sans le démontrer. Et l'on ne connaît aucune formule donnant la valeur
exacte $ f(n,k) $ à des paramètres concrets : même $ f(4,6) = 5 $ a exigé un calcul
exhaustif. Établissez ici quoi que ce soit de précis.
\end{enumerate}

\end{problem}

\noindent Le Mastermind est arrivé dans les magasins de jouets en 1971, mais les
mathématiques sont plus anciennes que le jeu : l'article d'Erdős et Rényi de 1963, \og{} On two
problems of information theory \fg{}, est exactement le cas à deux couleurs, posé comme question
de pesée de pièces, et il descend en droite ligne de PUZ-01 sur cette page. La note de Knuth
de 1976 comptait parmi les premiers exemples publiés d'utilisation d'un ordinateur pour
établir exactement, et non heuristiquement, un optimum de théorie des jeux. L'histoire
asymptotique a ensuite pris quatre décennies --- Chvátal, une longue série d'améliorations, et
enfin la borne linéaire de Martinsson et Su (arXiv 2020, publiée dans \textit{Combinatorics,
Probability and Computing} en 2024) --- et elle s'est achevée en fixant tous les ordres de
croissance et en laissant toutes les constantes ouvertes. Il y a aussi une seconde forme de
difficulté : décider s'il existe seulement un secret compatible avec une liste donnée
d'essais et de réponses est NP-complet, démontré par Michiel de Bondt en 2004 pour deux
couleurs et par Stuckman et Zhang en 2006 dans le cas général, de sorte que même tenir des
comptes honnêtes est intraitable. Notez ce que cela dit de PUZ-42 : le cas d'école est
résolu, les asymptotiques sont résolues, et l'entre-deux est un désert.

\begin{refs}
\item Contexte : Mastermind --- Wikipédia (\url{https://fr.wikipedia.org/wiki/Mastermind})
\item Article : Vašek Chvátal, \og{} Mastermind \fg{}, \textit{Combinatorica} 3 (1983), 325--329
\item Article : Benjamin Doerr, Carola Doerr, Reto Spöhel \& Henning Thomas, \og{} Playing Mastermind with many colors \fg{}, \textit{Journal of the ACM} 63 (2016), article 42
\item Article : Anders Martinsson \& Pascal Su, \og{} Mastermind with a linear number of queries \fg{} (arXiv:2011.05921) (\url{https://arxiv.org/abs/2011.05921}), \textit{Combinatorics, Probability and Computing} 33 (2024), 143--156
\end{refs}

\bigskip

\begin{problem}[{Les vingt questions avec un menteur --- Recherche}]
\textbf{PUZ-44.} \textbf{Problème.} Carole choisit $ x \in \{1, \ldots, n\} $. Paul pose des questions de
la forme \og{} $ x \in A $ ? \fg{} pour des parties $ A $, et Carole répond oui ou non, mais a le
droit de mentir au plus $ k $ fois dans toute la partie ; $ n $ et $ k $ sont connus des
deux, et Paul peut choisir chaque question après avoir entendu la réponse précédente. Soit
$ q_k(n) $ le plus petit nombre de questions au terme desquelles Paul est certain de
$ x $.
\begin{enumerate}
\item[(a)] Démontrez la borne de volume de Berlekamp : Paul ne peut pas gagner en $ q $ questions
à moins que
\[ n \sum_{i=0}^{k} \binom{q}{i} \;\le\; 2^{\,q}. \]
Démontrez ensuite que le jeu est exactement le problème du codage sur un canal binaire avec
au plus $ k $ erreurs et retour non bruité, et identifiez la borne ci-dessus à la borne de
Hamming pour de tels codes.
\item[(b)] Reconstituez les résultats exacts : Pelc (1987) a déterminé $ q_1(n) $ pour tout
$ n $, Guzicki (1990) a fait de même pour $ k = 2 $, et Spencer (1992) a montré que, pour
chaque $ k $ fixé, la borne de volume est atteinte dès que $ n $ est assez grand.
Réglez ensuite la question d'Ulam lui-même : $ n = 10^{6} $ avec au plus un mensonge.
\item[(c)] \textbf{Ouvert.} Remplacez le budget fixe de $ k $ mensonges par l'exigence que Carole
mente sur au plus une fraction $ r $ des réponses. Spencer et Winkler (1992) ont localisé
trois seuils en $ r $, qui diffèrent selon que la contrainte doit valoir à tout instant ou
seulement à la fin, et selon que les questions de Paul sont adaptatives ou non. Au-delà des
seuils, le tableau est inachevé : les asymptotiques exactes du plus grand espace de recherche
sur lequel Paul l'emporte face à un tel menteur linéairement borné ne sont pas connues.
Cooper et Ellis (2009) ont ramené la borne supérieure à un facteur
$ O\!\big(\sqrt{\log\log n}\,/(1-2f)\big) $ de la borne sphérique pour un rayon de
couverture $ fn $, et aucune minoration correspondante n'existe. Il n'existe pas non plus de
formule close unique pour $ q_k(n) $ valable pour tous $ n $ et $ k $ à la fois.
\end{enumerate}

\end{problem}

\noindent Alfréd Rényi a posé la question en 1961, en s'inspirant du Barkochba, la
forme hongroise des vingt questions ; Stanisław Ulam l'a reposée, indépendamment, dans son
autobiographie de 1976 \textit{Adventures of a Mathematician}, voulant savoir combien de
questions permettent de trouver un nombre inférieur au million lorsqu'une réponse peut être un
mensonge. Les deux noms voyagent ensemble depuis. Ce qui rend le sujet sérieux plutôt que
mignon, c'est la traduction de la question (a) : une stratégie pour Paul est un code
correcteur d'erreurs adaptatif, les mensonges de Carole sont le bruit du canal, et le retour
est la connaissance qu'a Paul des réponses déjà données. Elwyn Berlekamp a étudié précisément
cela dans sa thèse de doctorat de 1964 au MIT sur le codage par blocs avec retour non bruité,
et la borne de volume qu'il a obtenue est celle que rejoignent la plupart des résultats
exacts. Le panorama d'Andrzej Pelc, avec son merveilleux sous-titre \og{} cinquante ans à
composer avec les menteurs \fg{}, est la carte du territoire --- et c'est un territoire où les
questions propres, à $ k $ fixé, ont reçu réponse, et où le modèle le plus proche du bruit
réel n'en a pas reçu.

\begin{refs}
\item Contexte : Jeu d'Ulam --- Wikipédia (en anglais) (\url{https://en.wikipedia.org/wiki/Ulam%27s_game})
\item Article : Andrzej Pelc, \og{} Searching games with errors --- fifty years of coping with liars \fg{}, \textit{Theoretical Computer Science} 270 (2002), 71--109
\item Article : Joel Spencer \& Peter Winkler, \og{} Three thresholds for a liar \fg{}, \textit{Combinatorics, Probability and Computing} 1 (1992), 81--93
\item Article : Joshua Cooper \& Robert Ellis, \og{} Linearly bounded liars, adaptive covering codes, and deterministic random walks \fg{} (arXiv:0909.0029, 2009) (\url{https://arxiv.org/abs/0909.0029})
\end{refs}

\bigskip

\begin{problem}[{Le Chomp : une stratégie gagnante que personne ne sait nommer --- Recherche}]
\textbf{PUZ-56.} \textbf{Problème.} Le Chomp se joue sur une grille $ m \times n $ de carrés. Un coup
désigne un carré et le retire, ainsi que tous les carrés situés faiblement à sa droite et
faiblement au-dessous de lui. Le carré en haut à gauche est empoisonné, et le joueur contraint
de le manger perd.
\begin{enumerate}
\item[(a)] Démontrez, par un argument de vol de stratégie, que sur tout plateau rectangulaire autre
que $ 1 \times 1 $ le premier joueur possède une stratégie gagnante. Énoncez ensuite
exactement ce que votre démonstration livre et ne livre pas, et identifiez l'étape qui est non
constructive.
\item[(b)] Trouvez et démontrez des stratégies gagnantes explicites dans les cas que l'on comprend :
le plateau carré $ n \times n $, et le plateau à deux lignes $ 2 \times n $.
\item[(c)] \textbf{Problème ouvert.} Aucune stratégie gagnante explicite n'est connue pour le Chomp
$ m \times n $ général, et aucune n'est connue sous forme close, même pour le plateau à
trois lignes $ 3 \times n $. Doron Zeilberger, puis Xinyu Sun après lui, ont calculé
abondamment les positions perdantes du Chomp à trois lignes vers 2001-2002, et y ont trouvé de
fortes régularités mais aucune formule. Faites le point sur ce que l'on sait, et soit
produisez une stratégie, soit démontrez une minoration de la difficulté du calcul du coup
gagnant du premier joueur.
\item[(d)] Montrez que le même écart apparaît dans la version des diviseurs de Schuh : les positions
sont les diviseurs d'un entier fixé $ N $, un coup désigne un diviseur et le retire avec
tous ses multiples parmi les diviseurs, et le joueur contraint de prendre $ 1 $ perd.
Démontrez que le premier joueur gagne pour tout $ N > 1 $, et notez que là non plus aucune
stratégie générale n'est connue.
\end{enumerate}

\end{problem}

\noindent David Gale a donné la forme de la tablette de chocolat dans l'\textit{American
Mathematical Monthly} en 1974 ; Frederik Schuh avait publié plus tôt le jeu équivalent des
diviseurs. L'argument du vol de stratégie est le procédé même dont John Nash s'est servi pour
démontrer que le premier joueur gagne au Hex, et le Chomp est la plus nette réclame qui soit
pour la différence entre existence classique et existence constructive : nous savons avec
certitude qu'un premier coup gagnant existe sur le plateau $ 100 \times 100 $, et personne
sur Terre ne peut vous dire lequel. C'est un état du savoir réellement inconfortable, et il
revient dans toute la théorie des jeux combinatoires --- la liste courante de problèmes non
résolus de Richard Guy dans les volumes \textit{Games of No Chance} regorge de jeux dont on
connaît la classe d'issue sans en connaître les stratégies. Les calculs de Zeilberger ont
montré que le Chomp à trois lignes frôle la périodicité sans jamais être périodique, ce qui
lui a valu de sous-titrer un article voisin du mot \og{} chaos \fg{}.

\begin{refs}
\item Contexte : Chomp (jeu) --- Wikipédia (\url{https://fr.wikipedia.org/wiki/Chomp_%28jeu%29})
\item Contexte : Argument du vol de stratégie --- Wikipédia (en anglais) (\url{https://en.wikipedia.org/wiki/Strategy-stealing_argument})
\item Article : David Gale, \og{} A curious Nim-type game \fg{}, \textit{American Mathematical Monthly} 81 (1974), 876--879
\item Article : Doron Zeilberger, \og{} Three-rowed CHOMP \fg{}, \textit{Advances in Applied Mathematics} 26 (2001), exemplaire de l'auteur (\url{https://sites.math.rutgers.edu/\textasciitilde{}zeilberg/mamarim/mamarimPDF/chomp.pdf})
\end{refs}

\bigskip

\begin{problem}[{Le problème de Robbins : le problème du secrétaire que personne n'a résolu --- Recherche}]
\textbf{PUZ-57.} \textbf{Problème.} Vous observez des variables aléatoires indépendantes
$ X_1, \dots, X_n $, chacune uniforme sur $ [0,1] $, une à la fois. Après avoir vu
$ X_k $, vous devez soit vous arrêter --- en acceptant $ X_k $ irrévocablement ---, soit
continuer ; vous voyez les valeurs numériques réelles, et non seulement les rangs relatifs, et
vous ne pouvez pas revenir à une observation antérieure ; si vous atteignez $ X_n $ vous
devez l'accepter. Soit $ R $ le rang de la valeur acceptée parmi les $ n $, de sorte que
$ R = 1 $ signifie que vous avez accepté la plus petite, et soit $ V_n = \inf E[R] $,
l'infimum sur toutes les règles d'arrêt.
\begin{enumerate}
\item[(a)] Échauffez-vous sur le problème classique du secrétaire, où seuls les rangs relatifs sont
révélés et où l'objectif est de s'arrêter sur l'unique meilleur candidat. Démontrez qu'une
règle optimale observe un préfixe fixé puis accepte le premier record, et que la fraction
optimale du préfixe comme la probabilité de succès tendent toutes deux vers $ 1/e $.
\item[(b)] Restreignez maintenant le problème de Robbins aux règles \textit{sans mémoire}, dont la
décision à l'instant $ k $ peut dépendre de $ k $, de $ n $ et de $ X_k $, mais pas de
$ X_1, \dots, X_{k-1} $. Établissez les équations de récurrence rétrograde de la règle sans
mémoire optimale et calculez le rang espéré limite qu'elle atteint.
\item[(c)] Démontrez que la règle pleinement optimale dépend véritablement de l'histoire : exhibez un
$ n $, un instant $ k $, et deux histoires coïncidant en $ k $ et en $ X_k $ mais
différant sur les valeurs antérieures, pour lesquelles les décisions optimales diffèrent.
Expliquez pourquoi cela ruine la voie habituelle vers une solution, puisqu'aucune statistique
exhaustive de faible dimension n'est disponible.
\item[(d)] \textbf{Problème ouvert.} La valeur limite $ V = \lim_{n \to \infty} V_n $ est inconnue.
Les bornes habituellement citées sont
\[ 1{,}908 < V < 2{,}329 , \]
obtenues par Bruss et Ferguson à partir de règles de seuil sans mémoire. Améliorez l'une ou
l'autre borne, ou déterminez $ V $.
\end{enumerate}

\end{problem}

\noindent Herbert Robbins a posé ce problème lors d'une conférence à Amherst en 1990
et a conclu son allocution en disant qu'il aimerait le voir résolu avant de mourir. Il est
mort en 2001, et le problème est toujours ouvert. L'intérêt n'est pas la valeur numérique mais
la raison de sa résistance : le problème classique du secrétaire est soluble parce que les
rangs relatifs compriment le passé en un unique bit --- le candidat courant est-il un record ou
non ---, tandis que l'information complète en donne davantage à l'observateur et, ce faisant,
rend la règle optimale dépendante de tout ce qui a été vu. C'est un problème rare et posé avec
netteté, dans lequel plus d'information rend l'analyse plus difficile plutôt que plus facile,
et il se situe à peu de distance des questions de décision en ligne qui dominent aujourd'hui
l'ingénierie algorithmique des mécanismes.

\begin{refs}
\item Contexte : Problème de Robbins --- Wikipédia (\url{https://fr.wikipedia.org/wiki/Probl%C3%A8me_de_Robbins})
\item Contexte : Problème du secrétaire --- Wikipédia (\url{https://fr.wikipedia.org/wiki/Probl%C3%A8me_du_secr%C3%A9taire})
\item Article : F. T. Bruss \& T. S. Ferguson, \og{} Minimizing the expected rank with full information \fg{}, \textit{Journal of Applied Probability} 30 (1993), 616--626
\item Lecture : Thomas S. Ferguson, \textit{Optimal Stopping and Applications (\url{https://www.math.ucla.edu/\textasciitilde{}tom/Stopping/Contents.html})} --- notes de cours en libre accès (en anglais) ; voir aussi Optimisation et théorie des jeux (\url{applied-optimization.html}) pour la conjecture du secrétaire matroïdal
\end{refs}

\bigskip

\begin{problem}[{La Belle au bois dormant --- Recherche, short}]
\textbf{PUZ-58.} \textbf{Problème.} On expose à l'avance le protocole à la Belle au bois
dormant : le dimanche on l'endort, une pièce équilibrée est lancée, et on la réveille une
fois --- le lundi --- si elle tombe sur pile, mais deux fois --- le lundi puis de nouveau le mardi ---
si elle tombe sur face, une drogue effaçant chaque réveil de sorte que les réveils lui soient
indiscernables. Réveillée et interrogée sur sa crédence que la pièce est tombée sur pile, elle
doit donner un nombre : construisez le meilleur plaidoyer possible en faveur de $ 1/2 $ et en
faveur de $ 1/3 $, exhibez explicitement l'espace probabilisé sur lequel chacun d'eux est
une probabilité, et identifiez l'hypothèse sur la croyance auto-localisante qui sépare les
deux.
\end{problem}

\noindent Arnold Zuboff a formulé le problème dans des travaux non publiés au milieu
des années 1980 ; Robert Stalnaker lui a donné son nom, il s'est répandu par le forum
\textit{rec.puzzles} à la fin des années 1990, et la littérature moderne commence avec l'article
d'Adam Elga de 2000 plaidant pour $ 1/3 $ et la réponse de David Lewis de 2001 plaidant pour
$ 1/2 $. Il figure dans la bande Recherche de cette page pour une raison honnête : après un
quart de siècle, il n'y a pas de consensus, et le litige ne porte pas sur un calcul mais sur la
question de savoir si un agent rationnel peut mettre à jour ses croyances sur une information
purement auto-localisante --- apprendre \textit{où} ou \textit{quand} l'on est plutôt que \textit{comment
est le monde}. Les positions du \og{} moitié \fg{}, du \og{} tiers \fg{} et du \og{} double moitié \fg{}, ainsi que
l'opinion selon laquelle la question telle que posée est ambiguë, ont toutes des défenseurs
sérieux.
La même question anime l'argument de l'apocalypse et le raisonnement anthropique en
cosmologie, de sorte que la réponse ne se cantonne pas à un conte de fées, et la première
chose que doit faire une bonne tentative est de dire ce qui compterait comme une réponse.

\begin{refs}
\item Contexte : Problème de la Belle au bois dormant --- Wikipédia (\url{https://fr.wikipedia.org/wiki/Probl%C3%A8me_de_la_Belle_au_bois_dormant})
\item Contexte : Self-Locating Beliefs --- Stanford Encyclopedia of Philosophy (en anglais) (\url{https://plato.stanford.edu/entries/self-locating-beliefs/})
\item Article : Adam Elga, \og{} Self-locating belief and the Sleeping Beauty problem \fg{}, \textit{Analysis} 60 (2000), 143--147
\item Article : David Lewis, \og{} Sleeping Beauty: reply to Elga \fg{}, \textit{Analysis} 61 (2001), 171--176
\end{refs}

\bigskip

\begin{problem}[{L'effet Mpemba --- Recherche}]
\textbf{PUZ-71.} \textbf{Problème.} Fixez un protocole de refroidissement : un récipient de géométrie et de
matériau spécifiés, un bain maintenu à la température $ T_b $, et un critère de
\og{} congelé \fg{}. Pour chaque température initiale $ T_0 $, soit $ \tau(T_0) $ le temps mis
pour satisfaire le critère. L'effet Mpemba est l'affirmation que $ \tau $ n'est pas
croissante en $ T_0 $.
\begin{enumerate}
\item[(a)] Donnez une définition assez nette pour être testée : énoncez exactement ce qui doit être
maintenu fixe d'un terme à l'autre de la comparaison, ce qui compte comme spécification d'une
condition initiale, et ce que \og{} congelé \fg{} veut dire. Montrez ensuite que plusieurs
définitions en usage ne sont pas équivalentes, en exhibant un système qui satisfait l'une et
viole l'autre.
\item[(b)] Démontrez que l'effet est possible en principe. Construisez un processus de sauts
markovien explicite à états finis, muni d'une distribution initiale dépendant de la
température, dont la distance à l'équilibre froid n'est pas monotone en $ T_0 $, et
caractérisez quels spectres de relaxation admettent un tel comportement.
\item[(c)] Concevez une expérience sur l'eau qui établirait ou réfuterait l'effet selon votre
définition de la question (a), en contrôlant l'évaporation, les gaz dissous, la surfusion et
la nucléation, la convection, la couche de givre sous le récipient et la conductivité du
récipient. Indiquez combien d'essais votre dispositif exige pour une puissance statistique
donnée, et pourquoi.
\item[(d)] \textbf{Contesté.} Tranchez le cas de l'eau : soit produisez un protocole contrôlé sous
lequel l'eau manifeste l'effet de façon reproductible, avec un mécanisme énoncé et une analyse
préenregistrée, soit démontrez qu'à l'intérieur d'un modèle explicitement énoncé aucun tel
protocole ne peut exister.
\end{enumerate}

\end{problem}

\noindent Erasto Mpemba avait treize ans, à la Magamba Secondary School de Tanzanie
en 1963, lorsqu'il a remarqué qu'un mélange à crème glacée chaud gelait avant un mélange
froid, et que son professeur lui a répondu qu'il s'embrouillait. Denis Osborne, de passage
depuis l'University College de Dar es Salaam, a pris la question au sérieux, l'a reprise, et a
publié avec lui dans \textit{Physics Education} en 1969 sous le titre \og{} Cool? \fg{}. Aristote, Bacon
et Descartes avaient tous consigné des versions de cette affirmation. Le statut est
aujourd'hui honnêtement mitigé, et il faut le dire ainsi. Burridge et Linden ont cherché
sérieusement en 2016 et ont conclu qu'il n'existe aucun élément à l'appui d'observations
significatives de l'effet Mpemba dans de l'eau atteignant $ 0^\circ\mathrm{C} $. Pendant ce
temps, la reformulation moderne --- un temps de relaxation vers un bain froid peut-il ne pas
être monotone en la température initiale ? --- s'est révélée à la fois nette et affirmative :
Lu et Raz l'ont démontrée pour une dynamique markovienne en 2017, et Kumar et Bechhoefer l'ont
réalisée expérimentalement dans un piège colloïdal en 2020, avec une accélération
exponentielle. Depuis 2024, la question s'est étendue aux systèmes quantiques à N corps, des
effets Mpemba étant rapportés dans des ions piégés et des spins nucléaires. L'eau, le sujet
d'origine, est le cas le plus difficile, parce que chaque mécanisme proposé est aussi un
facteur de confusion ; une prépublication de 2025 de Klimov et Finkelstein soutient que, dans
l'eau pure, l'effet n'apparaît que lorsque le congélateur se trouve très près de la
température de nucléation de la glace, et qu'il est alors une conséquence du caractère
aléatoire de la nucléation plutôt que du refroidissement. Cette affirmation n'est pas encore
expertisée. Traitez tout le problème comme une étude de cas sur ce qu'il faut pour transformer
une observation populaire en observation testable.

\begin{refs}
\item Contexte : Effet Mpemba --- Wikipédia (\url{https://fr.wikipedia.org/wiki/Effet_Mpemba})
\item Article : E. B. Mpemba \& D. G. Osborne, \og{} Cool? \fg{}, \textit{Physics Education} 4 (1969), 172--175
\item Article : H. C. Burridge \& P. F. Linden, \og{} Questioning the Mpemba effect: hot water does not cool more quickly than cold \fg{}, \textit{Scientific Reports} 6 (2016), 37665 ; et A. Kumar \& J. Bechhoefer, \og{} Exponentially faster cooling in a colloidal system \fg{}, \textit{Nature} 584 (2020), 64--68
\item Article : Z. Lu \& O. Raz, \og{} Nonequilibrium thermodynamics of the Markovian Mpemba effect and its inverse \fg{}, \textit{PNAS} 114 (2017), 5083--5088 ; et (prépublication, non expertisée) A. A. Klimov \& A. V. Finkelstein, arXiv:2508.05607 (\url{https://arxiv.org/abs/2508.05607}) (2025)
\end{refs}

\bigskip

\begin{problem}[{La fontaine de chaîne --- Recherche, short}]
\textbf{PUZ-72.} \textbf{Problème.} Une longue chaîne à billes repose enroulée dans un bécher
posé sur une étagère haute ; une extrémité est tirée par-dessus le bord puis lâchée, et la
chaîne qui tombe extrait le reste si vite qu'une arche stationnaire s'élève bien au-dessus du
bord avant de redescendre jusqu'au sol. Construisez un modèle de la zone d'extraction qui
prédise la hauteur de l'arche en fonction de la hauteur de chute, de la masse et de
l'espacement des billes, de la rigidité des maillons et de la géométrie de l'enroulement ;
démontrez que votre modèle conserve ce qu'il doit conserver et dissipe ce qu'il doit dissiper ;
et déterminez à laquelle de ses composantes la hauteur prédite est réellement sensible.
\end{problem}

\noindent Steve Mould a filmé l'effet en 2013 et il en a pris le nom. John Biggins et
Mark Warner ont publié un modèle dans les \textit{Proceedings of the Royal Society A} en 2014,
où la chaîne reçoit une impulsion vers le haut de la part du récipient à mesure que chaque
maillon rigide est extrait du tas comme par un levier ; leur modèle reste le plus cité --- mais
il ne fait pas l'unanimité. L'analyse quantitative de Pantaleone dans l'\textit{American Journal
of Physics} (2017), l'article de Flekkøy, Moura et Måløy dans \textit{Frontiers in Physics}
(2018) et le \og{} Reexamining the chain fountain \fg{} d'Hiroshi Yokoyama (2018) parviennent à des
conclusions matériellement différentes sur la part de l'effet que fournit l'impulsion ;
Yokoyama soutient qu'extraire une chaîne de son propre mou est bien moins dissipatif que ne le
supposaient Biggins et Warner, et construit une fontaine qu'il attribue plutôt à la dynamique
de l'arche suspendue. Sous le désaccord se trouve un morceau réellement difficile de mécanique
des milieux continus : un corps unidimensionnel inextensible que l'on extrait d'un tas est un
problème à frontière libre dans lequel le bilan de quantité de mouvement et le bilan d'énergie
ne sont pas équivalents, et la différence se joue à l'intérieur d'une région que personne ne
peut observer. Cet échec est ancien --- le problème de la chaîne qui tombe de Cayley, en 1857,
est le cas classique où l'argument énergétique naïf et l'argument de quantité de mouvement
naïf se contredisent, et la contradiction est physique et non algébrique. Le statut est donc
le suivant : le phénomène est robuste et reproductible sur n'importe quelle table de cuisine,
et son mécanisme est ouvert.

\begin{refs}
\item Contexte : Fontaine de chaîne --- Wikipédia (\url{https://fr.wikipedia.org/wiki/Fontaine_de_cha%C3%AEne})
\item Contexte : Chaînette --- Wikipédia (\url{https://fr.wikipedia.org/wiki/Cha%C3%AEnette})
\item Article : J. S. Biggins \& M. Warner, \og{} Understanding the chain fountain \fg{}, \textit{Proceedings of the Royal Society A} 470 (2014), 20130689
\item Article : J. Pantaleone, \og{} A quantitative analysis of the chain fountain \fg{}, \textit{American Journal of Physics} 85 (2017) ; et (prépublication) H. Yokoyama, \og{} Reexamining the chain fountain \fg{}, arXiv:1810.13008 (\url{https://arxiv.org/abs/1810.13008}) (2018)
\end{refs}

\bigskip

\begin{problem}[{La spirale d'Ulam et les diagonales que personne ne sait expliquer --- Recherche}]
\textbf{PUZ-85.} \textbf{Problème.} Écrivez les entiers strictement positifs en spirale carrée --- $ 1 $ au
centre, puis $ 2, 3, 4, \dots $ en s'enroulant vers l'extérieur --- et marquez chaque nombre
premier. Des lignes diagonales de nombres premiers apparaissent de façon frappante, bien plus
marquées que dans un marquage aléatoire de même densité.
\begin{enumerate}
\item[(a)] Démontrez que toute diagonale, ligne ou colonne complète de la spirale est l'ensemble des
valeurs d'un polynôme quadratique $ f(x) = 4x^{2} + bx + c $ à $ b $ et $ c $ entiers ;
déterminez exactement quels couples $ (b,c) $ apparaissent. Déduisez-en qu'une ligne dont le
polynôme se factorise sur $ \mathbb{Z} $, ou dont les valeurs partagent un facteur premier
commun, ne porte qu'un nombre fini de nombres premiers, et donc que certaines lignes sont vides
pour des raisons triviales.
\item[(b)] Énoncez précisément la conjecture F de Hardy et Littlewood : pour
$ f(x) = ax^{2}+bx+c $ avec $ a > 0 $, avec $ a+b $ et $ c $ non tous deux pairs, et
avec $ D = b^{2}-4ac $ non carré parfait, le nombre de nombres premiers parmi
$ f(1), \dots, f(n) $ vaut asymptotiquement
\[ C \cdot \frac{\sqrt{n}}{\log n}, \]
où $ C $ est donné par un produit eulérien explicite sur les nombres premiers faisant
intervenir le symbole de Legendre de $ D $. Calculez $ C $ pour le polynôme d'Euler
$ x^{2}-x+41 $ et pour une diagonale de votre choix, et testez numériquement les deux
prédictions jusqu'à $ 10^{7} $.
\item[(c)] Établissez honnêtement ce qui est démontrable et ce qui ne l'est pas. Montrez qu'aucune
démonstration n'est connue qu'un seul polynôme quadratique à une variable représente une
infinité de nombres premiers --- le cas $ x^{2}+1 $ est le quatrième problème de Landau et il
est ouvert --- et reconstituez à la place ce que donnent les méthodes de crible : Iwaniec a
démontré en 1978 que $ x^{2}+1 $ prend une infinité de valeurs ayant au plus deux facteurs
premiers, et Friedlander et Iwaniec ont démontré en 1998 que la forme à deux variables
$ x^{2}+y^{4} $ représente une infinité de nombres premiers. Identifiez précisément l'étape
où un argument de crible pour une quadratique à une variable s'effondre.
\item[(d)] Rendez compte du polynôme d'Euler aussi loin qu'on le puisse : démontrez le critère de
Rabinowitsch de 1913, selon lequel $ x^{2}+x+p $ est premier pour tout
$ x = 0, 1, \dots, p-2 $ exactement lorsque le corps quadratique imaginaire
$ \mathbb{Q}(\sqrt{1-4p}) $ a un nombre de classes égal à un, et reliez-le aux nombres de
Heegner et à la valeur $ 163 $.
\item[(e)] \textbf{Ouvert.} Démontrez la conjecture F, ou démontrez qu'un certain polynôme quadratique
à une variable représente une infinité de nombres premiers, ou démontrez que les diagonales
visibles de la spirale persistent --- c'est-à-dire que le rapport entre la densité de nombres
premiers le long des lignes les plus riches et la densité moyenne ne tend pas vers $ 1 $.
Aucune des trois n'est connue.
\end{enumerate}

\end{problem}

\noindent Stanisław Ulam a dessiné la spirale en 1963 en griffonnant pendant ce qu'il
a décrit comme un long et très ennuyeux exposé lors d'une réunion scientifique ; il a remarqué
les diagonales, et a fait calculer et tracer le motif. Martin Gardner l'a mis en couverture de
\textit{Scientific American} en mars 1964 et il est célèbre depuis. La situation honnête mérite
d'être dite platement, car l'image invite à croire que quelque chose a été découvert. Ce que
les diagonales montrent est réel, quantitatif et complètement inexpliqué par le moindre
théorème : certaines quadratiques sont observablement plus riches en nombres premiers que
d'autres, Hardy et Littlewood ont écrit en 1923 exactement de combien chacune devrait l'être,
les prédictions collent au calcul sur autant de décimales que l'on veut bien en calculer --- et
pas une seule n'a été démontrée, ni personne n'a démontré qu'une quadratique à une variable
produit seulement une infinité de nombres premiers. La question (d) est le seul endroit où une
véritable explication existe, et elle est magnifique : le polynôme d'Euler est exceptionnel
parce qu'un certain corps quadratique imaginaire a la factorisation unique, fait qui porte sur
les nombres de classes et, en dernière analyse, sur les neuf nombres de Heegner. Cette
explication couvre exactement une ligne de l'image, et rien d'autre.

\begin{refs}
\item Contexte : Spirale d'Ulam --- Wikipédia (\url{https://fr.wikipedia.org/wiki/Spirale_d%27Ulam})
\item Contexte : Problèmes de Landau --- Wikipédia (\url{https://fr.wikipedia.org/wiki/Probl%C3%A8mes_de_Landau})
\item Contexte : Nombre de Heegner --- Wikipédia (\url{https://fr.wikipedia.org/wiki/Nombre_de_Heegner})
\item Article : G. H. Hardy \& J. E. Littlewood, \og{} Some problems of `Partitio numerorum' III: On the expression of a number as a sum of primes \fg{}, \textit{Acta Mathematica} 44 (1923) --- la conjecture F y est énoncée
\item Voir aussi : Problèmes ouverts (\url{open-problems.html}), Théorie des nombres (\url{pure-number-theory.html})
\end{refs}

\bigskip

\begin{problem}[{Le problème de Dürer : tout polyèdre convexe se déplie-t-il ? --- Recherche}]
\textbf{PUZ-86.} \textbf{Problème.} Soit $ P $ un polyèdre convexe. Un \textit{dépliage par arêtes} de
$ P $ découpe le long d'un ensemble d'arêtes et aplatit la surface dans le plan en une seule
pièce connexe ; c'est un \textit{patron} si la figure plane obtenue est simple, c'est-à-dire ne
se recouvre pas elle-même.
\begin{enumerate}
\item[(a)] Démontrez qu'un ensemble d'arêtes dont le retrait laisse la surface connexe et
aplatissable simplement est exactement un arbre couvrant du graphe sommets-arêtes de $ P $.
Comptez les arbres couvrants de chacun des cinq solides platoniciens, et dites ainsi combien
de dépliages par arêtes candidats chacun admet.
\item[(b)] Exhibez un polyèdre convexe accompagné d'un dépliage par arêtes qui se recouvre lui-même,
de sorte qu'être un dépliage par arêtes soit strictement plus faible qu'être un patron.
Exhibez ensuite un polyèdre \textit{non convexe} pour lequel \textit{aucun} dépliage par arêtes
n'est un patron, et démontrez qu'aucun ne l'est.
\item[(c)] Démontrez que tout polyèdre convexe possède bel et bien un dépliage \textit{général} sans
recouvrement, où les coupes n'ont pas à suivre les arêtes : construisez soit le dépliage
source, obtenu en découpant le lieu de coupure d'un point source choisi, soit le dépliage en
étoile, obtenu en découpant les plus courts chemins d'un point source générique vers tous les
sommets, et démontrez que le résultat ne se recouvre pas.
\item[(d)] \textbf{Problème ouvert.} Tout polyèdre convexe possède-t-il un patron --- un dépliage par
arêtes simple ? C'est le problème de Dürer, posé sous cette forme par G. C. Shephard en 1975
et non résolu. Ce que l'on sait : Mohammad Ghomi a démontré en 2014 que tout polyèdre convexe
admet un patron après une transformation affine convenable, et Alexander Barvinok et Ghomi ont
montré en 2019 que la généralisation naturelle aux graphes de pseudo-arêtes --- des réseaux de
géodésiques joignant les sommets --- est fausse, ce qui écarte toute une classe de démonstrations
tentées. Tranchez la conjecture, ou produisez un polyèdre convexe sans patron.
\item[(e)] Une question ouverte voisine, au cas où (d) résisterait : démontrez ou réfutez que tout
patron d'un polyèdre convexe peut être replié par un mouvement continu ne laissant jamais la
surface se traverser elle-même --- la question de l'\og{} éclosion \fg{}.
\end{enumerate}

\end{problem}

\noindent L'\textit{Underweysung der Messung} d'Albrecht Dürer, en 1525, est un manuel
de géométrie pratique pour les artistes et les artisans, et il dessine entre autres les solides
platoniciens et plusieurs solides archimédiens comme des patrons plans à découper et à plier ---
les premiers patrons de polyèdres sur le papier. Dürer les a dessinés ; il ne s'est pas demandé
si cela marche toujours. Personne ne l'a fait jusqu'à ce que G. C. Shephard mette la question
par écrit en 1975, et elle est ouverte depuis, ce qui est remarquable pour un problème dont
l'énoncé n'exige d'autre vocabulaire que des ciseaux et du carton. Le statut doit être énoncé
avec soin. La question (c) est un théorème : tout polyèdre convexe peut être déplié sans
recouvrement si l'on vous autorise à couper à travers les faces, de sorte que la difficulté de
(d) tient entièrement à la restriction aux arêtes. Le résultat affine de Ghomi, en 2014, dit
que l'obstruction, quelle qu'elle soit, n'est pas préservée par les applications affines, et le
contre-exemple de Barvinok et Ghomi, en 2019, dit que le renforcement le plus naturel est tout
simplement faux, de sorte que plusieurs voies plausibles sont fermées. Les recherches par
ordinateur sur les arbres couvrants de polyèdres convexes aléatoires trouvent constamment des
dépliages qui se recouvrent, et n'ont encore jamais trouvé de polyèdre dont tous les dépliages
se recouvrent. C'est l'extrémité \og{} recherche \fg{} de PUZ-76 (\url{#PUZ-76}), où déplier une
boîte était une astuce ; ici, c'est le sujet.

\begin{refs}
\item Contexte : Patron (géométrie) --- Wikipédia (\url{https://fr.wikipedia.org/wiki/Patron_%28g%C3%A9om%C3%A9trie%29}), y compris la section sur la conjecture de Dürer
\item Lecture : Erik D. Demaine \& Joseph O'Rourke, \textit{Geometric Folding Algorithms: Linkages, Origami, Polyhedra} (Cambridge, 2007), partie III sur le dépliage des polyèdres
\item Article : Mohammad Ghomi, \og{} Affine unfoldings of convex polyhedra \fg{}, \textit{Geometry \& Topology} 18 (2014)
\item Voir aussi : PUZ-76 (\url{#PUZ-76}) ci-dessus, Problèmes ouverts (\url{open-problems.html}), Géométrie (\url{pure-geometry.html})
\end{refs}

\bigskip

\begin{problem}[{L'ange et le démon --- Recherche}]
\textbf{PUZ-99.} \textbf{Problème.} On joue sur les cases d'un plateau infini $ \mathbb{Z}^{2} $. L'Ange de
puissance $ k $ occupe une case. Les deux joueurs alternent : le Démon détruit une case
quelconque non occupée à cet instant par l'Ange, et l'Ange se déplace ensuite vers toute case
atteignable en au plus $ k $ coups de roi d'échecs, à condition que cette case n'ait pas été
détruite --- il peut survoler les cases détruites mais ne peut pas s'y poser. Le Démon gagne si
l'Ange n'a plus de coup licite ; l'Ange gagne en se déplaçant à jamais.
\begin{enumerate}
\item[(a)] Démontrez que le Démon bat l'Ange de puissance $ 1 $ --- un roi d'échecs ordinaire --- en
exhibant une stratégie du Démon et en démontrant qu'elle piège le roi après un nombre fini de
coups.
\item[(b)] Démontrez que le Démon bat un Ange de puissance $ k $ quelconque à qui il est interdit
de jamais diminuer sa coordonnée $ y $. Démontrez la même chose pour un Ange tenu d'augmenter sa
distance à l'origine à chaque coup.
\item[(c)] \textbf{Résolu, 2006.} Après vingt-quatre ans, quatre démonstrations indépendantes sont
parues en quelques mois établissant qu'un Ange suffisamment puissant gagne : Brian Bowditch
pour la puissance $ 4 $, Oddvar Kloster et András Máthé chacun pour la puissance $ 2 $, et
Péter Gács pour une grande puissance. Reconstituez-en une intégralement --- celle de Kloster est
une stratégie d'Ange explicitement décrite et spécifiée de façon finie --- et combinez-la avec
(a) pour confirmer que le problème plan est désormais réglé sans réserve : le Démon gagne pour
$ k = 1 $ et l'Ange gagne pour tout $ k \ge 2 $.
\item[(d)] \textbf{Ouvert.} La variante tridimensionnelle dans laquelle l'Ange doit augmenter sa coordonnée
$ y $ à chaque coup tandis que le Démon ne peut détruire que des cases situées dans trois
plans fixés n'a été tranchée ni dans un sens ni dans l'autre. Plus généralement, on ne connaît
aucune caractérisation des graphes infinis sur lesquels un Ange de puissance finie gagne.
Tranchez l'une ou l'autre question, ou produisez une version quantitative de (c) : bornez le
nombre de coups dont l'Ange de puissance $ 2 $ a besoin pour s'échapper de toute région
bornée, en fonction de la taille de celle-ci.
\end{enumerate}

\end{problem}

\noindent John Conway a publié le problème en 1982 dans \textit{Winning Ways}, avec
Berlekamp et Guy, sous le nom de \og{} l'ange et le mangeur de cases \fg{}, et il y est revenu dans un
essai de 1996 du recueil \textit{Games of No Chance} de Richard Nowakowski, où il y a attaché des
prix : cent dollars pour une stratégie gagnante d'un Ange de puissance suffisamment élevée, et
mille pour une démonstration que le Démon bat tout Ange. Aucun des deux prix n'a été réclamé
pendant vingt-quatre ans, et la plus grosse somme était du côté perdant. Ce qui rend le
problème instructif, c'est la forme de la difficulté. Les cases détruites par le Démon
s'accumulent à jamais, de sorte qu'une stratégie d'Ange doit être démontrée correcte face à un
adversaire aux ressources accumulées non bornées, et toute politique gloutonne naturelle ---
courir dans une direction fixe, fuir la case détruite la plus proche, maximiser la liberté
immédiate --- peut être mise en échec, ce qui est exactement ce que la question (b) vous demande
de découvrir par vous-même. Le cas de puissance 1 est facile et a égaré la discipline pendant
deux décennies, en lui faisant attendre une victoire du Démon en général. Les arguments de
Kloster et de Máthé sont assez élémentaires pour se lire en un après-midi, ce qui est inhabituel
pour un problème qui a résisté si longtemps, et ils sont la meilleure réclame possible en
faveur de l'attaque des vieux problèmes par une comptabilité neuve plutôt que par une
machinerie neuve.

\begin{refs}
\item Contexte : Problème de l'ange --- Wikipédia (\url{https://fr.wikipedia.org/wiki/Probl%C3%A8me_de_l%27ange})
\item Article : Oddvar Kloster, \og{} A solution to the angel problem \fg{}, \textit{Theoretical Computer Science} 389 (2007), nos 1--2, 152--161
\item Article : András Máthé, \og{} The angel of power 2 wins \fg{}, \textit{Combinatorics, Probability and Computing} 16 (2007), no 3, 363--374
\item Article : Brian H. Bowditch, \og{} The angel game in the plane \fg{}, \textit{Combinatorics, Probability and Computing} 16 (2007), no 3, 345--362 --- voir aussi Problèmes ouverts (\url{open-problems.html})
\end{refs}

\bigskip

\begin{problem}[{Le problème de Newcomb --- Recherche}]
\textbf{PUZ-100.} \textbf{Problème.} Un Prédicteur, dont les prévisions des choix humains se sont révélées
exactes dans une longue série d'essais antérieurs consignés, place deux boîtes devant vous. La
boîte A est transparente et contient mille livres. La boîte B est opaque et contient soit un
million de livres, soit rien. Vous pouvez prendre les deux boîtes, ou la boîte B seule. Le
Prédicteur a placé l'argent hier, selon la règle suivante : un million dans la boîte B s'il a
prédit que vous prendriez la boîte B seule, rien dans la boîte B sinon. Le contenu est
maintenant fixé et il a quitté les lieux.
\begin{enumerate}
\item[(a)] Écrivez intégralement et précisément les deux arguments standard. Énoncez l'argument de
dominance, en identifiant la partition d'états sur laquelle la dominance est appliquée ;
énoncez l'argument d'utilité espérée, en donnant les probabilités conditionnelles qu'il
utilise. Démontrez que chacun est valide compte tenu de ses propres prémisses, et identifiez
l'unique prémisse sur laquelle ils divergent.
\item[(b)] Formalisez cette prémisse. Définissez les valuations évidentielle et causale
\[ V_{\mathrm{EDT}}(a) \;=\; \sum_{s} P(s \mid a)\, u(a,s), \qquad
 V_{\mathrm{CDT}}(a) \;=\; \sum_{K} P(K) \sum_{s} P_{K}(s \mid a)\, u(a,s), \]
où $ K $ parcourt les hypothèses sur la structure causale de la situation. Calculez les deux
pour ce problème comme fonctions de la précision $ q $ du Prédicteur, et déterminez les
valeurs de $ q $ auxquelles chaque règle change de recommandation.
\item[(c)] Montrez que le désaccord ne se cantonne pas aux prédicteurs fantaisistes. Construisez une
décision médicale ou économique ordinaire qui soit un problème de Newcomb déguisé et sur
laquelle les deux règles diffèrent. Analysez ensuite le cas de la lésion du fumeur, qui est
l'argument standard en sens inverse, et énoncez exactement pourquoi on le juge embarrassant
pour la règle évidentielle.
\item[(d)] \textbf{Non résolu.} Plus de cinquante ans après, il n'y a pas de consensus : l'enquête
PhilPapers de 2020 auprès des philosophes professionnels a trouvé environ $ 39\% $ en faveur
de la prise des deux boîtes et $ 31\% $ en faveur d'une seule, le reste étant indécis ou
rejetant la question. Engagez-vous. Énoncez une théorie de la décision avec assez de précision
pour que sa recommandation ici soit un théorème ; démontrez ce théorème ; puis exhibez un
problème de décision sur lequel votre théorie renvoie une réponse que vous êtes prêt à
défendre et sa rivale non.
\end{enumerate}

\end{problem}

\noindent William Newcomb, physicien théoricien au laboratoire Lawrence Livermore, a
inventé le problème vers 1960 en réfléchissant au dilemme du prisonnier. Il est parvenu à
l'imprimé dans l'essai de Robert Nozick de 1969, \og{} Newcomb's problem and two principles of
choice \fg{}, écrit pour les mélanges offerts à Carl Hempel, et la chronique de Martin Gardner dans
\textit{Scientific American} en 1974 lui a valu un courrier des lecteurs exceptionnellement
abondant. Le résumé qu'en donnait Nozick n'a jamais été surpassé : pour presque tout le monde,
il est parfaitement clair et évident ce qu'il faut faire ; l'ennui est que ces gens semblent se
partager à peu près à parts égales, un grand nombre estimant que la moitié adverse est tout
bonnement stupide. Les mathématiques qui en sont nées ne sont pas décoratives. La théorie
causale de la décision a été bâtie à la fin des années 1970 et au début des années 1980 par
Allan Gibbard et William Harper, David Lewis et Brian Skyrms, précisément pour rendre l'argument
de dominance respectable face à la théorie évidentielle de Richard Jeffrey de 1965, et le débat
est vivant : il resurgit aujourd'hui dans la conception d'agents automatisés, où savoir si une
procédure de décision doit traiter des copies corrélées d'elle-même comme des effets ou comme
des indices est une question d'ingénierie avec de l'argent à la clé. Ce problème figure dans
cette bande pour la même raison honnête que la Belle au bois dormant en PUZ-58 --- le désaccord
ne porte pas sur un calcul, et le premier devoir d'une tentative sérieuse est de dire ce qui
compterait comme trancher la question.

\begin{refs}
\item Contexte : Paradoxe de Newcomb --- Wikipédia (\url{https://fr.wikipedia.org/wiki/Paradoxe_de_Newcomb})
\item Contexte : Théorie causale de la décision --- Wikipédia (en anglais) (\url{https://en.wikipedia.org/wiki/Causal_decision_theory})
\item Article : Robert Nozick, \og{} Newcomb's problem and two principles of choice \fg{}, dans Nicholas Rescher (dir.), \textit{Essays in Honor of Carl G. Hempel} (Reidel, 1969), 114--146
\item Article : David Lewis, \og{} Prisoners' dilemma is a Newcomb problem \fg{}, \textit{Philosophy \& Public Affairs} 8 (1979), 235--240 --- voir aussi PUZ-58 (\url{#PUZ-58}) ci-dessus
\end{refs}

\bigskip


\section*{Pour aller plus loin}


\vfill
\noindent\emph{Hors de la Caverne} --- des probl\`emes, pas de r\'eponses.
\end{document}
